\documentclass[a4paper,12pt,leqno]{article}
\usepackage[T1]{fontenc}
\usepackage[german]{babel}
\usepackage{epsf}

\renewcommand{\neq}{\mathrel{\vcenter{\halign{\hfil$##$\hfil\cr
   >\cr\noalign{\kern-8pt}<\cr}}}}

\renewcommand{\leq}{\mathrel{\vcenter{\halign{\hfil$##$\hfil\cr
   <\cr\noalign{\kern-8pt}=\cr}}}}

\renewcommand{\geq}{\mathrel{\vcenter{\halign{\hfil$##$\hfil\cr
   >\cr\noalign{\kern-8pt}=\cr}}}}


\renewcommand{\theequation}{\arabic{equation}.}

\begin{document}

\thispagestyle{empty}
\begin{center}
\Large\bfseries
Theorie der Abel'schen Functionen.\\[12 pt]
Bernhard Riemann\\[12 pt]
[Journal f\"{u}r die reine und angewandte Mathematik,
Bd.~54.  S.~101--155.  1857.]\\[24 pt]
\large\mdseries
Transcribed by D. R. Wilkins\\[12 pt]
Preliminary Version: December 1998\\
Corrected: April 2000
\end{center}

\newpage
\setcounter{page}{1}

\title{Theorie der Abel'schen Functionen.}
\author{Bernhard Riemann}
\date{[Journal f\"{u}r die reine und angewandte Mathematik,
Bd.~54.  S.~101--155.  1857.]}

\maketitle

\begin{center}
\large\bfseries
1.  Allgemeine Voraussetzungen und H\"{u}lfsmittel f\"{u}r
die Untersuchung von Functionen unbeschr\"{a}nkt
ver\"{a}nderlicher Gr\"{o}ssen.
\end{center}

Die Absicht den Lesern dieses Journals Untersuchungen \"{u}ber
verschiedene Transcendenten, inbesondere auch \"{u}ber
\emph{Abel}'sche Functionen vorzulegen, macht es mir
w\"{u}nschenswerth, um Wiederholungen zu vermeiden, eine
Zusammenstellung der allgemeinen Voraussetzungen, von denen ich
bei ihrer Behandlung ausgehen werde, in einem besonderen Aufsatze
voraufzuschicken.

F\"{u}r die unabh\"{a}ngig ver\"{a}nderliche Gr\"{o}sse setze ich
stets die jetzt allgemein bekannte \emph{Gauss}'sche
geometrische Repr\"{a}sentation voraus, nach welcher eine
complexe Gr\"{o}sse $z = x + yi$ vertreten wird durch einen Punkt
einer unendlichen Ebene, dessen rechtwinklige Coordinaten $x$, $y$
sind; ich werde dabei die complexen Gr\"{o}ssen und die sie
repr\"{a}sentirenden Punkte durch dieselben Buchstaben bezeichnen.
Als Function von $x + yi$ betrachte ich jede Gr\"{o}sse~$w$, die
sich mit ihr der Gleichung
\[ i \frac{\partial w}{\partial x}
   = \frac{\partial w}{\partial y} \]
gem\"{a}ss \"{a}ndert, ohne einen Ausdruck von $w$ durch $x$ und
$y$ vorauszusetzen.  Aus dieser Differentialgleichung folgt nach
einem bekannten Satze, dass die Gr\"{o}sse~$w$ durch eine nach
ganzen Potenzen von $z - a$ fortschreitende Reihe von der Form
$\sum\limits_{n=0}^{n = \infty} a_n (z - a)^n$
darstellbar ist, sobald sie in der Umgebung von $a$ allenthalben
\emph{einen} bestimmten mit $z$ stetig sich \"{a}ndernden Werth
hat, und dass diese Darstellbarkeit stattfindet bis zu einem
Abstande von $a$ oder Modul von $z - a$, f\"{u}r welchen eine
Unstetigkeit eintritt.  Es ergiebt sich aber aus den
Betrachtungen, welche der Methode der unbestimmten Coefficienten
zu Grunde liegen, dass die Coefficienten $a_n$ v\"{o}llig
bestimmt sind, \label{`Linie'} wenn $w$ in einer endlichen
\"{u}brigens beliebig kleinen von $a$ ausgehenden Linie gegeben
ist.

Beide Ueberlegungen verbindend, wird man sich leicht von der
Richtigkeit des Satzes \"{u}berzeugen:

{\itshape
Eine Function von $x + yi$, die in einem Theile der
$(x, y)$-Ebene gegeben ist, kann dar\"{u}ber hinaus nur auf Eine
Weise stetig fortgesetzt werden.}

Man denke sich nun die zu untersuchende Function nicht durch
irgend welche $z$ enthaltende analytische Ausdr\"{u}cke oder
Gleichungen bestimmt, sondern dadurch, dass der Werth der
Function in einem beliebig begrenzten Theile der $z$-Ebene
gegeben ist und sie von dort aus stetig (der partiellen
Differentialgleichung
\[ i \frac{\partial w}{\partial x}
   = \frac{\partial w}{\partial y} \]
gem\"{a}ss) fortgesetzt wird.  Diese Fortsetzung ist nach den
obigen S\"{a}tzen eine v\"{o}llig bestimmte, vorausgesetzt, dass
sie nicht in blossen Linien geschieht, wobei eine partielle
Differentialgleichung nicht zur Anwendung kommen k\"{o}nnte,
sondern durch Fl\"{a}chenstreifen von endlicher Breite.  Je
nach der Beschaffenheit der fortzusetzenden Function wird nun
entweder die Function f\"{u}r denselben Werth von $z$ immer
wieder denselben Werth annehmen, auf welchem Wege auch die
Fortsetzung geschehen sein m\"{o}ge, oder nicht.  Im ersteren
Falle nenne ich sie \emph{einwerthig}, sie bildet dann eine
f\"{u}r jeden Werth von $z$ v\"{o}llig bestimmte und nicht
l\"{a}ngs einer Linie unstetige Function.  Im letzteren Falle, wo
sie \emph{mehrwerthig} heissen soll, hat man um ihren Verlauf
aufzufassen vor Allem seine Aufmerksamkeit auf gewisse Punkte
der $z$-Ebene zu richten, um welche herum sich die Function in
eine andere fortsetzt.  Ein solcher Punkt ist z.~B.\ bei der
Function $\log (z - a)$ der Punkt~$a$.  Denkt man sich von diesem
Punkte $a$ aus eine beliebige Linie gezogen, so wird man in der
Umgebung von $a$ den Werth der Function so w\"{a}hlen k\"{o}nnen,
dass sie sich ausser dieser Linie \"{u}berall stetig \"{a}ndert;
zu beiden Seiten dieser Linie nimmt sie aber dann verschiedene
Werthe an, auf der negativen\footnote{\label{`positiv'}Im
Anschlusse an die von \emph{Gauss} vorgeschlagene Benennung
positiv laterale Einheit f\"{u}r $+i$ werde ich als positive
Seitenrichtung zu einer gegebenen Richtung diejenige bezeichnen,
welche zu ihr ebenso liegt, wie $+i$ zu $1$.}
einen um $2\pi i$ gr\"{o}sseren, als auf der positiven.  Die
Fortsetzung der Function von einer Seite dieser Linie aus,
z.~B.\ von der negativen, \"{u}ber sie hin\"{u}ber in das
jenseitige Gebiet giebt dann offenbar eine von der dort schon
vorhandenen verschiedene Function und zwar im hier betrachteten
Falle eine allenthalben um $2\pi i$ gr\"{o}ssere.

Zur bequemeren Bezeichnung dieser Verh\"{a}ltnisse sollen die
verschiedenen Fortsetzungen \emph{einer} Function f\"{u}r
denselben Theil der $z$-Ebene \emph{Zweige} dieser Function
genannt werden und ein Punkt, um welchen sich ein Zweig einer
Function in einen andern fortsetzt, eine
\emph{Verzweigungsstelle} dieser Function; wo keine Verzweigung
stattfindet, heisst die Function \emph{ein\"{a}ndrig} oder
\emph{monodrom}.

Ein Zweig einer Function von mehreren unabh\"{a}ngig
ver\"{a}nderlichen Gr\"{o}ssen, $z$, $s$, $t,\ldots$ ist
\emph{ein\"{a}ndrig} in der Umgebung eines bestimmten
Werthensystems
$z = a$, $s = b$, $t = c,\ldots$,
wenn allen Werthencombinationen bis zu einem endlichen Abstande
von demselben (oder bis zu einer bestimmten endlichen Gr\"{o}sse
der Moduln von
$z - a$, $s - b$, $t - c,\ldots$)
ein bestimmter mit den ver\"{a}nderlichen Gr\"{o}ssen stetig sich
\"{a}ndernder Werth dieses Zweiges der Function entspricht.  Eine
Verzweigungsstelle oder eine Stelle, um welche sich ein Zweig in
einen andern fortsetzt, wird bei einer Function von mehreren
Ver\"{a}nderlichen durch s\"{a}mmtliche einer Gleichung zwischen
ihnen gen\"{u}gende Werthe der unabh\"{a}ngig ver\"{a}nderlichen
Gr\"{o}ssen gebildet.

Nach einem oben angef\"{u}hrten bekannten Satze ist die
Ein\"{a}ndrigkeit einer Function identisch mit ihre
Entwickelbarkeit, ihre Verzweigung mit ihrer
Nichtentwickelbarkeit nach ganzen positiven oder negativen
Potenzen der Aenderungen der ver\"{a}nderlichen Gr\"{o}ssen.  Es
scheint aber nicht zweckm\"{a}ssig, jene von ihrer
Darstellungsweise unabh\"{a}ngigen Eigenschaften durch diese an
eine bestimmte Form ihres Ausdrucks gekn\"{u}pften Merkmale
auszudr\"{u}cken.

F\"{u}r manche Untersuchungen, namentlich f\"{u}r die
Untersuchung algebraischer und \emph{Abel}'scher Functionen ist
es vortheilhaft, die Verzweigungsart einer mehrwerthigen Function
in folgender Weise geometrisch darzustellen.  Man denke sich in
der $(x, y)$-Ebene eine andere mit ihr zusammenfallende Fl\"{a}che
\label{`Flaeche'}
(oder auf der Ebene einen unendlich d\"{u}nnen K\"{o}rper)
ausgebreitet, welche sich so weit und nur so weit erstreckt, als
die Function gegeben ist.  Bei Fortsetzung dieser Function wird
also diese Fl\"{a}che ebenfalls weiter ausgedehnt werden.  In
einem Theile der Ebene, f\"{u}r welchem zwei oder mehrere
Fortsetzungen der  Function vorhanden sind, wird die Fl\"{a}che
doppelt oder mehrfach sein; sie wird dort aus zwei oder mehreren
Bl\"{a}ttern bestehen, deren jedes einen Zweig der Function
vertritt.  Um einen Verzweigungspunkt der Function herum wird
sich ein Blatt der Fl\"{a}che in ein anderes fortsetzen, so dass
in der Umgebung eines solchen Punktes die Fl\"{a}che als eine
Schraubenfl\"{a}che mit einer in diesem Punkte auf der
$(x, y)$-Ebene senkrechten Axe und unendlich kleiner H\"{o}he des
Schraubenganges betrachtet werden kann.  Wenn die Function nach
mehreren Uml\"{a}ufen des $z$ um den Verzweigungswerth ihren
vorigen Werth wieder erh\"{a}lt
(wie z.~B.\ $(z - a)^{\frac{m}{n}}$,
wenn $m$, $n$ relative Primzahlen sind, nach $n$ Uml\"{a}ufen von
$z$ um $a$), muss man dann freilich annehmen, das sich das
oberste Blatt der Fl\"{a}che durch die \"{u}brigen hindurch in
das unterste fortsetzt.

Die mehrwerthige Function hat f\"{u}r jeden Punkt einer solchen
ihre Verzweigungsart darstellenden Fl\"{a}che nur \emph{einen}
bestimmten Werth und kann daher als eine v\"{o}llig bestimmte
Function des Orts in dieser Fl\"{a}che angesehen werden.

G\"{o}ttingen, 1857.

\begin{center}
\large\bfseries
2.  Lehrs\"{a}tze aus der analysis situs f\"{u}r die
Theorie der Integrale von zweigliedrigen vollst\"{a}ndigen
Differentialien.
\end{center}

Bei der Untersuchung der Functionen, welche aus der Integration
vollst\"{a}ndiger Differentialien entstehen, sind einige der
analysis situs angeh\"{o}rige S\"{a}tze fast unentbehrlich.  Mit
diesem von \emph{Leibnitz}, wenn auch vielleicht nicht ganz in
derselben Bedeutung, gebrauchten Namen darf wohl ein Theil der
Lehre von den stetigen Gr\"{o}ssen bezeichnet werden, welcher
die Gr\"{o}ssen nicht als unabh\"{a}ngig von der Lage existierend
und durch einander messbar betrachtet, sondern von den
Massverh\"{a}ltnissen ganz absehend, nur ihre Orts- und
Gebietsverh\"{a}ltnisse der Untersuchung unterwirft.  Indem ich
eine von Massverh\"{a}ltnissen ganz abstrahirende Behandlung
dieses Gegenstandes mir vorbehalte, werde ich hier nur die bei der
Integration zweigliedriger vollst\"{a}ndiger Differentialien
n\"{o}thigen S\"{a}tze in einem geometrischen Gewande darstellen.

Es sei eine in der $(x,y)$-Ebene einfach oder mehrfach
ausgebreitete Fl\"{a}che $T$ gegeben\footnote{Man sehe die
vorhergehende Abhandlung S.~\pageref{`Flaeche'}.}
und $X$, $Y$ seien solche stetige Functionen des Orts in dieser
Fl\"{a}che, dass in ihr allenthalben $X \, dx + Y \, dy$ ein
vollst\"{a}ndiges Differential, also
\[ \frac{\partial X}{\partial y} - \frac{\partial Y}{\partial x} = 0 \]
ist.  Bekanntlich ist dann
\[ \int (X \, dx + Y \, dy),\]
um einen Theil der Fl\"{a}che~$T$ positiv oder negativ
herum---d.~h.\ durch die ganze Begrenzung entweder allenthalben
nach der positiven oder allenthalben nach der negativen Seite
gegen die Richtung von Innen nach Aussen (siehe die Anmerkung 
Seite~\pageref{`positiv'} der vorhergehenden
Abhandlung)---erstreckt, $= 0$, da dies Integral dem \"{u}ber
diesen Theil ausgedehnten Fl\"{a}chenintegrale
\[ \int  \left(
            \frac{\partial Y}{\partial x}
          - \frac{\partial X}{\partial y}
         \right) \, dT \]
identisch im ersteren Falle gleich, im zweiten entgegengesetzt
ist.  Das Integral
\[ \int (X \, dx + Y \, dy) \]
hat daher zwischen zwei festen Punkten auf zwei verschiedenen
Wegen erstreckt denselben Werth, wenn diese beiden Wege
zusammengenommen die ganze Begrenzung eines Theils der
Fl\"{a}che~$T$ bilden.  Wenn also jede im Innern von $T$ in sich
zur\"{u}cklaufende Curve die ganze Begrenzung eines Theils von $T$
bildet, so hat das Integral von einem festen Anfangspunkte bis zu
einem und demselben Endpunkte erstreckt immer denselben Werth und
ist eine von dem Wege der Integration unabh\"{a}ngige
allenthalben in $T$ stetige Function von der Lage des Endpunkts.
Dies veranlasst zu einer Unterscheidung der Fl\"{a}chen in
einfach zusammenh\"{a}ngende, in welchen jede geschlossene Curve
einen Theil der Fl\"{a}che vollst\"{a}ndig begrenzt---wie
z.~B.\ ein Kreis---, und mehrfach zusammenh\"{a}ngende, f\"{u}r
welche dies nicht stattfindet,---wie z.~B.\ eine durch zwei
concentrische Kreise begrenzte Ringfl\"{a}che.  Eine mehrfach
zusammenh\"{a}ngende l\"{a}sst sich durch
Zerschneidung in eine einfach zusammenh\"{a}ngende verwandeln
(s.\ die durch Zeichnungen erl\"{a}uterten Beispiele am Schluss
dieser Abhandlung).  Da diese Operation wichtige Dienste bei der
Untersuchung der Integrale algebraischer Functionen leistet, so
sollen die darauf bez\"{u}glichen S\"{a}tze kurz zusammengestellt
werden; sie gelten f\"{u}r beliebig im Raume liegende
Fl\"{a}chen.

Wenn in einer Fl\"{a}che~$F$ zwei Curvensysteme $a$ und $b$
zusammengenommen einen Theil dieser Fl\"{a}che vollst\"{a}ndig
begrenzen, so bildet jedes andere Curvensystem, das mit $a$
zusammen einen Theil von $F$ vollst\"{a}ndig begrenzt, auch mit
$b$ die ganze Begrenzung eines Fl\"{a}chentheils, der aus den
beiden ersteren Fl\"{a}chentheilen l\"{a}ngs $a$ (durch Addition
oder Subtraction, jenachdem sie auf entgegengesetzter oder auf
gleicher Seite von $a$ liegen) zusammengesetzt ist.  Beide
Curvensysteme leisten daher f\"{u}r vollige Begrenzung eines
Theils von $F$ dasselbe und k\"{o}nnen f\"{u}r die Erf\"{u}llung
dieser Forderung einander ersetzen.

{\itshape
\label{`mehrfachzusammenhaengende'}
Wenn in einer
Fl\"{a}che sich $n$ geschlossene Curven
$a_1, a_2,\ldots, a_n$
ziehen lassen, welche weder f\"{u}r sich noch mit einander einen
Theil dieser Fl\"{a}che $F$ vollst\"{a}ndig begrenzen, mit deren
Zuziehung aber jede andere geschlossene Curve die
vollst\"{a}ndige Begrenzung eines Theils der Fl\"{a}che~$F$
bilden kann, so heisst die Fl\"{a}che eine ${(n + 1)} \,$\-fach
zusammenh\"{a}ngende.}

Dieser Charakter der Fl\"{a}che ist unabh\"{a}ngig von der Wahl
des Curvensystems $a_1, a_2,\ldots, a_n$, da je $n$ andere
geschlossene Curven $b_1, b_2,\ldots, b_n$, welche zu
v\"{o}lliger Begrenzung eines Theils dieser Fl\"{a}che nicht
ausreichen, ebenfalls mit jeder andern geschlossenen Curve
zusammengenommen einen Theil von $F$ v\"{o}llig begrenzen.

In der That, da $b_1$ mit Linien $a$ zusammengenommen einen Theil
von $F$ vollst\"{a}ndig begrenzt, so kann eine dieser Curven $a$
durch $b_1$ und die \"{u}brigen Curven $a$ ersetzt werden.  Es
ist daher mit $b_1$ und diesen $n - 1$ Curven $a$ jede andere
Curve, und folglich auch $b_2$, zu v\"{o}lliger Begrenzung eines
Theils von $F$ ausreichend, und es kann eine dieser $n - 1$
Curven $a$ durch $b_1$, $b_2$ und die \"{u}brigen $n - 2$ Curven
$a$ ersetzt werden.  Dieses Verfahren kann offenbar, wenn wie
vorausgesetzt die Curven $b$ zu vollst\"{a}ndiger Begrenzung
eines Theils von $F$ nicht ausreichen, so lange fortgesetzt
werden, bis s\"{a}mmtliche $a$ durch die $b$ ersetzt worden sind.

{\itshape
Eine ${(n + 1)} \,$\-fach zusammenh\"{a}ngende Fl\"{a}che~$F$ kann
durch einen Quer\-schnitt---d.~h.\ eine von einem
Begrenzungspunkte durch das Innere bis zu einem Begrenzungspunkte
gef\"{u}hrte Schnittlinie---in eine $n \,$\-fach
zusammenh\"{a}ngende $F'$ verwandelt werden.  Es gelten dabei die
durch die Zerschneidung entstehenden Begrenzungstheile schon
w\"{a}hrend der weiteren Zerschneidung als Begrenzung, so dass
ein Querschnitt keinen Punkt mehrfach durchschneiden, aber in
einem seiner fr\"{u}heren Punkte enden kann.}

Da die Linien $a_1, a_2,\ldots, a_n$ zu v\"{o}lliger Begrenzung
eines Theils von $F$ nicht ausreichen, so muss, wenn man sich $F$
durch diese Linien zerschnitten denkt, sowohl das auf der
rechten, als das auf der linken Seite von $a_n$ anliegende
Fl\"{a}chenst\"{u}ck noch andere von den Linien $a$ verschiedene
und also zur Begrenzung von $F$ geh\"{o}rige Begrenzungstheile
enthalten.  Man kann daher von einem Punkte von $a_n$ sowohl in
dem einen, als in dem andern dieser Fl\"{a}chenst\"{u}cke eine
die Curven $a$ nicht schneidende Linie bis zur Begrenzung von $F$
ziehen.  Diese beiden Linien $q'$ und $q''$ zusammengenommen
bilden alsdann einen Querschnitt $q$ der Fl\"{a}che $F$, welcher
das Verlangte leistet.

In der That sind in der durch diesen Querschnitt aus $F$
entstehenden Fl\"{a}che $F'$ die Linien
$a_1, a_2,\ldots, a_{n-1}$ im Innern von $F'$ verlaufende
geschlossene Curven, welche zur Begrenzung eines Theils von $F$,
also auch von $F'$, nicht hinreichen.  Jede andere im Innern von
$F'$ verlaufende geschlossene Curve~$l$ aber bildet mit ihnen die
ganze Begrenzung eines Theils von $F'$.  Denn die Linie~$l$
bildet mit einem Complex aus den Linien
$a_1, a_2,\ldots, a_n$ die ganze Begrenzung eines Theils $f$
von $F$.  Es l\"{a}sst sich aber zeigen, dass in der Begrenzung
desselben $a_n$ nicht vorkommen kann; denn dann w\"{u}rde, je
nachdem $f$ auf der linken oder rechten Seite von $a_n$ l\"{a}ge,
$q'$ oder $q''$ aus dem Innern von $f$ nach einem
Begrenzungspunkte von $F$, also nach einem ausserhalb $f$
gelegenen Punkte, f\"{u}hren und also die Begrenzung von $f$
schneiden m\"{u}ssen gegen die Voraussetzung, dass $l$ sowohl als
die Linien~$a$, den Durchschnittspunkt von $a_n$ und $q$
ausgenommen, stets im Innern von $F'$ bleiben.

Die Fl\"{a}che $F'$, in welche $F$ durch den Querschnitt $q$
zerf\"{a}llt, ist demnach, wie verlangt, eine $n \,$\-fach
zusammenh\"{a}ngende.

Es soll jetzt bewiesen werden, dass die Fl\"{a}che~$F$ durch
jeden Querschnitt~$p$, welcher sie nicht in getrennte St\"{u}cke
zerf\"{a}llt, in eine $n \,$\-fach zusammenh\"{a}ngende $F'$
verwandelt wird.  Wenn die zu beiden Seiten des Querschnitts~$p$
angrenzenden Fl\"{a}chentheile zusammenh\"{a}ngen, so l\"{a}sst
sich eine Linie $b$ von der einen Seite desselben durch das Innere
von $F'$ auf die andere Seite zum Anfangspunkte zur\"{u}ck
ziehen.  Diese Linie~$b$ bildet eine im Innern von $F$ in sich
zur\"{u}cklaufende Linie, welche, da der Querschnitt von ihr aus
nach beiden Seiten zu einem Begrenzungspunkte f\"{u}hrt, von
keinem der beiden Fl\"{a}chenst\"{u}cke, in welche sie $F$
zerschneidet, die ganze Begrenzung bildet.  Man kann daher eine
der Curven~$a$ durch die Curve~$b$ und jede der \"{u}brigen
$n - 1$ Curven $a$ durch eine im Innern von $F'$ verlaufende
Curve und wenn n\"{o}thig die Curve $b$ ersetzen, worauf der
Beweis, dass $F'$ $n \,$\-fach zusammenh\"{a}ngend ist, durch
dieselben Schl\"{u}sse, wie vorhin, gef\"{u}hrt werden kann.

{\itshape
Eine ${(n + 1)} \,$\-fach zusammenh\"{a}ngende Fl\"{a}che wird
daher durch jeden sie nicht in St\"{u}cke zerschneidenden
Querschnitt in eine $n \,$\-fach zusammenh\"{a}ngende verwandelt.}

Die durch einen Querschnitt entstandene Fl\"{a}che kann durch
einen neuen Querschnitt weiter zerlegt werden, und bei $n \,$\-maliger
Wiederholung dieser Operation wird eine ${(n + 1)} \,$\-fach
zusammenh\"{a}ngende Fl\"{a}che durch $n$ nach einander gemachte
sie nicht zerst\"{u}ckelnde Querschnitte in eine einfach
zusammenh\"{a}ngende verwandelt.

Um diese Betrachtungen auf eine Fl\"{a}che ohne Begrenzung, eine
geschlossene Fl\"{a}che, anwendbar zu machen, muss diese durch
Ausscheidung eines beliebigen Punktes in eine begrenzte
verwandelt werden, so dass die erste Zerlegung durch diesen Punkt
und einen in ihm anfangenden und endenden Querschnitt, also durch
eine geschlossene Curve, geschieht.  Die Oberfl\"{a}che eines
Ringes z.~B., welche eine dreifach zusammenh\"{a}ngende ist, wird
durch eine geschlossene Curve und einen Querschnitt in eine
einfach zusammenh\"{a}ngende verwandelt.

Auf das im Eingange betrachtete Integral des vollst\"{a}ndigen
Differentials $X \, dx + Y \, dy$ wird nun die eben behandelte
Zerschneidung der mehrfach zusammenh\"{a}ngenden Fl\"{a}chen in
einfach zusammenh\"{a}ngende, wie folgt, angewandt.  Ist die die
$(x, y)$-Ebene bedeckende Fl\"{a}che~$T$, in welcher $X$, $Y$
allenthalben stetige der Gleichung
\[ \frac{\partial X}{\partial y} - \frac{\partial Y}{\partial x}
   = 0 \]
gen\"{u}gende Functionen des Orts sind, $n \,$\-fach
zusammenh\"{a}ngend, so wird sie durch $n$ Querschnitte in eine
einfach zusammenh\"{a}ngende $T'$ zerschnitten.  Die Integration
von $X \, dx + Y \, dy$ von einem festen Anfangspunkte aus durch
Curven im Innern von $T'$ liefert dann einen nur von der Lage des
Endpunkts abh\"{a}ngigen Werth, welcher als Function von dessen
Coordinaten betrachtet werden kann.  Substituirt man f\"{u}r die
Coordinaten die Gr\"{o}ssen $x$, $y$, so erh\"{a}lt man eine
Function
\[ z = \int (X \, dx + Y \, dy) \]
von $x$, $y$, welche f\"{u}r jeden Punkt von $T'$ v\"{o}llig
bestimmt ist und sich innerhalb $T'$ allenthalben stetig, beim
Ueberschreiten eines Querschnitts aber allgemein zu reden um eine
endliche von einem Knotenpunkte des Schnittnetzes zum andern
constante Gr\"{o}sse \"{a}ndert.  Die Aenderungen beim
Ueberschreiten der Querschnitte sind von einer der Zahl der
Querschnitte gleichen Anzahl von einander unabh\"{a}ngiger
Gr\"{o}ssen abh\"{a}ngig; denn wenn man das Schnittsystem
r\"{u}ckw\"{a}rts,---die sp\"{a}teren Theile zuerst---,
durchl\"{a}uft, so ist diese Aenderung \"{u}berall bestimmt, wenn
ihr Werth beim Beginn jedes Querschnitts gegeben wird; letztere
Werthe aber sind von einander unabh\"{a}ngig.

G\"{o}ttingen, 1857.

\bigbreak

\centerline{\vbox{\hrule width 144 pt}}

\nobreak\bigskip

Um das, was oben (S.~\pageref{`mehrfachzusammenhaengende'}) unter
einer $n \,$\-fach zusammenh\"{a}ngenden Fl\"{a}che verstanden wird,
anschaulicher zu machen, folgen in den Zeichnungen auf nachstehenden
Seiten Beispiele von einfach, zweifach und dreifach
zusammenh\"{a}ngenden Fl\"{a}chen.

\newpage

\centerline{\bfseries Einfach zusammenh\"{a}ngende Fl\"{a}che.}

\begin{figure}[h]
\centerline{\epsfbox{abfnfig.1}}
\end{figure}

Sie wird durch jeden Querschnitt in getrennte St\"{u}cke
zerf\"{a}llt, und es bildet in ihr jede geschlossene Curve die
ganze Begrenzung eines Theils der Fl\"{a}che.

\medskip

\centerline{\bfseries Zweifach zusammenh\"{a}ngende Fl\"{a}che.}

\begin{figure}[h]
\centerline{\epsfbox{abfnfig.2}}
\end{figure}

Sie wird durch jeden sie nicht zerst\"{u}ckelnden Querschnitt $q$
in eine einfach zusammenh\"{a}ngende zerschnitten.  Mit Zuziehung
der Curve $a$ kann in ihr jede geschlossene Curve die ganze
Begrenzung eines Theils der Fl\"{a}che bilden.

\newpage

\centerline{\bfseries Dreifach zusammenh\"{a}ngende Fl\"{a}che.}

\begin{figure}[h]
\centerline{\epsfbox{abfnfig.3}}
\end{figure}

In dieser Fl\"{a}che kann jede geschlossene Curve mit Zuziehung
der Curven $a_1$ und $a_2$ die ganze Begrenzung eines Theils der
Fl\"{a}che bilden.  Sie zerf\"{a}llt durch jeden sie nicht
zerst\"{u}ckelnden Querschnitt in eine zweifach
zusammenh\"{a}ngende und durch zwei solche Querschnitte, $q_1$
und $q_2$, in eine einfach zusammenh\"{a}ngende.


\begin{figure}[h]
\centerline{\epsfbox{abfnfig.4}}
\end{figure}

In dem Theile $\alpha$~$\beta$~$\gamma$~$\delta$ der Ebene ist
die Fl\"{a}che doppelt.  Der $a_1$ enthaltende Arm der Fl\"{a}che
ist als unter dem andern fortgehend betrachtet und daher durch
punktirte Linien angedeutet.

\newpage

\begin{center}
\large\bfseries
3.  Bestimmung einer Function einer ver\"{a}nderlichen
complexen Gr\"{o}sse durch Grenz- und Unstetigkeitsbedingungen.
\end{center}

Wenn in einer Ebene, in welcher die rechtwinkligen Coordinaten
eines Punkts $x$, $y$ sind, der Werth einer Function von $x + yi$
in einer endlichen Linie gegeben ist, so kann diese von dort aus
nur auf eine Weise stetig fortgesetzt werden und ist also dadurch
v\"{o}llig bestimmt (siehe oben S.~\pageref{`Linie'}).  Sie kann
aber auch in dieser Linie nicht willk\"{u}rlich angenommen
werden, wenn sie von ihr aus einer stetigen Fortsetzung in die
anstossenden Fl\"{a}chentheile nach beiden Seiten hin f\"{a}hig
sein soll, da sie durch ihren Verlauf in einem noch so kleinen
endlichen Theile dieser Linie schon f\"{u}r den \"{u}brigen Theil
bestimmt ist.  Bei dieser Bestimmungsweise einer Function sind
also die zu ihrer Bestimmung dienenden Bedingungen nicht von
einander unabh\"{a}ngig.

Als Grundlage f\"{u}r die Untersuchung einer Transcendenten ist
es vor allen Dingen n\"{o}thig, ein System zu ihrer Bestimmung
hinreichender von einander unabh\"{a}ngiger Bedingungen
aufzustellen.  Hierzu kann in vielen F\"{a}llen, namentlich bei
den Integralen algebraischer Functionen und ihren inversen
Functionen, ein Princip dienen, welches \emph{Dirichlet} zur
L\"{o}sung dieser Aufgabe f\"{u}r eine der \emph{Laplace}'schen
partiellen Differentialgleichung gen\"{u}gende Function von drei
Ver\"{a}nderlichen,---wohl durch einen \"{a}hnlichen Gedanken von
\emph{Gauss} veranlasst---in seinen Vorlesungen \"{u}ber die
dem umgekehrten Quadrat der Entfernung proportional wirkenden
Kr\"{a}fte seit einer Reihe von Jahren zu geben pflegt.  F\"{u}r
diese Anwendung auf die Theorie von Transcendenten ist jedoch
gerade ein Fall besonders wichtig, auf welchen dies Princip in
seiner dortigen einfachsten Form nicht anwendbar ist, und welcher
dort als von ganz untergeordneter Bedeutung unber\"{u}cksichtigt
bleiben kann.  Dieser Fall ist der, wenn die Function an gewissen
Stellen des Gebiets, wo sie zu bestimmen ist, vorgeschriebene
Unstetigkeiten annehmen soll; was so zu verstehen ist, dass sie
an jeder solchen Stelle der Bedingung unterworfen ist, unstetig
zu werden, wie eine dort gegebene unstetige Function, oder sich
nur um eine dort stetige Function von ihr zu unterscheiden.  Ich
werde hier das Prinzip in der f\"{u}r die beabsichtigte Anwendung
erforderlichen Form darstellen und erlaube mir dabei in Betreff
einiger Nebenuntersuchungen auf die in meiner Doctordissertation
(Grundlagen f\"{u}r eine allgemeine Theorie der Functionen einer
ver\"{a}nderlichen complexen Gr\"{o}sse.   G\"{o}ttingen 1851)
gegebene Darstellung desselben zu verweisen.

Man nehme an, dass eine die $(x,y)$-Ebene einfach oder mehrfach
bedeckende beliebig begrenzte Fl\"{a}che~$T$ und in derselben
zwei f\"{u}r jeden ihrer Punkte eindeuting bestimmte reelle
Functionen von $x$, $y$, die Functionen $\alpha$ und $\beta$
gegeben seien, und bezeichne das durch die Fl\"{a}che ausgedehnte
Integral
\[ \int  \left(
            \left(
               \frac{\partial \alpha}{\partial x}
             - \frac{\partial \beta}{\partial y}
            \right)^2
          + \left(
               \frac{\partial \alpha}{\partial y}
             + \frac{\partial \beta}{\partial x}
            \right)^2
         \right) \, dT \]
durch $\Omega(\alpha)$, wobei die Functionen $\alpha$ und $\beta$
beliebige Unstetigkeiten besitzen k\"{o}nnen, wenn nur das
Integral dadurch nicht unendlich wird.  Es bleibt dann auch
$\Omega (\alpha - \lambda)$ endlich, wenn $\lambda$ allenthalben
stetig ist und endliche Differentialquotienten hat.  Wird diese
stetige Function~$\lambda$ der Bedingung unterworfen, nur in
einem unendlich kleinen Theile der Fl\"{a}che~$T$ von einer
unstetigen Function~$\gamma$ verschieden zu sein, so wird
$\Omega (\alpha - \lambda)$ unendlich gross, wenn $\gamma$
l\"{a}ngs einer Linie unstetig ist oder in einem Punkte so
unstetig ist, dass
\[ \int  \left(
            \left( \frac{\partial \gamma}{\partial x} \right)^2
          + \left( \frac{\partial \gamma}{\partial y} \right)^2
         \right) \, dT \]
unendlich wird (Meine Inaug.\ Diss.\ Art.\ 17); es bleibt aber
$\Omega (\alpha - \lambda)$ endlich, wenn $\gamma$ nur in
einzelnen Punkten und nur so unstetig ist, dass
\[ \int  \left(
            \left( \frac{\partial \gamma}{\partial x} \right)^2
          + \left( \frac{\partial \gamma}{\partial y} \right)^2
         \right) \, dT \]
durch die Fl\"{a}che $T$ erstreckt endlich bleibt, wie
z.~B.\ wenn $\gamma$ in der Umgebung eines Punktes im
Abstande~$r$ von demselben $= (- \log r)^\varepsilon$ und
$0 < \varepsilon < \frac{1}{2}$ ist.  Zur Abk\"{u}rzung m\"{o}gen
hier die Functionen, in welche $\lambda$ unbeschadet der
Endlichkeit von $\Omega (\alpha - \lambda)$ \"{u}bergehen kann,
unstetig von der ersten Art, die Functionen, f\"{u}r welche dies
nicht m\"{o}glich ist, unstetig von der zweiten Art genannt
werden.  Denkt man sich nun in $\Omega (\alpha - \mu)$ f\"{u}r
$\mu$ alle stetigen oder von der ersten Art unstetigen Functionen
gesetzt, welche an der Grenze verschwinden, so erh\"{a}lt dies
Integral immer einen endlichen, aber seiner Natur nach nie einen
negativen Werth, und es muss daher wenigstens einmal, f\"{u}r
$\alpha - \mu = u$, ein Minimumwerth eintreten, so dass $\Omega$
f\"{u}r jede Function $\alpha - \mu$, die unendlich wenig von $u$
verschieden ist, gr\"{o}sser als $\Omega(u)$ wird.

Bezeichnet daher $\sigma$ eine beliebige stetige oder von erster
Art unstetige Function des Orts in der Fl\"{a}che $T$, die an der
Grenze allenthalben gleich $0$ ist, und $h$ eine von $x$, $y$
unabh\"{a}ngige Gr\"{o}sse, so muss $\Omega( u + h \sigma)$
sowohl f\"{u}r ein positives, als f\"{u}r ein negatives
hinreichend kleines $h$ gr\"{o}sser als $\Omega(u)$ werden, und
daher in der Entwicklung dieses Ausdrucks nach Potenzen von $h$
der Coefficient von $h$ verschwinden.  Ist dieser $0$, so ist
\[ \Omega (u + h \sigma)
   = \Omega(u) + h^2 \int
         \left(
            \left( \frac{\partial \sigma}{\partial x} \right)^2
          + \left( \frac{\partial \sigma}{\partial y} \right)^2
         \right) \, dT \]
und folglich $\Omega$ immer ein Minimum.  Das Minimum tritt nur f\"{u}r
eine einzige Function $u$ ein; denn f\"{a}nde auch ein Minimum
f\"{u}r $u + \sigma$ statt, so k\"{o}nnte $\Omega( u + \sigma)$
nicht $> \Omega(u)$ sein, weil sonst
\[ \Omega(u + h \sigma) < \Omega (u + \sigma) \]
f\"{u}r $h < 1$ w\"{u}rde; also k\"{o}nnte $\Omega(u + \sigma)$
nicht kleiner als die anliegenden Werthe sein.  Ist aber
$\Omega(u + \sigma) = \Omega(u)$, so muss $\sigma$ constant, also
da es in der Begrenzung~$0$ ist, \"{u}berall $0$ sein.  Es wird
daher nur f\"{u}r eine einzige Function~$u$ das Integral $\Omega$
ein Minimum und die Variation erster Ordnung oder das
$h$~proportionale Glied in $\Omega(u + h \sigma)$,
\[ 2 h \int dT \,
         \left(
            \left(
               \frac{\partial u}{\partial x}
             - \frac{\partial \beta}{\partial y}
            \right)
            \frac{\partial \sigma}{\partial x}
          + \left(
               \frac{\partial u}{\partial y}
             + \frac{\partial \beta}{\partial x}
            \right)
            \frac{\partial \sigma}{\partial y}
         \right)
   = 0.\]

Aus dieser Gleichung folgt, dass das Integral
\[ \int  \left(
            \left(
               \frac{\partial \beta}{\partial x}
             + \frac{\partial u}{\partial y}
            \right) \, dx
          + \left(
               \frac{\partial \beta}{\partial y}
             - \frac{\partial u}{\partial x}
            \right) \, dy
         \right) \]
durch die ganze Begrenzung eines Theils der Fl\"{a}che~$T$
erstreckt stets $= 0$ ist. Zerlegt man nun (nach der
vorhergehenden Abhandlung) die Fl\"{a}che~$T$, wenn sie eine
mehrfach zusammenh\"{a}ngende ist, in eine einfach
zusammenh\"{a}ngende $T'$, so liefert die Integration durch das
Innere von $T'$ von einem festen Anfangspunkte bis zum Punkte
$(x,y)$ eine Function von $x$,~$y$,
\[ \nu = \int
         \left(
            \left(
               \frac{\partial \beta}{\partial x}
             + \frac{\partial u}{\partial y}
            \right) \, dx
          + \left(
               \frac{\partial \beta}{\partial y}
             - \frac{\partial u}{\partial x}
            \right) \, dy
         \right)
      + \mbox{const.},\]
welche in $T'$ \"{u}berall stetig oder unstetig von der ersten
Art ist und sich beim Ueberschreiten der Querschnitte um endliche
von einem Knotenpunkte des Schnittnetzes zum andern constante
Gr\"{o}ssen \"{a}ndert.  Es gen\"{u}gt dann $v = \beta - \nu$ den
Gleichungen
\[     \frac{\partial v}{\partial x}
   = - \frac{\partial u}{\partial y},\quad
       \frac{\partial v}{\partial y}
   =   \frac{\partial u}{\partial x},\]
und folglich ist $u + vi$ eine L\"{o}sung der
Differentialgleichung
\[      \frac{\partial (u + vi)}{\partial y}
    - i \frac{\partial (u + vi)}{\partial x} = 0 \]
oder eine Function von $x + yi$.

Man erh\"{a}lt auf diesem Wege den in der erw\"{a}hnten
Abhandlung Art.~18 ausgesprochenen Satz:

{\itshape
\label{`Bestimmung'}
Ist in einer zusammenh\"{a}ngenden durch Querschnitte
in eine einfach zusammenh\"{a}ngende $T'$ zerlegten Fl\"{a}che
$T$ eine complexe Function $\alpha + \beta i$ von $x$ und $y$
gegeben, f\"{u}r welche
\[ \int  \left(
            \left(
               \frac{\partial \alpha}{\partial x}
             - \frac{\partial \beta}{\partial y}
            \right)^2
          + \left(
               \frac{\partial \alpha}{\partial y}
             + \frac{\partial \beta}{\partial x}
            \right)^2
         \right) \, dT \]
durch die ganze Fl\"{a}che ausgedehnt einen endlichen Werth hat,
so kann sie immer und nur auf Eine Art in eine Function von
$x + yi$ verwandelt werden durch Subtraction einer Function
$\mu + \nu i$ von $x$,~$y$, welche folgenden Bedingungen
gen\"{u}gt:

1) $\mu$ ist am Rande $= 0$ oder doch nur in einzelnen Punkten
davon verschieden, $\nu$ in Einem Punkte beliebig gegeben.

2) Die Aenderungen von $\mu$ sind in $T$, von $\nu$ in $T'$ nur
in einzelnen Punkten und nur so unstetig, dass
\[ \int  \left(
            \left( \frac{\partial \mu}{\partial x} \right)^2
          + \left( \frac{\partial \mu}{\partial y} \right)^2
         \right) \, dT \]
und
\[ \int  \left(
            \left( \frac{\partial \nu}{\partial x} \right)^2
          + \left( \frac{\partial \nu}{\partial y} \right)^2
         \right) \, dT,\]
durch die ganze Fl\"{a}che erstreckt, endlich bleiben, und
letztere l\"{a}ngs der Querschnitte beiderseits gleich.}

Wenn die Function $\alpha + \beta i$, wo ihre
Differentialquotienten unendlich werden, unstetig wird, wie eine
gegebene dort unstetige Function von $x + yi$, und keine durch
eine Ab\"{a}nderung ihres Werthes in einem einzelnen Punkte
hebbare Unstetigkeit besitzt, so bleibt $\Omega(\alpha)$ endlich,
und es wird $\mu + \nu i$ in $T'$ allenthalben stetig.  Denn da
eine Function von $x + yi$ gewisse Unstetigkeiten, wie z.~B.
Unstetigkeiten erster Art, gar nicht annehmen kann
(Meine Diss.\ Art.\ 12), so muss die Differenz zweier solcher
Functionen stetig sein, sobald sie nicht von der zweiten Art
unstetig ist.

Nach den eben bewiesenen Satze l\"{a}sst sich daher eine Function
von $x + yi$ so bestimmen, dass sie im Innern von $T$, von der
Unstetigkeit des imagin\"{a}ren Theils in den Querschnitten
abgesehen, gegebene Unstetigkeiten annimmt, und ihr reeller Theil
and der Grenze einen dort allenthalben beliebig gegebenen Werth
erh\"{a}lt; wenn nur f\"{u}r jeden Punkt, wo ihre
Differentialquotienten unendlich werden sollen, die
vorgeschriebene Unstetigkeit die einer gegebenen dort unstetigen
Function von $x + yi$ ist.  Die Bedingung an der Grenze kann man,
wie leicht zu sehen, ohne eine wesentliche Aenderung der
gemachten Schl\"{u}sse durch manche andere ersetzen.

G\"{o}ttingen, 1857.

\begin{center}
\large\bfseries
4.  Theorie der \emph{Abel}'schen Functionen.
\end{center}

In der folgenden Abhandlung habe ich die \emph{Abel}'schen
Functionen nach einer Methode behandelt, deren Principien in
meiner Inauguraldissertation\footnote{Grundlagen f\"{u}r eine
allgemeine Theorie der Functionen einer ver\"{a}nderlichen
complexen Gr\"{o}sse.  G\"{o}ttingen 1851.}
aufgestellt und in einer etwas ver\"{a}nderten Form in den drei
vorhergehenden Aufs\"{a}tzen dargestellt worden sind.  Zur
Erleichterung der Uebersicht schicke ich eine kurze Inhaltsangabe
vorauf.

Die erster Abtheilung enth\"{a}lt die Theorie eines Systems von
gleichverzweigten algebraischen Functionen und ihren Integralen,
soweit f\"{u}r dieselbe nicht die Betrachtung von
$\vartheta$-Reihen massgebend ist, und handelt im \S.~1--5 von
der Bestimmung dieser Functionen durch ihre Verzweigungsart und
ihre Unstetigkeiten, im \S.~6--10 von den rationalen
Ausdr\"{u}cken derselben in zwei durch eine algebraische
Gleichung verkn\"{u}pfte ver\"{a}nderliche Gr\"{o}ssen; und im
\S.~11--13 von der Transformation dieser Ausdr\"{u}cke durch
rationale Substitutionen.  Der bei dieser Untersuchung sich
darbietende Begriff eine \emph{Klasse} von algebraischen
Gleichungen, welche sich durch rationale Substitutionen in
einander transformiren lassen, d\"{u}rfte auch f\"{u}r andere
Untersuchungen wichtig und die Transformation einer solchen
Gleichung in Gleichungen niedrigsten Grades ihrer Klasse (\S.~13)
auch bei anderen Gelegenheiten von Nutzen sein.  Diese Abtheilung
behandelt endlich im \S.~14--16 zur Vorbereitung der folgenden
die Anwendung des \emph{Abel}'schen Additionstheorems f\"{u}r
ein beliebiges System allenthalben endlicher Integrale von
gleichverzweigten algebraischen Functionen zur Integration eines
Systems von Differentialgleichungen.

In der zweiter Abtheilung werden f\"{u}r ein beliebiges System
von immer endlichen Integralen gleichverzweigter, algebraischer,
$\overline{2p + 1} \,$\-fach zusammenh\"{a}ngender Functionen die
\emph{Jacobi}'schen Umkehrungsfunctionen von $p$
ver\"{a}nderlichen Gr\"{o}ssen durch $p \,$\-fach unendliche
$\vartheta$-Reihen ausgedr\"{u}ckt, d.~h.\ durch Reihen von der
Form
\[ \vartheta(v_1, v_2,\ldots, v_p)
   = \left( \sum_{-\infty}^\infty \right)^p
      e^{\left( \sum\limits_1^p \right)^2
               a_{\mu, \mu'} m_\mu m_\mu'
         + 2 \sum\limits_1^p v_\mu m_\mu},\]
worin die Summationen im Exponenten sich auf $\mu$ und $\mu'$,
die \"{a}usseren Summationen auf $m_1, m_2,\ldots, m_p$ beziehen.
Es ergiebt sich, dass zur allgemeinen L\"{o}sung dieser Aufgabe
eine---wenn $p > 3$ specielle---Gattung von $\vartheta$-Reihen
ausreicht, in denen zwischen den
$\displaystyle \frac{p(p + 1)}{2}$
Gr\"{o}ssen
\[ \displaystyle a \frac{(p - 2)(p - 3)}{1 \mathbin{.} 2} \]
Relationen stattfinden, so dass nur $3p - 3$ willk\"{u}rlich
bleiben.  Dieser Theil der Abhandlung bildet zugleich eine
Theorie dieser speciellen Gattung von $\vartheta$-Functionen; die
allgemeinen $\vartheta$-Functionen bleiben hier ausgeschlossen,
lassen sich jedoch nach einer ganz \"{a}hnlichen Methode
behandeln.

Das hier erledigte \emph{Jacobi}'sche Umkehrungsproblem ist
f\"{u}r die hyperelliptischen Integrale schon auf mehreren Wegen
durch die beharrlichen mit so sch\"{o}nem Erfolge gekr\"{o}nten
Arbeiten von \emph{Weierstrass} gel\"{o}st worden, von denen
eine Uebersicht in 47.~Bande d.~J.\ (S.~289)
mitgetheilt worden ist.  Es ist jedoch bis jetzt nur von
dem Theile dieser Arbeiten, welcher in den \S\S.~1 und 2 und der
ersten die elliptischen Functionen betreffenden H\"{a}lfte des
\S~3 der angef\"{u}hrten Abhandlung skizzirt wird, die wirkliche
Ausf\"{u}hrung ver\"{o}ffentlicht
(Bd.~52, S.~285 d.~J.);
in wie weit zwischen den sp\"{a}teren Theilen dieser Arbeiten und
meinen hier dargestellten eine Uebereinstimmung nicht bloss in
Resultaten, sondern auch in den zu ihnen f\"{u}hrenden Methoden
stattfindet, wird grossentheils erst die versprochene
ausf\"{u}hrliche Darstellung derselben ergeben k\"{o}nnen.

Die gegenw\"{a}rtige Abhandlung bildet mit Ausnahme der beiden
letzten \S\S.~26 und 27, deren Gegenstand damals nur kurz
angedeutet werden konnte, einen Auszug aus einem Theile meiner
von Michaelis 1855 bis Michaelis 1856 zu G\"{o}ttingen gehaltenen
Vorlesungen.  Was die Auffindung der einzelnen Resultate
betrifft, so wurde ich auf das im \S.~1--5, 9 und 12 Mitgetheilte
und die dazu n\"{o}thigen vorbereitenden S\"{a}tze, welche
sp\"{a}ter Behufs der Vorlesungen so, wie es in dieser Abhandlung
geschehen ist, weiter ausgef\"{u}hrt wurden, im Herbste 1851 und
zu Anfang 1852 durch Untersuchungen \"{u}ber die conforme
Abbildung mehrfach zusammenh\"{a}ngender Fl\"{a}chen gef\"{u}hrt,
ward aber dann durch einen andern Gegenstand von dieser
Untersuchung abgezogen.  Erst um Ostern 1855 wurde sie wieder
aufgenommen und in den Oster- und Michaelisferien jenes Jahres
bis zu \S.~21 incl.\ fortgef\"{u}hrt; das Uebrige wurde bis
Michaelis 1856 hinzugef\"{u}gt.  Einzelne erg\"{a}nzende
Zus\"{a}tze sind an manchen Stellen w\"{a}hrend der Ausarbeitung
hinzugekommen.

\begin{center}
\large\bfseries
Erste Abtheilung.
\end{center}
 
\centerline{1.}

\nobreak\medskip

Ist $s$ die Wurzel einer irreductibeln Gleichung $n \,$\-ten Grades,
deren Coefficienten ganze Functionen $m \,$\-ten Grades von $z$
sind, so entsprechen jedem Werthe von $z$ $n$~Werthe von $s$, die
sich mit $z$ \"{u}berall, wo sie nicht unendlich werden, stetig
\"{a}ndern.  Stellt man daher (nach S.~\pageref{`Flaeche'}) die
Verzweigungsart dieser Function durch eine in der $z$-Ebene
ausgebreitete unbegrenzte Fl\"{a}che $T$ dar, so ist diese in
jedem Theile der Ebene $n \,$\-fach, und $s$ ist dann eine
einwerthige Function des Orts in dieser Fl\"{a}che.  Eine
unbegrenzte Fl\"{a}che kann entweder als eine Fl\"{a}che mit
unendlich weit entfernter Begrenzung oder als eine geschlossene
angesehen werden, und Letzteres soll bei der Fl\"{a}che~$T$
geschehen, so dass dem Werthe $z = \infty$ in jedem der $n$
Bl\"{a}tter der Fl\"{a}che \emph{ein} Punkt entspricht, wenn
nicht etwa f\"{u}r $z = \infty$ eine Verzweigung stattfindet.

Jede rationale Function von $s$ und $z$ ist offenbar ebenfalls
eine einwerthige Function des Orts in der Fl\"{a}che~$T$ und
besitzt also dieselbe Verzweigungsart wie die Function~$s$, und
es wird sich unten ergeben, dass auch das Umgekehrte gilt.

Durch Integration einer solchen Function erh\"{a}lt man eine
Function, deren verschiedene Fortsetzungen f\"{u}r denselben
Theil der Fl\"{a}che~$T$ sich nur um Constanten unterscheiden, da
ihre Derivirte f\"{u}r denselben Punkt dieser Fl\"{a}che immer
denselben Werth wieder annimmt.

Ein solches System von gleichverzweigten algebraischen Functionen
und Integralen dieser Functionen bildet zun\"{a}chst den
Gegenstand unserer Betrachtung; statt aber von diesen
Ausdr\"{u}cken dieser Functionen auszugehen, werden wir sie mit
Anwendung des \emph{Dirichlet}'schen Princips
(S.~\pageref{`Bestimmung'}) durch ihre Unstetigkeiten definiren.

\medbreak

\centerline{2.}

\nobreak\medskip

Zur Vereinfachung des Folgenden heisse eine Function
\emph{f\"{u}r einen Punkt der Fl\"{a}che $T$ unendlich klein von
der ersten Ordnung}, wenn ihr Logarithmus bei einem positiven
Umlaufe um ein diesen Punkt umgebendes Fl\"{a}chenst\"{u}ck, in
welchem sie endlich und von Null verschiedenen bleibt, um
$2\pi i$ w\"{a}chst.  Es ist demnach f\"{u}r einen Punkt, um den
die Fl\"{a}che~$T$ sich $\mu$ mal windet, wenn dort $z$ einen
endlichen Werth~$a$ hat,
$\displaystyle (z - a)^{\frac{1}{\mu}}$,
also
$\displaystyle (dz)^{\frac{1}{\mu}}$,
wenn aber $z = \infty$,
$\displaystyle \left( \frac{1}{z} \right)^{\frac{1}{\mu}}$
unendlich klein von der ersten Ordnung.  Der Fall, wo eine
Function in einem Punkte der Fl\"{a}che~$T$ unendlich klein oder
unendlich gross von der $\nu \,$\-ten Ordnung wird, kann so
betrachtet werden, als wenn die Function in $\nu$ dort
zusammenfallenden (oder unendlich nahen) Punkten unendlich klein
oder unendlich gross von der ersten Ordnung wird, wie in der
Folge bisweilen geschehen soll.

Die Art und Weise, wie jene hier zu betrachtenden Functionen
unstetig werden, kann so ausgedr\"{u}ckt werden.  Wird eine von
ihnen in einem Punkte der Fl\"{a}che $T$ unendlich, so kann sie,
wenn $r$ eine beliebige Function bezeichnet, die in diesem Punkte
unendlich klein von der ersten Ordnung wird, stets durch
Subtraction eines endlichen Ausdrucks von der Form
\[ A \log r + B r^{-1} + C r^{-2} + \cdots \]
in eine dort stetige verwandelt werden, wie sich aus den
bekannten---nach \emph{Cauchy} oder durch die
\emph{Fourier}'sche Reihe zu beweisenden---S\"{a}tzen \"{u}ber
die Entwicklung einer Function in Potenzreihen ergiebt.

\medbreak

\centerline{3.}

\nobreak\medskip

Man denke sich jetzt eine in der $z$-Ebene allenthalben $n \,$\-fach
ausgebreitete unbegrenzte und nach dem Obigen als geschlossen zu
betrachtende zusammenh\"{a}ngende Fl\"{a}che~$T$ gegeben und
diese in eine einfach zusammenh\"{a}ngende $T'$ zerschnitten.  Da
die Begrenzung einer einfach zusammenh\"{a}ngenden Fl\"{a}che aus
\emph{einem} St\"{u}cke besteht, eine geschlossene Fl\"{a}che
aber durch eine ungerade Anzahl von Schnitten eine gerade Zahl
von Begrenzungsst\"{u}cken, durch eine gerade eine ungerade
erh\"{a}lt, so ist zu dieser Zerschneidung eine gerade Anzahl von
Schnitten erforderlich.  Die Anzahl dieser Querschnitte sei
$= 2p$.  Die Zerschneidung werde zur Vereinfachung des Folgenden
so ausgef\"{u}hrt, dass jeder sp\"{a}tere Schnitt von einem
Punkte eines fr\"{u}heren bis zu dem anstossenden Punkte auf der
anderen Seite desselben geht: wenn sich dann eine Gr\"{o}sse
l\"{a}ngs der ganzen Begrenzung von $T'$ stetig \"{a}ndert und im
ganzen Schnittsysteme zu beiden Seiten gleiche Aenderungen
erleidet, so ist die Differenz der beiden Werthe, die sie in
demselben Punkte des Schnittnetzes annimmt, in allen Theilen
\emph{eines} Querschnitts derselben Constanten gleich.

Man setze nun $z = x + yi$ und nehme in $T$ eine Function
$\alpha + \beta i$ von $x$,~$y$ folgendermassen an:

In der Umgebung der Punkte $\varepsilon_1, \varepsilon_2,\ldots$
bestimmte man sie gleich gegebenen in diesen Punkten unendlich
werdenden Functionen von $x + yi$, und zwar um $\varepsilon_\nu$,
indem man eine beliebige Function von $z$, die in
$\varepsilon_\nu$ unendlich klein von der ersten Ordnung wird,
durch $r_\nu$ bezeichnet, gleich einem endlichen Ausdrucke von
der Form
\[ A_\nu \log r_\nu + B_\nu r_\nu^{-1} + C_\nu r_\nu^{-2} + \cdots
   = \varphi_\nu(r_\nu),\]
worin $A_\nu, B_\nu, C_\nu,\ldots$ willk\"{u}rliche Constanten
sind.  Man ziehe ferner nach einem beliebigen Punkte von allen
Punkten $\varepsilon$, f\"{u}r welche die Gr\"{o}sse~$A$ von Null
verschieden ist, einander nicht schneidende Linien durch das
Innere von $T'$, von $\varepsilon_\nu$ die Linie $l_\nu$.  Man
nehme endlich die Function in der ganzen noch \"{u}brigen
Fl\"{a}che~$T$ so an, dass sie ausser den Linien~$l$ und den
Querschnitten \"{u}berall stetig, auf der positiven (linken)
Seite der Linie~$l_\nu$ um $-2\pi i A_\nu$ und auf der positiven
Seite des $\nu \,$\-ten Querschnitts um die gegebene Constante
$h^{(\nu)}$ gr\"{o}sser ist, als auf der andern, und dass das
Integral
\[ \int  \left(
            \left(
               \frac{\partial \alpha}{\partial x}
             - \frac{\partial \beta}{\partial y}
            \right)^2
          + \left(
               \frac{\partial \alpha}{\partial y}
             + \frac{\partial \beta}{\partial x}
            \right)^2
         \right) \, dT \]
durch die Fl\"{a}che~$T$ ausgedehnt einen endlichen Werth
erh\"{a}lt.  Dies ist wie leicht zu sehen immer m\"{o}glich, wenn
die Summe s\"{a}mmtlicher Gr\"{o}ssen~$A$ gleich Null ist, aber
auch nur unter dieser Bedingung, weil nur dann die Function nach
einem Umlaufe um das System der Linien~$l$ den vorigen Werth
wieder annehmen kann.

Die Constanten $h^{(1)}, h^{(2)},\ldots, h^{(2p)}$, um welche
eine solche Function auf der positiven Seite der Querschnitte
gr\"{o}sser ist, als auf der andern, sollen die
\emph{Periodicit\"{a}tsmoduln} dieser Function genannt werden.

Nach dem \emph{Dirichlet}'schen Princip kann nun die Function
$\alpha + \beta i$ in eine Function $\omega$ von $x + yi$
verwandelt werden durch Subtraction einer \"{a}hnlichen in $T'$
allenthalben stetigen Function von $x$,~$y$ mit rein
imagin\"{a}ren Periodicit\"{a}tsmoduln, und diese ist bis auf
eine additive Constante v\"{o}llig bestimmt.  Die
Function~$\omega$ stimmt dann mit $\alpha + \beta i$ in den
Unstetigkeiten im Innern von $T'$ und in den reellen Theilen der
Periodicit\"{a}tsmoduln \"{u}berein.  F\"{u}r $\omega$ k\"{o}nnen
daher die Functionen $\varphi_\nu$ und die reellen Theile ihrer
Periodicit\"{a}tsmoduln willk\"{u}rlich gegeben werden.  Durch
diese Bedingungen ist sie bis auf eine additive Constante
v\"{o}llig bestimmt, folglich auch der imagin\"{a}re Theil ihrer
Periodicit\"{a}tsmoduln.

Es wird sich zeigen, dass diese Function $\omega$ s\"{a}mmtliche
im \S.~1 bezeichneten Functionen als specielle F\"{a}lle unter
sich enth\"{a}lt.

\medbreak

\centerline{4.}

\nobreak\medskip

{\itshape
Allenthalben endliche Functionen $\omega$.  (Integrale erster
Gattung.)}

\nobreak\medskip

Wir wollen jetzt die einfachsten von ihnen betrachteten und zwar
zuerst diejenigen, die immer endlich bleiben und also im Innern
von $T'$ allenthalben stetig sind.  Sind $w_1, w_2,\ldots, w_p$
solche Functionen, so ist auch
\[ w = \alpha_1 w_1 + \alpha_2 w_2 + \cdots
         + \alpha_p w_p + \mbox{const.},\]
worin $\alpha_1, \alpha_2,\ldots, \alpha_p$ beliebige Constanten
sind, eine solche Function.  Es seien die Periodicit\"{a}tsmoduln
der Functionen $w_1, w_2,\ldots, w_p$ f\"{u}r den
$\nu \,$\-ten Querschnitt
$k_1^{(\nu)}, k_2^{(\nu)},\ldots, k_p^{(\nu)}$.  Der
Periodicit\"{a}tsmodul von $w$ f\"{u}r diesen Querschnitt ist
dann
$\alpha_1 k_1^{(\nu)} + \alpha_2 k_2^{(\nu)} + \cdots
      + \alpha_p k_p^{(\nu)} = k^{(\nu)}$;
und setzt man die Gr\"{o}ssen $\alpha$ in die Form
$\gamma + \delta i$, so sind die reellen Theile der $2p$
Gr\"{o}ssen $k^{(1)}, k^{(2)},\ldots, k^{(2p)}$ lineare
Functionen der Gr\"{o}ssen
$\gamma_1, \gamma_2,\ldots, \gamma_p$,
$\delta_1, \delta_2,\ldots, \delta_p$.
Wenn nun zwischen den Gr\"{o}ssen
$w_1, w_2,\ldots, w_p$
keine lineare Gleichung mit constanten Coefficienten stattfindet,
so kann die Determinante dieser linearen Ausdr\"{u}cke nicht
verschwinden; denn es liessen sich sonst die Verh\"{a}ltnisse der
Gr\"{o}ssen~$\alpha$ so bestimmen, dass die
Periodicit\"{a}tsmoduln des reellen Theils von $w$ s\"{a}mmtlich
$0$ w\"{u}rden, folglich der reelle Theil von $w$ und also auch
$w$ selbst nach dem \emph{Dirichlet}'schen Princip eine
Constante sein m\"{u}sste.  Es k\"{o}nnen daher dann die $2p$
Gr\"{o}ssen $\gamma$ und $\delta$ so bestimmt werden, dass die
reellen Theile der Periodicit\"{a}tsmoduln gegebene Werthe
erhalten; und folglich kann $w$ jede immer endlich bleibende
Function $\omega$ darstellen, wenn $w_1, w_2,\ldots, w_p$ keiner
linearen Gleichung mit constanten Coefficienten gen\"{u}gen.
Diese Functionen lassen sich aber immer dieser Bedingung
gem\"{a}ss w\"{a}hlen; denn so lange $\mu < p$, finden zwischen
den Periodicit\"{a}tsmoduln des reellen Theils von
\[ \alpha_1 w_1 + \alpha_2 w_2 + \cdots
      + \alpha_\mu w_\mu + \mbox{const.},\]
lineare Bedingungsgleichungen statt; es ist daher $w_{\mu + 1}$
nicht in dieser Form enthalten, wenn man, was nach dem Obigen
immer m\"{o}glich ist, die Periodicit\"{a}tsmoduln des reellen
Theils dieser Function so bestimmt, dass sie diesen
Bedingungsgleichungen nicht gen\"{u}gen.

\medbreak

{\itshape
Functionen $\omega$, die f\"{u}r einen Punkt der Fl\"{a}che~$T$
unendlich von der ersten Ordnung werden.  (Integrale zweiter
Gattung.)}

\nobreak\medskip

Es sei $\omega$ nur f\"{u}r einen Punkt~$\varepsilon$ der
Fl\"{a}che~$T$ unendlich, und f\"{u}r diesen seien alle
Coefficienten in $\varphi$ ausser $B$ gleich~$0$.  Eine solche
Function ist dann bis auf eine additive Constante bestimmt durch
die Gr\"{o}sse~$B$ und die reellen Theile ihrer
Periodicit\"{a}tsmoduln.  Bezeichnet $t^0(\varepsilon)$ irgend
eine solche Function, so k\"{o}nnen in dem Ausdrucke
\[ t(\varepsilon)
   = \beta t^0(\varepsilon)
      + \alpha_1 w_1 + \alpha_2 w_2 + \cdots
      + \alpha_p w_p + \mbox{const.} \]
die Constanten
$\beta, \alpha_1, \alpha_2,\ldots, \alpha_p$
immer so bestimmt werden, dass f\"{u}r ihn die Gr\"{o}sse~$B$ und
die reellen Theile der Periodicit\"{a}tsmoduln beliebig gegebene
Werthe erhalten.  Dieser Ausdruck stellt also jede solche
Function dar.

\medbreak

{\itshape
Functionen $\omega$, welche f\"{u}r zwei Punkte der
Fl\"{a}che~$T$ logarithmisch unendlich werden.  (Integrale dritte
Gattung.)}

\nobreak\medskip

Betrachten wir drittens den Fall, wo die Function~$\omega$ nur
logarithmisch unendlich wird, so muss dies, da die Summe der
Gr\"{o}ssen $A$ gleich $0$ sein muss, wenigstens f\"{u}r zwei
Punkte der Fl\"{a}che~$T$, $\varepsilon_1$ und $\varepsilon_2$,
geschehen und $A_2 = - A_1$ sein.  Ist von den Functionen, bei
denen dies statt hat und die beiden letztern Gr\"{o}ssen $= 1$
sind, irgend eine
$\varpi^0(\varepsilon_1, \varepsilon_2)$,
so sind nach \"{a}hnliche Schl\"{u}ssen, wie oben, alle
\"{u}brigen in der Form
\[ \varpi(\varepsilon_1, \varepsilon_2)
   = \varpi^0(\varepsilon_1, \varepsilon_2)
      + \alpha_1 w_1 + \alpha_2 w_2 + \cdots
      + \alpha_p w_p + \mbox{const.} \]
enthalten.

F\"{u}r die folgenden Bemerkungen nehmen wir zur Vereinfachung
an, dass die Punkte $\varepsilon$ keine Verzweigungspunkte sind
und nicht im Unendlichen liegen.  Man kann dann
$r_\nu = z - z_\nu$ setzen, indem man durch $z_\nu$ den Werth von
$z$ in $\varepsilon_\nu$ bezeichnet.  Wenn man dann
$\varpi(\varepsilon_1, \varepsilon_2)$
so nach $z_1$ differentiirt, dass die reellen Theile der
Periodicit\"{a}tsmoduln (oder auch $p$ von den
Periodicit\"{a}tsmoduln) und der Werth von
$\varpi(\varepsilon_1, \varepsilon_2)$
f\"{u}r einen beliebigen Punkt der Fl\"{a}che $T$ constant
bleiben, so erh\"{a}lt man eine Function $t(\varepsilon_1)$, die
in $\varepsilon_1$ unstetig wie
$\displaystyle \frac{1}{z - z_1}$
wird.  Umgekehrt ist, wenn $t(\varepsilon_1)$ eine solche
Function ist,
$\displaystyle \int\limits_{z_2}^{z_3} t(\varepsilon_1) \, \partial z_1$,
durch eine beliebige in $T$ von $\varepsilon_2$ nach
$\varepsilon_3$ f\"{u}hrende Linie genommen, gleich einer
Function
$\varpi(\varepsilon_2, \varepsilon_3)$.
Auf \"{a}hnliche Art erh\"{a}lt man durch $n$ successive
Differentiationen eines solchen $t(\varepsilon_1)$ nach $z_1$
Functionen $\omega$, welche im Punkte $\varepsilon_1$ wie
$n! {(z - z_1)^{-n-1}}$
unstetig werden und \"{u}brigens endlich bleiben.

F\"{u}r die ausgeschlossenen Lagen der Punkte $\varepsilon$
bed\"{u}rfen diese S\"{a}tze einer leichten Modification.

Offenbar kann nun ein mit constanten Coefficienten aus Functionen
$w$, aus Functionen $\varpi$ und ihren Derivirten nach
den Unstetigkeitswerthen gebildete linearer Ausdruck so bestimmt
werden, dass er im Innern von $T'$ beliebig gegebene
Unstetigkeiten von der Form, wie $\omega$, erh\"{a}lt, und die
reellen Theile seiner Periodicit\"{a}tsmoduln beliebig gegebene
Werthe annehmen.  Durch einen solchen Ausdruck kann also jede
gegebene Function $\omega$ dargestellt werden.

\medbreak

\centerline{5.}

\nobreak\medskip

Der allgemeine Ausdruck einer Function~$\omega$, die f\"{u}r $m$
Punkte der Fl\"{a}che $T$,
$\varepsilon_1, \varepsilon_2,\ldots, \varepsilon_m$
unendlich gross von der ersten Ordnung wird, ist nach dem Obigen
\[ s =  \beta_1 t_1  + \beta_2 t_2 + \cdots + \beta_m t_m
      + \alpha_1 w_1 + \alpha_2 w_2 + \cdots
      + \alpha_p w_p + \mbox{const.},\]
worin $t_\nu$ eine beliebige Function $t(\varepsilon_\nu)$ und
die Gr\"{o}ssen $\alpha$ und $\beta$ Constanten sind.  Wenn von
den $m$ Punkten $\varepsilon$ eine Anzahl~$\varrho$ in denselben
Punkt~$\eta$ der Fl\"{a}che~$T$ zusammenfallen, so sind die
$\varrho$ diesen Punkten zugeh\"{o}rigen Functionen $t$ zu
ersetzen durch eine Function $t(\eta)$ und deren $\varrho - 1$
erste Derivirte nach ihrem Unstetigkeitswerthe (\S.~2).

Die $2p$ Periodicit\"{a}tsmoduln dieser Function~$s$ sind lineare
homogene Functionen der $p + m$ Gr\"{o}ssen $\alpha$ und $\beta$.
Wenn $m \geq p + 1$, lassen sich also $2p$ von den Gr\"{o}ssen
$\alpha$ und $\beta$ als lineare homogene Functionen der
\"{u}brigen so bestimmen, dass die Periodicit\"{a}tsmoduln
s\"{a}mmtlich $0$ werden.  Die Function enth\"{a}lt dann noch
$m - p + 1$ willk\"{u}rliche Constanten, von denen sie eine
lineare homogene Function ist, und kann als ein linearer
Ausdruck von $m - p$ Functionen betrachtet werden, deren jede nur
f\"{u}r $p + 1$ Werthe unendlich von der ersten Ordnung wird.

Wenn $m = p + 1$ ist, so sind die Verh\"{a}ltnisse der $2p + 1$
Gr\"{o}ssen $\alpha$ und $\beta$ bei jeder Lage der $p + 1$
Punkte~$\varepsilon$ v\"{o}llig bestimmt.  Es k\"{o}nnen jedoch
f\"{u}r besondere Lagen dieser Punkte einige der Gr\"{o}ssen
$\beta$ gleich $0$ werden.  Die Anzahl dieser Gr\"{o}ssen sei
$= m - \mu$, so dass die Function nur f\"{u}r $\mu$ Punkte
unendlich von der ersten Ordnung wird.  Diese $\mu$ Punkte
m\"{u}ssen dann eine solche Lage haben, dass von den $2p$
Bedingungsgleichungen zwischen den $p + \mu$ \"{u}brigen
Gr\"{o}ssen $\beta$ und $\alpha$ $p + 1 - \mu$ eine identische
Folge der \"{u}brigen sind, und es k\"{o}nnen daher  nur
$2\mu - p - 1$ von ihnen beliebig gew\"{a}hlt werden.  Ausserdem
enth\"{a}lt die Function noch 2 willk\"{u}rliche Constanten.

Es sei nun $s$ so zu bestimmen, dass $\mu$ m\"{o}glichst klein
wird.  Wenn $s$ $\mu \,$\-mal unendlich von der ersten Ordnung wird,
so ist dies auch mit jeder rationalen Function ersten Grades von
$s$ der Fall; man kann daher f\"{u}r die L\"{o}sung dieser
Aufgabe einen der $\mu$ Punkte beliebig w\"{a}hlen.  Die Lage der
\"{u}brigen muss dann so bestimmt werden, dass $p + 1 - \mu$ von
den Bedingungsgleichungen zwischen den Gr\"{o}ssen $\alpha$ und
$\beta$ eine identische Folge der \"{u}brigen sind; es muss also,
wenn die Verzweigungswerthe der Fl\"{a}che~$T$ nicht besondern
Bedingungsgleichungen gen\"{u}gen, $p + 1 - \mu \leq \mu - 1$ oder
$\mu \geq \frac{1}{2}p + 1$ sein.

Die Anzahl der in einer Function~$s$, die nur f\"{u}r $m$ Punkte
der Fl\"{a}che~$T$ unendlich von der ersten Ordnung wird und
\"{u}brigens stetig bleibt, enthaltenen willk\"{u}rlichen
Constanten ist in allen F\"{a}llen $= 2m - p + 1$.

{\itshape
\label{`Wurzel'}
Eine solche Function ist die Wurzel einer Gleichung
$n^{\mathrm{ten}}$ Grades, deren Coefficienten ganze Functionen
$m^{\mathrm{ten}}$ Grades von $z$ sind.}

Sind $s_1, s_2,\ldots, s_n$ die $n$ Werthe der Function~$s$
f\"{u}r dasselbe $z$, und bezeichnet $\sigma$ eine beliebige
Gr\"{o}sse, so ist
$(\sigma - s_1) (\sigma - s_2) \cdots (\sigma - s_n)$
eine einwerthige Function von $z$, die nur f\"{u}r einen Punkt
der $z$-Ebene, der mit einem Punkte~$\varepsilon$
zusammenf\"{a}llt, unendlich wird und unendlich von einer so
hohen Ordnung, als Punkte~$\varepsilon$ auf ihn fallen.  In der
That wird f\"{u}r jeden auf ihn fallenden Punkt~$\varepsilon$,
der kein Verzweigungspunkt ist, nur \emph{ein} Factor dieses
Products von einer um $1$ h\"{o}heren Ordnung unendlich, f\"{u}r
einen Punkt~$\varepsilon$, um den die Fl\"{a}che $T$ sich
$\mu \,$\-mal windet, aber $\mu$ Factoren von einer um
$\displaystyle\frac{1}{\mu}$ h\"{o}heren Ordnung.  Bezeichnet man
nun die Werthe von $z$ in \emph{den} Punkten~$\varepsilon$, wo
$z$ nicht unendlich ist, durch
$\zeta_1, \zeta_2,\ldots, \zeta_\nu$
und
$(z - \zeta_1) (z - \zeta_2) \cdots (z - \zeta_\nu)$
durch $a_0$, so ist
$a_0 (\sigma - s_1) (\sigma - s_2) \cdots (\sigma - s_n)$
eine einwerthige Function von $z$, die f\"{u}r alle endlichen
Werthe von $z$ endlich ist, und f\"{u}r $z = \infty$ unendlich
von der $m^{\mathrm{ten}}$ Ordnung wird, also eine ganze Function
$m^{\mathrm{ten}}$ Grades von $z$.  Sie ist zugleich eine ganze
Function $n^{\mathrm{ten}}$ Grades von $\sigma$, dir f\"{u}r
$\sigma = s$ verschwindet.  Bezeichnet man sie durch $F$ und, wie
wir in der Folge thun wollen, \emph{eine ganze Function $F$
$n^{\mathrm{ten}}$ Grades von $\sigma$ und $m^{\mathrm{ten}}$
Grades von $z$} durch
$F( {\buildrel n \over \sigma}, {\buildrel m \over z} )$,
so ist $s$ die Wurzel der Gleichung
$F( {\buildrel n \over s}, {\buildrel m \over z} ) = 0$.

Die Function~$F$ ist eine Potenz einer
unzerf\"{a}llbaren---d.~h.\ nicht als ein Product aus ganzen
Functionen von $\sigma$ und $z$ darstellbaren---Function.  Den
jeder ganze rationale Factor von $F(\sigma, z)$ bildet, da er
f\"{u}r einige der Wurzeln $s_1, s_2,\ldots, s_n$ verschwinden
muss, f\"{u}r $\sigma = s$ eine Function von $z$, die in einem
Theile der Fl\"{a}che $T$ verschwindet und folglich, da diese
Fl\"{a}che zusammenh\"{a}ngend ist, in der ganzen Fl\"{a}che $0$
sein muss.  Zwei unzerf\"{a}llbare Factoren von $F(\sigma, z)$
k\"{o}nnten aber nur f\"{u}r eine endliche Anzahl von
Werthenpaaren zugleich verschwinden, wenn die eine nicht durch
Multiplication mit einer Constanten aus der andern erhalten
werden k\"{o}nnte.  Folglich muss $F$ eine Potenz einer
unzerf\"{a}llbaren Function sein.

Wenn der Exponent~$\nu$ dieser Potenz $> 1$ ist, so wird die
Verzweigungsart der Function $s$ nicht dargestellt durch die
Fl\"{a}che $T$, sondern durch eine in der $z$-Ebene allenthalben
$\displaystyle\frac{n}{\nu} \,$\-fach
ausgebreitete Fl\"{a}che~$\tau$, in welche die Fl\"{a}che~$T$
allenthalben $\nu \,$\-fach ausgebreitet ist.  Es kann dann zwar $s$
als eine wie~$T$ verzweigte Function betrachtet werden, nicht aber
umgekehrt~$T$ als verzweigt, wie~s.

Eine solche nur in einzelnen Punkten von $T$ unstetige Function,
wie $s$, ist auch
$\displaystyle\frac{\partial \omega}{\partial z}$.
Denn diese Function nimmt zu beiden Seiten der Querschnitte und
der Linien~$l$ denselben Werth an, da die Differenz der beiden
Werthe von $\omega$ in diesen Linien l\"{a}ngs denselben constant
ist; sie kann nur unendlich werden, wo $\omega$ unendlich wird,
und in den Verzweigungspunkten der Fl\"{a}che und ist sonst
allenthalben stetig, da die Derivirte einer ein\"{a}ndrig und
endlich bleibenden Function ebenfalls ein\"{a}ndrig und endlich
bleibt.

Es sind daher s\"{a}mmtliche Functionen $\omega$ algebraische wie
$T$ verzweigte Functionen von $z$ oder Integrale solcher
Functionen.  Dieses System von Functionen ist bestimmt, wenn die
Fl\"{a}che~$T$ gegeben ist und h\"{a}ngt nur von der Lage ihrer
Verzweigungspunkte ab.

\medbreak

\centerline{6.}

\nobreak\medskip

Es sei jetzt die irreductible Gleichung
$F( {\buildrel n \over s}, {\buildrel m \over z} ) = 0$
gegeben und die Art der Verzweigung der Function~$s$ oder der die
darstellenden Fl\"{a}che~$T$ zu bestimmen.  Wenn f\"{u}r einen
Werth~$\beta$ von $z$ $\mu$ Zweige der Function
zusammenh\"{a}ngen, so dass einer dieser Zweige sich erst nach
$\mu$~Umlaufen des $z$ um $\beta$ wieder in sich selbst
fortsetzt, so k\"{o}nnen diese $\mu$ Zweige der Function (wie
nach \emph{Cauchy} oder durch die \emph{Fourier}'sche Reihe
leicht bewiesen werden kann) dargestellt werden durch eine Reihe
nach steigenden rationalen Potenzen von $z - \beta$ mit
Exponenten vom kleinsten gemeinschaftlichen Nenner~$\mu$, und
umgekehrt.

Ein Punkt der Fl\"{a}che~$T$, in welchem nur zwei Zweige einer
Function zusammenh\"{a}ngen, so dass sich um diesen Punkt der
erste in den zweiten und dieser in jenen fortsetzt, heisse ein
\emph{einfacher Verzweigungspunkt}.

Ein Punkt der Fl\"{a}che, um welchen sie sich ${(\mu + 1)} \,$\-mal
windet, kann dann angesehen werden als $\mu$ zusammengefallene
(oder unendlich nahe) einfache Verzweigungspunkte.

Um dies zu zeigen, seien in einem diesen Punkt umgebenden
St\"{u}cke der $z$-Ebene $s_1, s_2,\ldots, s_{\mu + 1}$
ein\"{a}ndrige Zweige der Function~$s$ und in der Begrenz\-ung
desselben, bei positiver Umschreibung auf einander folgend,
$a_1, a_2,\ldots, a_{\mu}$ einfache Verzweigungspunkte.  Durch
einen positiven Umlauf um $a_1$ werde $s_1$ mit $s_2$, um $a_2$
$s_1$ mit $s_3$,\ldots, um $a_\mu$ $s_1$ mit $s_{\mu + 1}$
vertauscht.  Es gehen dann nach nach einem positiven Umlaufe um
ein alle diese Punkte (und keinen andern Verzweigungspunkt)
enthaltenes Gebiet
\[ s_1, s_2,\ldots, s_\mu, s_{\mu + 1} \]
\[ \mbox{in } s_2, s_3,\ldots, s_{\mu + 1}, s_1, \mbox{ \"{u}ber,} \]
und es entsteht daher, wenn sie zusammenfallen, ein $\mu \,$\-facher
Windungspunkt.

Die Eigenschaften der Functionen $\omega$ h\"{a}ngen wesentlich
davon ab, wie vielfach zusammenh\"{a}ngend die Fl\"{a}che~$T$
ist.  Um dies zu entscheiden, wollen wir zun\"{a}chst die Anzahl
der einfachen Verzweigungspunkte der Function~$s$ bestimmen.

In einem Verzweigungspunkte nehmen die dort zusammenh\"{a}ngenden
Zw\-eige der Function denselben Werth an, und es werden daher zwei
oder mehrere Wurzeln der Gleichung
\[ F(s) = a_0 s^n + a_1 s^{n-1} + \cdots + a_n = 0 \]
einander gleich.  Dies kann nur geschehen, wenn
\[ F'(s) = a_0 n s^{n-1} + a_1 \overline{n-1} s^{n-2} + \cdots + a_{n-1} \]
oder die einwerthige Function von $z$,
$F'(s_1) \, F'(s_2) \, \ldots \, F'(s_n)$,
verschwindet.  Diese Function wird f\"{u}r endliche Werthe von
$z$ nur unendlich, wenn $s = \infty$, also $a_0 = 0$ ist und
muss, um endlich zu bleiben, mit $a_0^{n-2}$ multiplicirt werden.
Sie wird dann eine einwerthige, f\"{u}r eine endliches $z$
endliche Function von $z$, welche f\"{u}r $z = \infty$ unendlich
von der ${2m(n - 1)} \,$\-ten Ordnung wird, also eine ganze Function
${2m(n - 1)} \,$\-ten Grades.  Die Werthe von $z$, f\"{u}r welche
$F'(s)$ und $F(s)$ gleichzeitig verschwinden, sind also die
Wurzeln der Gleichung ${2m(n - 1)} \,$\-ten Grades
\begin{eqnarray*}
Q(z) &=& a_0^{n-2} \prod_i F'(s_i) = 0 \\
     & & \mbox{ oder auch, da }
         F'(s_i) = a_0 \prod_{i'} (s_i - s_{i'}),
         \enspace (i \neq i'), \\
     &=& a_0^{2(n-1)} \prod_{i,i'} (s_i - s_{i'}) = 0,
         \enspace (i \neq i'),
\end{eqnarray*}
welche durch Elimination von $s$ aus $F'(s) = 0$ und $F(s) = 0$
gebildet werden kann.

Wird $F(s,z) = 0$ f\"{u}r $s = \alpha$, $z = \beta$, so ist
\begin{eqnarray*}
F(s,z) &=& \frac{\partial F}{\partial s}(s - \alpha)
           + \frac{\partial F}{\partial z}(z - \beta) \\
       & & + {\textstyle\frac{1}{2}}
               \left\{
                  \frac{\partial^2 F}{\partial s^2} (s - \alpha)^2
                + 2 \frac{\partial^2 F}{\partial s \, \partial z}
                     (s - \alpha) (z - \beta)
                + \frac{\partial^2 F}{\partial z^2} (z - \beta)^2
               \right\} \\
       & & + \cdots ,\\
F'(s)  &=& \frac{\partial F}{\partial s}
           + \frac{\partial^2 F}{\partial s^2}(s - \alpha)
           + \frac{\partial^2 F}{\partial s \, \partial z} (z - \beta)
           + \cdots
\end{eqnarray*}
Ist also f\"{u}r ($s = \alpha$, $z = \beta$)
$\displaystyle \frac{\partial F}{\partial s} = 0$
und verschwinden
$\displaystyle \frac{\partial F}{\partial z}$,
$\displaystyle \frac{\partial^2 F}{\partial s^2}$
dann nicht, so wird $s - \alpha$ unendlich klein, wie
$(z - \beta)^{\frac{1}{2}}$, und findet also ein einfacher
Verzweigungspunkt statt.  Es werden zugleich in dem Producte
$\prod\limits_i F'(s_i)$ zwei Factoren unendlich klein wie
$(z - \beta)^{\frac{1}{2}}$, und $Q(z)$ erh\"{a}lt dadurch den
Factor $(z - \beta)$.  In dem Falle, dass
$\displaystyle \frac{\partial F}{\partial z}$
und
$\displaystyle \frac{\partial^2 F}{\partial s^2}$
nie verschwinden, wenn gleichzeitig $F = 0$ und
$\displaystyle \frac{\partial F}{\partial s} = 0$
werden, entspricht demnach jedem linearen Factor von $Q(z)$
\emph{ein} einfacher Verzweigungspunkt, und die Anzahl dieser
Punkte ist also $= 2m(n - 1)$.

Die Lage der Verzweigungspunkte h\"{a}ngt von den Coefficienten
der Potenzen von $z$ in den Functionen $a$ ab und \"{a}ndert sich
stetig mit denselben.

Wenn diese Coefficienten solche Werthe annehmen, dass zwei
demselben Zweigepaar angeh\"{o}rige einfache Verzweigungspunkte
zusammenfallen, so heben diese sich auf, und es werden zwei
Wurzeln von $F(s)$ einander gleich, ohne dass eine Verzweigung
stattfindet.  Setzt sich um jeden von ihnen $s_1$ in $s_2$ und
$s_2$ in $s_1$ fort, so geht durch einen Umlauf um ein beide
enthaltendes St\"{u}ck der $z$-Ebene $s_1$ in $s_1$ und $s_2$ in
$s_2$ \"{u}ber, und beide Zweige werden ein\"{a}ndrig, wenn sie
zusammenfallen.  Es bleibt dann also auch ihre Derivirte
$\displaystyle\frac{\partial s}{\partial z}$
ein\"{a}ndrig und endlich, und folglich wird
$\displaystyle\frac{\partial F}{\partial z}
   = - \frac{\partial s}{\partial z} \frac{\partial F}{\partial s} = 0$.

Wird
$\displaystyle F
   = \frac{\partial F}{\partial s}
   = \frac{\partial F}{\partial z}
   = 0$
f\"{u}r $s = \alpha$, $z = \beta$, so ergeben sich aus den drei
folgenden Gliedern der Entwicklung von $F(s,z)$ zwei Werthe
f\"{u}r
$\displaystyle\frac{s - \alpha}{s- \beta} = \frac{\partial s}{\partial z}$,
($s = \alpha$, $z = \beta$).  Sind diese Werthe ungleich und
endlich, so k\"{o}nnen die beiden Zweige der Function~$s$, denen
sie angeh\"{o}ren, dort nicht zusammenh\"{a}ngen und sich nicht
verzweigen.  Es wird dann
$\displaystyle \frac{\partial F}{\partial s}$
f\"{u}r beide unendlich klein wie $z - \beta$, und $Q(z)$
erh\"{a}lt dadurch den Factor $(z - \beta)^2$; es fallen also nur
zwei einfache Verzweigungspunkte zusammen.

Um in jeden Falle, wenn f\"{u}r $z = \beta$ mehrere Wurzeln der
Gleichung $F(s) = 0$ gleich $\alpha$ werden, zu entscheiden, wie
viele einfache Verzweigungspunkte f\"{u}r ($s = \alpha$, $z =
\beta$) zusammenfallen, und wie viele von diesen sich aufheben,
muss man diese Wurzeln (nach dem Verfahren von \emph{Lagrange})
soweit nach steigenden Potenzen von $z - \beta$ entwickeln, bis
diese Entwicklungen s\"{a}mmtlich von einander verschieden werden,
wodurch sich die wirklich noch stattfindenen Verzweigungen
ergeben.  Und man muss dann untersuchen, von welcher Ordnung
$F'(s)$ f\"{u}r jede dieser Wurzeln unendlich klein wird, um die
Anzahl der ihnen zugeh\"{o}rigen linearen Factoren von $Q(z)$
oder der f\"{u}r ($s = \alpha$, $z = \beta$) zusammengefallenen
einfachen Verzweigungspunkte zu bestimmen.

Bezeichnet die Zahl~$\varrho$, wie oft sich die Fl\"{a}che~$T$
um den Punkt $(s,z)$ windet, so wird im Punkte ($z$) $F'(s)$ so
oft unendlich klein von der ersten Ordnung, als dort einfache
Verzweigungspunkte zusammenfallen,
$dz^{1 - \frac{1}{\varrho}}$ so oft, als deren wirklich
stattfinden, folglich $F'(s) \, dz^{\frac{1}{\varrho} - 1}$ so
oft, als von ihnen sich aufheben.

Ist die Anzahl der wirklich stattfindenden einfachen
Verzweigungen~$\mathrm{w}$, die Anzahl der sich aufhebenden $2r$,
so ist
\[ \mathrm{w} + 2r = 2 (n - 1) m.\]
Nimmt man an, dass die Verzweigungspunkte nur paarweise und sich
aufhebend zusammenfallen, so ist f\"{u}r $r$ Werthenpaare
\label{`Werthenpaare'}
($s= \gamma_\varrho$, $z = \delta_\varrho$)
\[ F = \frac{\partial F}{\partial s}
     = \frac{\partial F}{\partial z} = 0
      \quad\mbox{und}\quad
   \frac{\partial^2 F}{\partial z^2} \frac{\partial^2 F}{\partial s^2}
   - \left( \frac{\partial^2 F}{\partial s \, \partial z} \right)^2 \]
nicht Null und f\"{u}r $\mathrm{w}$ Werthenpaare von $s$ und $z$
$F = 0$,
$\displaystyle \frac{\partial F}{\partial s} = 0$,
$\displaystyle \frac{\partial F}{\partial z}$
nicht Null und
$\displaystyle \frac{\partial^2 F}{\partial s^2}$
nicht Null.

Wir beschr\"{a}nken uns meistens auf die Behandlung dieses
Falles, da sich die Resultate auf die \"{u}brigen als
Grenzf\"{a}lle desselben leicht ausdehnen lassen, und wir
k\"{o}nnen dies hier um so mehr thun, da wir die Theorie dieser
Functionen auf eine von der Ausdrucksform unabh\"{a}ngige, keinen
Ausnahmef\"{a}llen unterworfene Grundlage gest\"{u}tzt haben.

\medbreak

\centerline{7.}

\nobreak\medskip

Es findet nun bei einer einfach zusammenh\"{a}ngenden, \"{u}ber
einen endlichen Theil der $z$-Ebene ausgebreiteten Fl\"{a}che
zwischen der Anzahl ihrer einfachen Verzweigungspunkte und der
Anzahl der Umdrehungen, welche die Richtung ihrer Begrenzlinie
macht, die Relation statt, dass die letztere um eine Einheit
gr\"{o}sser ist, als die erstere; und aus dieser ergiebt sich
f\"{u}r eine mehrfach zusammenh\"{a}ngende Fl\"{a}che eine
Relation zwischen diesen Anzahlen und der Anzahl der
Querschnitte, welche sie in eine einfach zusammenh\"{a}ngende
verwandeln.  Wir k\"{o}nnen diese Relation, welche im Grunde von
Massverh\"{a}ltnissen unabh\"{a}ngig ist und der \emph{analysis
situs} angeh\"{o}rt, hier f\"{u}r die Fl\"{a}che~$T$ so ableiten.

Nach dem \emph{Dirichlet}'schen Princip l\"{a}sst sich in der
einfach zusammenh\"{a}ngenden Fl\"{a}che $T'$ die Function
$\log \zeta$ von $z$ so bestimmen, dass $\zeta$ f\"{u}r einen
beliebigen Punkt im Innern derselben unendlich klein von der
ersten Ordnung wird, und $\log \zeta$ l\"{a}ngs einer beliebigen
sich nicht schneidenden, von dort nach der Begrenzung
f\"{u}hrenden Linie auf der positiven Seite um $-2\pi i$
gr\"{o}sser, als auf der negativen, \"{u}brigens aber
allenthalben stetig und l\"{a}ngs der Begrenzung von $T'$ rein
imagin\"{a}r ist.  Es nimmt dann die Function $\zeta$ jeden
Werth, dessen Modul $< 1$, einmal an; die Gesamtheit ihrer Werthe
wird folglich durch eine \"{u}ber einen Kreis in der
$\zeta$-Ebene einfach ausgebreitete Fl\"{a}che vertreten.  Jedem
Punkte von $T'$ entspricht ein Punkt des Kreises, und umgekehrt.
Es wird daher f\"{u}r einen beliebigen Punkt der Fl\"{a}che, wo
$z = z'$, $\zeta = \zeta'$, die Function $\zeta - \zeta'$
unendlich klein von der ersten Ordnung, und folglich bleibt dort,
wenn die Fl\"{a}che~$T'$ sich ${(\mu + 1)} \,$\-mal um ihn windet,
bei endlichem~$z'$
\[ (\mu + 1) \frac{z - z'}{(\zeta - \zeta')^{\mu + 1}}
   = \frac{dz}{d\zeta (\zeta - \zeta')^\mu},\]
bei unendlichem $z'$ aber
\[ (\mu + 1) \frac{z^{-1}}{(\zeta - \zeta')^{\mu + 1}}
   = - \frac{\partial z}{z z \partial \zeta (\zeta - \zeta')^\mu} \]
endlich.  Das Integral
$\displaystyle \int \partial \log \frac{\partial z}{\partial \zeta}$,
um den ganzen Kreis positiv herumgenommen, ist gleich der Summe
der Integrale um die Punkte, wo
$\displaystyle \frac{\partial z}{\partial \zeta}$
unendlich oder Null wird, und also $= 2\pi i (\mathrm{w} - 2n)$.
Bezeichnet $s$ ein St\"{u}ck der Begrenzung von $T'$ von einem
und demselben bestimmten Punkte bis zu einem ver\"{a}nderlichen
Punkte der Begrenzung, und $\sigma$ das entsprechende St\"{u}ck
auf dem Kreisumfange, so ist
\[ \log \frac{\partial z}{\partial \zeta}
   =  \log \frac{\partial z}{\partial s}
      + \log \frac{\partial s}{\partial \sigma}
      - \log \frac{\partial \zeta}{\partial \sigma},\]
und, durch die ganze Begrenzung ausgedehnt,
\begin{eqnarray*}
\int \partial \log \frac{\partial z}{\partial s}
   &=& (2p - 1) 2 \pi i,\\
\int \partial \log \frac{\partial s}{\partial \sigma}
   &=& 0,\\
- \int \partial \log \frac{\partial \zeta}{\partial \sigma}
   &=& - 2 \pi i,
\end{eqnarray*}
also
\[ \int \partial \log \frac{\partial z}{\partial \zeta} = (2p - 2) 2 \pi i.\]
Es ergiebt sich demnach $\mathrm{w} - 2n = 2 (p - 1)$.  Da nun
\[ \mathrm{w} = 2 ( (n - 1) m - r ),\]
so ist
\[ p = (n - 1) (m - 1) - r.\]

\medbreak

\centerline{8.}

\nobreak\medskip

Der allgemeine Ausdruck der wie $T$ verzweigten Functionen $s'$
von $z$, die f\"{u}r $m'$ beliebig gegebene Punkte von $T$
endlich von der ersten Ordnung werden und \"{u}brigens stetig
bleiben, enh\"{a}lt nach dem Obigen $m' - p + 1$ willk\"{u}rliche
Constanten und ist eine lineare Function derselben (\S.~5).
Lassen sich also, wie jetzt gezeigt werden soll, rationale
Ausdr\"{u}cke von $s$ und $z$ bilden, die f\"{u}r $m'$ beliebig
gegebene, der Gleichung $F = 0$ gen\"{u}gende Werthenpaare von
$s$ und $z$ unendlich von der ersten Ordnung werden und lineare
Functionen von $m' - p + 1$ willk\"{u}rlichen Constanten sind, so
kann durch diese Ausdr\"{u}cke jede Function $s'$ dargestellt
werden.

Damit der Quotient zweier ganzen Functionen $\chi(s,z)$ und
$\psi(s,z)$ f\"{u}r $s = \infty$ und $z = \infty$ beliebige
endliche Werthe annehmen kann, m\"{u}ssen beide von gleichem
Grade sein; der Ausdruck, durch welchen $s'$ dargestellt werden
soll, sei daher von der Form
$\displaystyle\frac{
   \psi( {\buildrel \nu \over s}, {\buildrel \mu \over z} )}{
   \chi( {\buildrel \nu \over s}, {\buildrel \mu \over z} )}$,
und \"{u}berdies sei $\nu \geq n - 1$, $\mu \geq m - 1$.  Wenn
zwei Zweige der Function~$s$ ohne zusammenzuh\"{a}ngen einander
gleich werden, also f\"{u}r zwei verschiedene Punkte der
Fl\"{a}che $T$ $z = \gamma$ und $s = \delta$ wird, so wird $s'$
allgemein zu reden in diesen beiden Punkten verschiedene Werthe
annehmen; soll also $\psi - s' \chi$ allenthalben $= 0$ sein, so
muss f\"{u}r zwei verschiedene Werthe von $s'$
$\psi(\gamma, \delta) - s' \chi(\gamma, \delta) = 0$
sein, folglich $\chi(\gamma, \delta) = 0$ und
$\psi(\gamma, \delta) = 0$.  Es m\"{u}ssen also die Functionen
$\chi$ und $\psi$ f\"{u}r die $r$ Werthenpaare
$s = \gamma_\varrho$, $z = \delta_\varrho$
(S.~\pageref{`Werthenpaare'}) verschwinden\footnote{Es ist hier,
wie gesagt, nur der Fall ber\"{u}cksichtigt, wo die
Verzweigungspunkte der Function~$s$ nur paarweise und sich
aufhebend zusammenfallen.  Im Allgemeinen m\"{u}ssen in einem
Punkte von $T$, wo nach der Auffassung im \S.~6 sich aufhebende
Verzweigungspunkte zusammenfallen, $\chi$ und $\psi$, wenn $T$
sich um diesen Punkt~$\varrho$ mal windet, unendlich klein
werden, wie $F'(s) \, dz^{\frac{1}{\varrho} - 1}$, damit die
ersten Glieder in der Entwicklung der darzustellenden
Function nach ganzen Potenzen von
$(\Delta z)^{\frac{1}{\varrho}}$ beliebige Werthe annehmen
k\"{o}nnen}.

Die Function $\chi$ verschwindet f\"{u}r einen Werth von $z$,
f\"{u}r welchen die einwerthige und f\"{u}r ein endliches~$z$
endliche Function von $z$
\[ K(z) = a_0^\nu \chi(s_1) \chi(s_2) \ldots \chi(s_n) = 0 \]
ist; diese Function wird f\"{u}r ein unendliches $z$ unendlich
von der Ordnung $m \nu + n \mu$ und ist also eine ganze Function
${(m \nu + n \mu)} \,$\-ten Grades.  Da f\"{u}r die Werthenpaare
$(\gamma, \delta)$ zwei Factoren des Products
$\prod\limits_i \chi(s_i)$ unendlich klein von der ersten Ordnung
werden, also $K(z)$ unendlich klein von der zweiten Ordnung, so
wird $\chi$ ausserdem noch unendlich klein von der ersten Ordnung
f\"{u}r
\[ i = m \nu + n \mu - 2r \]
Werthenpaare von $s$ und $z$ oder Punkte von $T$.

Ist $\nu > n - 1$, $\mu > m - 1$, so bleibt der Werth der
Function~$\chi$ unge\"{a}ndert, wenn man
\[ \chi( {\buildrel \nu \over s}, {\buildrel \mu \over z} )
      + \varrho( {\buildrel \nu - n \over s}, {\buildrel \mu - m \over z} )
        F( {\buildrel n \over s}, {\buildrel m \over z} ),\]
worin $\varrho$ beliebig ist, f\"{u}r
$\chi( {\buildrel \nu \over s}, {\buildrel \mu \over z} )$
setzt; es k\"{o}nnen also
\[ (\nu - n + 1) (\mu - m + 1) \]
von den Coefficienten dieses Ausdrucks willk\"{u}rlich angenommen
werden.  Werden nun von den
\[ (\mu + 1)(\nu + 1) - (\nu - n + 1)(\mu - m + 1) \]
noch \"{u}brigen $r$ als lineare Functionen der \"{u}brigen so
bestimmt, dass $\chi$ f\"{u}r die $r$ Werthenpaare
$(\gamma, \delta)$ verschwindet, so enth\"{a}lt die Function
$\chi$ noch
\begin{eqnarray*}
\varepsilon
   &=& (\mu + 1)(\nu + 1) - (\nu - n + 1)(\mu - m + 1) - r \\
   &=& n \mu + m \nu - (n - 1)(m - 1) - r + 1
\end{eqnarray*}
willk\"{u}rliche Constanten.  Es ist also
\[ i - \varepsilon = (n - 1)(m - 1) - r - 1 = p - 1.\]

Nimmt man $\mu$ und $\nu$ so an, dass $\varepsilon > m'$ ist, so
kann man $\chi$ so bestimmen, dass es f\"{u}r $m'$ beliebig
gegebene Werthenpaare unendlich klein von der ersten Ordnung
wird, und dann, wenn $m' > p$, $\psi$ so einrichten, dass
$\displaystyle \frac{\psi}{\chi}$
f\"{u}r alle \"{u}brigen Werthe endlich bleibt.  In der That ist
$\psi$ ebenfalls eine lineare homogene Function von $\varepsilon$
willk\"{u}rlichen Constanten, und es lassen sich also, wenn
$\varepsilon - i + m' > 1$ ist, $i - m'$ von ihnen als lineare
Functionen der \"{u}brigen so bestimmen, dass $\psi$ f\"{u}r die
$i - m'$ Werthenpaare von $s$ und $z$, f\"{u}r welche $\chi$ noch
unendlich klein von der ersten Ordnung wird, ebenfalls
verschwindet.  Die Function $\psi$ enth\"{a}lt demnach
$\varepsilon - i + m' = m' - p + 1$
willk\"{u}rliche Constanten, und
$\displaystyle \frac{\psi}{\chi}$
kann also jede Function $s'$ darstellen.

\medbreak

\centerline{9.}

\nobreak\medskip

Da die Functionen
$\displaystyle \frac{\partial \omega}{\partial z}$
algebraische wie $s$ verzweigte Functionen von $z$ sind (\S.~5),
so lassen sie sich zufolge des eben bewiesenen Satzes rational in
$s$ und $z$ ausdr\"{u}cken, und s\"{a}mmtliche Functionen
$\omega$ als Integrale rationaler Functionen von $s$ und $z$.

Ist $w$ eine allenthalben endliche Function~$\omega$, so wird
$\displaystyle \frac{\partial w}{\partial z}$
unendlich von der ersten Ordnung f\"{u}r jeden einfachen
Verzweigungspunkt der Fl\"{a}che~$T$, da $dw$ und
$(dz)^{\frac{1}{2}}$ dort unendlich klein von der ersten Ordnung
sind, bleibt aber sonst allenthalben stetig und wird f\"{u}r
$z = \infty$ unendlich klein von der zweiten Ordnung.  Umgekehrt
bleibt das Integral einer Function, die sich so verh\"{a}lt,
allenthalben endlich.

Um diese Function
$\displaystyle \frac{\partial w}{\partial z}$
als Quotient zweier ganzen Functionen von $s$ und $z$
auszudr\"{u}cken, muss man (nach \S.~8) zum Nenner eine Function
nehmen, die verschwindet in den Verzweigungspunkten und f\"{u}r
die $r$ Werthenpaare $(\gamma, \delta)$.  Dieser Bedingung
gen\"{u}gt man am einfachsten durch eine Function, die nur
f\"{u}r diese Werthe $0$ wird.  Eine solche ist
\[ \frac{\partial F}{\partial s}
   = a_0 n s^{n-1} + a_1 \overline{n - 1} s^{n-2} + \cdots + a_{n-1}.\]
Diese wird f\"{u}r ein unendliches $s$ unendlich von der
${(n - 2)} \,$\-ten Ordnung (da $a_0$ dann unendlich klein von der
ersten Ordung wird) und f\"{u}r ein unendliches $z$ unendlich von
der $m \,$\-ten Ordnung.  Damit
$\displaystyle \frac{\partial w}{\partial z}$
ausser den Verzweigungspunkten endlich und f\"{u}r ein unendliches
$z$ unendlich klein von der zweiten Ordnung ist, muss also der
Z\"{a}hler eine ganze Function
$\varphi( {\buildrel n - 2 \over s}, {\buildrel m - 2 \over z} )$
sein, die f\"{u}r die $r$ Werthenpaare $(\gamma, \delta)$
(S.~\pageref{`Werthenpaare'}) verschwindet.  Demnach ist
\[ w = \int
   \frac{\varphi( {\buildrel n - 2 \over s}, {\buildrel m - 2 \over z} )
   \, \partial z}{\displaystyle\frac{\partial F}{\partial s}}
   = - \int
   \frac{\varphi( {\buildrel n - 2 \over s}, {\buildrel m - 2 \over z} )
   \, \partial s}{\displaystyle\frac{\partial F}{\partial z}},\]
worin $\varphi = 0$ f\"{u}r
$s = \gamma_\varrho$, $z = \delta_\varrho$,
$\varrho = 1,2,\ldots, r$.

Die Function $\varphi$ enh\"{a}lt $(n - 1)(m - 1)$ constante
Coefficienten, und wenn $r$ von ihnen als lineare Functionen der
\"{u}brigen so bestimmt werden, dass $\varphi = 0$ f\"{u}r die
$r$ Werthenpaare $s = \gamma$, $z = \delta$, so bleiben noch
$(m - 1)(n - 1) - r$ oder $p$ willk\"{u}rlich, und es erh\"{a}lt
$\varphi$ die Form
\[ \alpha_1 \varphi_1 + \alpha_2 \varphi_2 + \cdots
      + \alpha_p \varphi_p,\]
worin $\varphi_1, \varphi_2,\ldots, \varphi_p$ besondere
Functionen $\varphi$, von denen keine eine lineare Function der
\"{u}brigen ist, und $\alpha_1, \alpha_2,\ldots, \alpha_p$
beliebige Constanten sind.  Als allgemeiner Ausdruck von $w$
ergiebt sich, wie oben auf anderem Wege
\[ \alpha_1 w_1 + \alpha_2 w_2 + \cdots + \alpha_p w_p
      + \mbox{const.} \]

Die nicht allenthalben endlich bleibenden Function~$\omega$ und
also die Integrale zweiter und dritter Gattung lassen sich nach
denselben Principien rational in $s$ und $z$ ausdr\"{u}cken,
wobei wir indess hier nicht verweilen, da die allgemeinen Regeln
des vorigen Paragraphen keiner weitern Erl\"{a}uterung
bed\"{u}rfen und zur Betrachtung bestimmter Formen dieser
Integrale erst die Theorie der $\vartheta$-Functionen Anlass
giebt.

\medbreak

\centerline{10.}

\nobreak\medskip

Die Function $\varphi$ wird ausser f\"{u}r die $r$ Werthenpaare
$(\gamma, \delta)$ noch f\"{u}r
\[ m(n - 2) + n(m - 2)- 2r \]
oder $2(p - 1)$ der Gleichung $F = 0$ gen\"{u}gende Werthenpaare
von $s$ und $z$ unendlich klein von der ersten Ordung.  Sind nun
\[ \varphi^{(1)}
   = \alpha_1^{(1)} \varphi_1 + \alpha_2^{(1)} \varphi_2 + \cdots
      + \alpha_p^{(1)} \varphi_p \]
und
\[ \varphi^{(2)}
   = \alpha_1^{(2)} \varphi_1 + \alpha_2^{(2)} \varphi_2 + \cdots
      + \alpha_p^{(2)} \varphi_p \]
zwei beliebige Functionen $\varphi$, so kann man in dem Ausdrucke
$\displaystyle\frac{\varphi^{(2)}}{\varphi^{(1)}}$
den Nenner so bestimmen, dass er f\"{u}r $p - 1$ beliebig
gegebene der Gleichung $F = 0$ gen\"{u}gende Werthenpaare von $s$
und $z$ gleich Null wird, und dann den Z\"{a}hler so, dass er
f\"{u}r $p - 2$ von den \"{u}brigen Werthenpaaren, f\"{u}r welche
$\varphi^{(1)}$ noch gleich $0$ wird, gleichfalls verschwindet.
Er ist dann noch eine lineare Function von zwei willk\"{u}rlichen
Constanten und folglich ein allgemeiner Ausdruck einer Function,
die nur f\"{u}r $p$ Punkte der Fl\"{a}che~$T$ unendlich von der
ersten Ordnung wird.  Eine Function, die f\"{u}r weniger als $p$
Punkte unendlich wird, bildet einen speciellen Fall dieser
Function; es lassen sich daher alle Functionen, die f\"{u}r
weniger als $p + 1$ Punkte der Fl\"{a}che~$T$ unendlich von der
ersten Ordnung werden, in der Form
$\displaystyle\frac{\varphi^{(2)}}{\varphi^{(1)}}$
oder in der form
$\displaystyle\frac{dw^{(2)}}{dw^{(1)}}$,
wenn $w^{(1)}$ und $w^{(2)}$ zwei allenthalben endliche Integrale
rationaler Functionen von $s$ und $z$ sind, darstellen.

\medbreak

\centerline{11.}

\nobreak\medskip

Eine wie $T$ verzweigte Function $z_1$ von $z$, die f\"{u}r $n_1$
Punkte dieser Fl\"{a}che unendlich von der ersten Ordnung wird,
ist nach dem Fr\"{u}heren (S.~\pageref{`Wurzel'}) die Wurzel einer
Gleichung von der Form
$G( {\buildrel n \over z_1}, {\buildrel n_1 \over z} ) = 0$
und nimmt daher jeden Werth f\"{u}r $n_1$ Punkte der Fl\"{a}che
$T$ an.  Wenn man sich also jeden Punkt von $T$ durch einen den
Werth von $z_1$ in diesem Punkte geometrisch repr\"{a}sentirenden
Punkt einer Ebene abgebildet denkt, so bildet die Gesammtheit
dieser Punkte eine in der $z_1$-Ebene allenthalben $n_1 \,$\-fach
ausgebreitete und die Fl\"{a}che $T$---bekanntlich in den
kleinsten Theilen \"{a}hnlich---abbildende Fl\"{a}che~$T_1$.
Jedem Punkt in der einen Fl\"{a}che entspricht dann \emph{ein}
Punkt in der andern.  Die Functionen $\omega$ oder die Integrale
wie $T$ verzweigter Functionen von $z$ gehen daher, wenn man
f\"{u}r $z$ als unabh\"{a}ngig ver\"{a}nderliche Gr\"{o}sse~$z_1$
einf\"{u}hrt, in Functionen \"{u}ber, welche in der
Fl\"{a}che~$T_1$ allenthalben \emph{einen} bestimmten Werth und
dieselben Unstetigkeiten haben, wie die Functionen~$\omega$ in den
entsprechenden Punkten von $T$, und welche folglich Integrale wie
$T_1$ verzweigter Functionen von $z_1$ sind.

Bezeichnet $s_1$ irgend eine andere wie $T$ verzweigte Function
von $z$, die f\"{u}r $m_1$ Punkte von $T$ und also auch von $T_1$
unendlich von der ersten Ordnung wird, so findet (\S.~5) zwischen
$s_1$ und $z_1$ eine Gleichung von der Form
\[ F_1( {\buildrel n_1 \over s_1}, {\buildrel m_1 \over z_1} ) = 0 \]
statt, worin $F_1$ eine Potenz einer unzerf\"{a}llbaren ganzen
Function von $s_1$ und $z_1$ ist, und es lassen sich, wenn diese
Potenz die erste ist, alle wie $T_1$ verzweigten Functionen von
$z_1$, folglich alle rationalen Functionen von $s$ und $z$
rational in $s_1$ und $z_1$ ausdr\"{u}cken (\S.~8).

Die Gleichung
$F( {\buildrel n \over s}, {\buildrel m \over z} ) = 0$
kann also durch eine rationale Substitution in
$F_1( {\buildrel n_1 \over s_1}, {\buildrel m_1 \over z_1} ) = 0$
und diese in jene transformirt werden.

Die Gr\"{o}ssengebiete $(s,z)$ und $(s_1, z_1)$ sind
gleichvielfach zusammenh\"{a}ngend, da jedem Punkte des einen
\emph{ein} Punkt des andern entspricht.  Bezeichnet daher $r_1$
die Anzahl der F\"{a}lle, in welchen $s_1$ und $z_1$ f\"{u}r zwei
verschiedene Punkte des Gr\"{o}ssengebiets $(s_1, z_1)$ beide
denselben Werth annehmen und folglich gleichzeitig $F_1$,
$\displaystyle \frac{\partial F_1}{\partial s_1}$
und
$\displaystyle \frac{\partial F_1}{\partial z_1}$
gleich $0$ und
\[ \frac{\partial^2 F_1}{\partial s_1^2}
      \frac{\partial^2 F_1}{\partial z_1^2}
    - \left( \frac{\partial^2 F_1}{\partial s_1 \, \partial z_1} \right)^2 \]
nicht Null ist, so muss
\[ (n_1 - 1)(m_1 - 1) - r_1 = p = (n - 1)(m - 1) - r \]
sein.

\medbreak

\centerline{12.}

\nobreak\medskip

Man betrachte nun als zu Einer \emph{Klasse} geh\"{o}rend
\emph{alle irreductiblen algebraischen Gleichungen zwischen zwei
ver\"{a}nderlichen Gr\"{o}ssen, welche sich durch rationale
Substitutionen in einander transformiren lassen}, so dass
\[ F(s,z) = 0 \mbox{ und } F_1(s_1, z_1) = 0 \]
zu derselben Klasse geh\"{o}ren, wenn sich f\"{u}r $s$ und $z$
solche rationale Functionen von $s_1$ und $z_1$ setzen lassen,
dass $F(s,z) = 0$ in $F_1(s_1, z_1) = 0$ \"{u}bergeht und
zugleich $s_1$ und $z_1$ rationale Functionen von $s$ und $z$
sind.

Die rationalen Functionen von $s$ und $z$ bilden, als Functionen
von irgend einer von ihnen $\zeta$ betrachtet, ein System
gleichverzweigter algebraischer Functionen.  Auf diese Weise
f\"{u}hrt jede Gleichung offenbar zu einer Klasse von Systemen
gleichverzweigter algebraischer Functionen, welche sich durch
Einf\"{u}hrung einer Function des Systems als unabh\"{a}ngig
ver\"{a}nderlicher Gr\"{o}sse in einander transformiren lassen
und zwar alle Gleichungen Einer Klasse zu derselben Klasse von
Systemen algebraischer Functionen, und umgekehrt f\"{u}hrt
(\S.~11) jede Klasse von solchen Systemen zu Einer Klasse von
Gleichungen.

Ist das Gr\"{o}ssengebiet $(s,z)$ $\overline{2p + 1} \,$\-fach
zusammenh\"{a}ngend und die Function $\zeta$ in $\mu$ Punkten
desselben unendlich von der ersten Ordnung, so ist die Anzahl der
Verzweigungswerthe der gleichverzweigten Functionen von $\zeta$,
welche durch die \"{u}brigen rationalen Functionen von $s$ und
$z$ gebildet werden, $2(\mu + p - 1)$, und die Anzahl der
willk\"{u}rlichen Constanten in der Function~$\zeta$
$2\mu - p + 1$ (\S.~5).  Diese lassen sich so bestimmen, dass
$2\mu - p + 1$ Verzweigungswerthe gegebene Werthe annehmen, wenn
diese Verzweigungswerthe von einander unabh\"{a}ngige Functionen
von ihnen sind, und zwar nur auf eine endliche Anzahl Arten, da
die Bedingungsgleichungen algebraisch sind.  In jeder Klasse von
Systemen gleichverzweigter $\overline{2p + 1} \,$\-fach
zusammenh\"{a}ngender Functionen giebt es daher eine endliche
Anzahl von Systemen $\mu \,$\-werthiger Functionen, in welchen
$2\mu - p + 1$ Verzweigungswerthe gegebene Werthe annehmen.  Wenn
andererseits die $2(\mu + p - 1)$ Verzweigungspunkte einer die
$\zeta$-Ebene allenthalben $\mu \,$\-fach bedeckenden
$\overline{2p + 1} \,$\-fach zusammenh\"{a}ngenden Fl\"{a}che beliebig
gegeben sind, so giebt es (\S\S.~3--5) immer ein System wie diese
Fl\"{a}che verzweigter algebraischer Functionen von $\zeta$.  Die
$3p - 3$ \"{u}brigen Verzweigungswerthe in jenen Systemen
gleichverzweigter $\mu \,$\-werthiger Functionen k\"{o}nnen daher
beliebige Werthe annehmen; und es h\"{a}ngt also eine Klasse von
Systemen gleichverzweigter $\overline{2p + 1} \,$\-fach
zusammenh\"{a}ngender Functionen und die zu ihr geh\"{o}rende
Klasse algebraischer Gleichungen von $3p - 3$ stetig
ver\"{a}nderlichen Gr\"{o}ssen ab, welche die Moduln dieser
Klasse genannt werden sollen.

Diese Bestimmung der Anzahl der Moduln einer Klasse
$\overline{2p + 1} \,$\-fach zusammenh\"{a}ngender algebraischer Functionen
gilt jedoch nur unter der Voraussetzung, dass es $2\mu - p + 1$
Verzweigungswerthe giebt, welche von einander unabh\"{a}ngige
Functionen der willk\"{u}rlichen Constanten in der Function
$\zeta$ sind.  Diese Voraussetzung trifft nur zu, wenn $p > 1$,
und die Anzahl der Moduln ist nur dann $= 3p - 3$, f\"{u}r
$p = 1$ aber $= 1$.  Die directe Untersuchung derselben wird
indess schwierig durch die Art und Weise, wie die
willk\"{u}rlichen Constanten in $\zeta$ enthalten sind.  Man
f\"{u}hre deshalb in einem Systeme gleichverzweigter
$\overline{2p + 1} \,$\-fach zusammenh\"{a}ngender Functionen, um die
Anzahl der Moduln zu bestimmen, als unabh\"{a}ngig
ver\"{a}nderliche Gr\"{o}sse nicht eine dieser Functionen,
sondern ein allenthalben endliches Integral einer solchen
Function ein.

Die Werthe, welche die Function $w$ von $z$ innerhalb der
Fl\"{a}che $T'$ annimmt, werden geometrisch repr\"{a}sentirt durch
eine einen endlichen Theil der $w$-Ebene einfach oder mehrfach
bedeckende und die Fl\"{a}che~$T'$ (in den kleinsten Theilen
\"{a}hnlich) abbildende Fl\"{a}che, welche durch $S$ bezeichnet
werden soll.  Da $w$ auf der positiven Seite des $\nu \,$\-ten
Querschnitts um die Constante $k^{(\nu)}$ gr\"{o}sser ist, als
auf der negativen, so besteht die Begrenzung von $S$ aus Paaren
von parallelen Curven, welche denselben Theil des $T'$
begrenzenden Schnittsystems abbilden, und es wird die
Ortsverschiedenheit der entsprechenden Punkte in den parallelen,
den $\nu \,$\-ten Querschnitt abbildenden Begrenzungstheilen von $S$
durch die complexe Gr\"{o}sse $k^{(\nu)}$ ausgedr\"{u}ckt.  Die
Anzahl der einfachen Verzweigungspunkte der Fl\"{a}che~$S$ ist
$2p - 2$, da $dw$ in $2p - 2$ Punkten der Fl\"{a}che~$T$
unendlich klein von der zweiten Ordnung wird.  Die rationalen
Functionen von $s$ und $z$ sind dann Functionen von $w$, welche
f\"{u}r jeden Punkt von $S$ Einen bestimmten, wo sie nicht
unendlich werden, stetig sich \"{a}ndernden Werth haben und in
den entsprechenden Punkten paralleler Begrenzungstheile denselben
Werth annehmen.  Sie bilden daher ein System gleichverzweigter
und $2p \,$\-fach periodischer Functionen von $w$.  Es l\"{a}sst
sich nun (auf \"{a}hnlichem Wege, wie in den \S\S.~3--5) zeigen,
dass, die $2p - 2$ Verzweigungspunkte und die $2p$
Ortsverschiedenheiten paralleler Begrenzungs\-theile der
Fl\"{a}che~$S$ als willk\"{u}rlich gegeben vorausgesetzt, immer
ein System wie diese Fl\"{a}che verzweigter Functionen existirt,
welche in den entsprechenden Punkten paralleler Begrenzungstheile
denselben Werth annehmen und also $2p \,$\-fach periodisch sind, und
die, als Functionen von einer von ihnen betrachtet, ein System
gleichverzweigter $\overline{2p + 1} \,$\-fach zusammenh\"{a}ngender
algebraischer Functionen bilden, folglich zu einer Klasse von
$\overline{2p + 1} \,$\-fach zusammenh\"{a}ngenden algebraischen
Functionen f\"{u}hren.  In der That ergiebt sich nach dem
\emph{Dirichlet}'schen Princip, dass in der Fl\"{a}che~$S$ eine
Function von $w$ bis auf eine additive Constante bestimmt ist
durch die Bedingungen, im Innern von $S$ beliebig gegebene
Unstetigkeiten von der Form wie $\omega$ in $T'$ anzunehmen und
in den entsprechenden Punkten paralleler Begrenzungstheile um
Constanten, deren reeller Theil gegeben ist, verschiedene Werthe
zu erhalten.  Hieraus schliesst man, \"{a}hnlich wie im \S.~5,
die M\"{o}glichkeit von Functionen, welche nur in einzelnen
Punkten von $S$ unstetig werden und in den entsprechenden
Punkten paralleler Begrenzungstheile denselben Werth annehmen.
Wird eine solche Function~$z$ in $n$ Punkten von $S$ unendlich
von der ersten Ordnung und sonst nicht unstetig, so nimmt sie
jeden complexen Werth in $n$ Punkten von $S$ an; denn wenn $a$
eine beliebige Constante ist, so ist
$\int \partial \log (z - a)$,
um $S$ erstreckt, $= 0$, da die Integration durch parallele
Begrenzungstheile sich aufhebt, und es wird daher $z - a$ in $S$
ebenso oft unendlich klein, als unendlich von der ersten Ordnung.
Die Werthe, welche $z$ annimmt, werden folglich durch eine
\"{u}ber die $z$-Ebene allenthalben $n \,$\-fach ausgebreitete
Fl\"{a}che repr\"{a}sentirt, und die \"{u}brigen ebenso
verzweigten und periodischen Functionen von $w$ bilden daher ein
System wie diese Fl\"{a}che verzweigter $\overline{2p + 1} \,$\-fach
zusammenh\"{a}ngender algebraischer Functionen von~$z$,
w.~z.~b.~w.

F\"{u}r eine beliebig gegebene Klasse $\overline{2p + 1} \,$\-fach
zusammenh\"{a}ngender algebraischer Functionen kann man nun in
dem als unabh\"{a}ngig ver\"{a}nderliche Gr\"{o}sse
einzuf\"{u}hrenden
\[ w = \alpha_1 w_1 + \alpha_2 w_2 + \cdots + \alpha_p w_p + c \]
die Gr\"{o}ssen $\alpha$ so bestimmen, dass $p$ von den $2p$
Periodicit\"{a}tsmoduln gegebene Werthe annehmen, und $c$ wenn
$p > 1$ so, dass einer von den $2p - 2$ Verzeigungswerthen der
periodischen Functionen von $w$ einen gegebenen Werth erh\"{a}lt.
Dadurch ist $w$ v\"{o}llig bestimmt, und also sind es auch die
$3p - 3$ \"{u}brigen Gr\"{o}ssen, von denen die Verzweigungsart
und Periodicit\"{a}t jener Functionen von $w$ abh\"{a}ngt; und da
jedweden Werthen dieser $3p - 3$ Gr\"{o}ssen eine Klasse von
$\overline{2p + 1} \,$\-fach zusammenh\"{a}ngenden algebraischen Functionen
entspricht, so h\"{a}ngt eine solche von $3p - 3$ unabh\"{a}ngig
ver\"{a}nderlichen Gr\"{o}ssen ab.

Wenn $p = 1$ ist, so ist kein Verzweigungspunkt vorhanden, und es
l\"{a}sst sich in
\[ w = \alpha_1 w_1 + c \]
die Gr\"{o}sse $\alpha_1$ so bestimmen, dass \emph{ein}
Periodicit\"{a}tsmodul einen gegebenen Werth erh\"{a}lt und
dadurch ist der andere Periodicit\"{a}tsmodul bestimmt.  Die
Anzahl der Moduln einer Klasse ist also dann $= 1$.

\medbreak

\centerline{13.}

\nobreak\medskip

Nach den obigen (im \S.~11 entwickelten) Principien der
Transformation muss man, um eine beliebig gegebene Gleichung
$F(s,z) = 0$ durch eine rationale Substitution in eine Gleichung
derselben Klasse
$F_1( {\buildrel n_1 \over s_1}, {\buildrel m_1 \over z_1} ) = 0$
von m\"{o}glichst niedrigem Grade zu transformiren, zuerst
f\"{u}r $z_1$ einen rationalen Ausdruck in $s$ und $z$, $r(s,z)$,
so bestimmen, dass $n_1$ m\"{o}glichst klein wird, und dann $s_1$
gleich einem andern rationalen Ausdrucke $r'(s,z)$ so, dass $m_1$
m\"{o}glichst klein wird und zugleich die zu einem beliebigen
Werthe von $z_1$ geh\"{o}rigen Werthe von $s_1$ nicht in Gruppen
unter einander gleicher zerfallen, so dass
$F_1( {\buildrel n_1 \over s_1}, {\buildrel m_1 \over z_1} )$
nicht eine h\"{o}here Potenz einer unzerf\"{a}llbaren Function
sein kann.

Wenn das Gr\"{o}ssengebiet $(s,z)$ $\overline{2p + 1} \,$\-fach
zusammenh\"{a}ngend ist, so ist der kleinste Werth, den $n_1$
annehmen kann, allgemein zu reden,
$\displaystyle \geq \frac{p}{2} + 1$
(\S.~5) und die Anzahl der F\"{a}lle, in denen $s_1$ und $z_1$
f\"{u}r zwei verschiedene Punkte des Gr\"{o}ssengebiets beide
denselben Werth annehmen,
\[  = (n_1 - 1)(m_1 - 1) - p.\]

In einer Klasse von algebraischen Gleichungen zwischen zwei
ver\"{a}nderlichen Gr\"{o}ssen haben demnach, wenn ihre Moduln
nicht besonderen Bedingungsgleichungen gen\"{u}gen, die
Gleichungen niedrigsten Grades folgende Form:
\[ \begin{array}{lllll}
   & \mbox{f\"{u}r} & p = 1,
   & F( {\buildrel 2 \over s}, {\buildrel 2 \over z} ) = 0,
   & r = 0 \\
   &  & p = 2,
   & F( {\buildrel 2 \over s}, {\buildrel 3 \over z} ) = 0,
   & r = 0 \\
   &  & p = 2\mu - 3,
   & F( {\buildrel \mu \over s}, {\buildrel \mu \over z} ) = 0,
   & r = (\mu - 2)^2 \\
p > 2 \\
   &  & p = 2\mu - 2,
   & F( {\buildrel \mu \over s}, {\buildrel \mu \over z} ) = 0,
   & r = (\mu - 1)(\mu - 3).
\end{array} \]

Von den Coefficienten der Potenzen von $s$ und $z$ in den ganzen
Functionen $F$ m\"{u}ssen $r$ als lineare homogene Functionen der
\"{u}brigen so bestimmt werden, dass
$\displaystyle \frac{\partial F}{\partial s}$
und
$\displaystyle \frac{\partial F}{\partial z}$
f\"{u}r $r$ der Gleichung $F = 0$ gen\"{u}gende Werthenpaare
gleichzeitig verschwinden.  Die rationalen Functionen von $s$ und
$z$, als Functionen von einer von ihnen betrachtet, stellen dann
alle Systeme $\overline{2p + 1} \,$\-fach zusammenh\"{a}ngender
algebraischer Functionen dar.

\medbreak

\centerline{14.}

\nobreak\medskip

Ich benutze nun nach \emph{Jacobi} (dieses Journal Bd.~9
Nr.~32 \S.~8) das \emph{Abel}'sche Additionstheorem zur
Integration eines Systems von Differentialgleichungen; ich werde
mich dabei auf das beschr\"{a}nken, was in dieser Abhandlung
sp\"{a}ter n\"{o}thig ist.

F\"{u}hrt man in einem allenthalben endlichen Integrale $w$ einer
rationalen Function von $s$ und $z$ als unabh\"{a}ngig
ver\"{a}nderliche Gr\"{o}sse eine rationale Function von $s$ und
$z$, $\zeta$, ein, die f\"{u}r $m$ Werthenpaare von $s$ und $z$
unendlich von der ersten Ordnung wird, so ist
$\displaystyle \frac{\partial w}{\partial z}$
eine $m \,$\-werthige Function von $\zeta$.  Bezeichnet man die $m$
Werthe von $w$ f\"{u}r dasselbe $\zeta$ durch
$w^{(1)}, w^{(2)},\ldots, w^{(m)}$, so ist
\[    \frac{\partial w^{(1)}}{\partial \zeta}
    + \frac{\partial w^{(2)}}{\partial \zeta} + \cdots
    + \frac{\partial w^{(m)}}{\partial \zeta} \]
eine einwerthige Function von $\zeta$, deren Integral
allenthalben endlich bleibt, und folglich ist auch
\[ \int \partial ( w^{(1)} + w^{(2)} + \cdots + w^{(m)} ) \]
allenthalben einwerthig und endlich, mithin constant.  Auf
\"{a}hnliche Weise findet sich, wenn
$\omega^{(1)}, \omega^{(2)},\ldots, \omega^{(m)}$
die demselben $\zeta$ entsprechenden Werthe eines beliebigen
Integrals $\omega$ einer rationalen Function von $s$ und $z$
bezeichnen,
\[ \int \partial ( \omega^{(1)} + \omega^{(2)} + \cdots + \omega^{(m)}) \]
bis auf eine additive Constante aus den Unstetigkeiten von
$\omega$ und zwar als Summe von einer rationalen Function und mit
constanten Coefficienten versehenen Logarithmen rationaler
Functionen von $\zeta$.

Mittelst dieses Satzes lassen sich, wie jetzt gezeigt werden
soll, folgende $p$ gleichzeitige Differentialgleichungen
zwischen den $p + 1$ der Gleichung $F(s, z) = 0$ gen\"{u}genden
Werthenpaaren von $s$ und $z$,
$(s_1, z_1)$, $(s_2, z_2),\ldots$, $(s_{p+1}, z_{p+1})$
\[       \frac{\varphi_\pi(s_1, z_1) \, \partial z_1}{\displaystyle
            \frac{\partial F(s_1, z_1)}{\partial s_1}}
       + \frac{\varphi_\pi(s_2, z_2) \, \partial z_2}{\displaystyle
            \frac{\partial F(s_2, z_2)}{\partial s_2}}
       + \frac{\varphi_\pi(s_{p+1}, z_{p+1}) \, \partial z_{p+1}}{\displaystyle
            \frac{\partial F(s_{p+1}, z_{p+1})}{\partial s_{p+1}}}
   = 0 \]
f\"{u}r $\pi = 1,2,\ldots, p$, allgemein oder vollst\"{a}ndig
(complete) integriren.

Durch diese Differentialgleichungen sind $p$ von den
Gr\"{o}ssenpaaren $(s_\mu, z_\mu)$ als Functionen des einen noch
\"{u}brigen v\"{o}llig bestimmt, wenn f\"{u}r einen beliebigen
Werth des letzteren die Werthe der \"{u}brigen gegeben werden.
Wenn man also diese $p + 1$ Gr\"{o}ssenpaare als Functionen
\emph{einer} ver\"{a}nderlichen Gr\"{o}sse $\zeta$ so bestimmt,
dass sie f\"{u}r denselben Werth $0$ dieser Gr\"{o}sse
\emph{beliebig} gegebene Anfangswerthe
$(s_1^0, z_1^0)$, $(s_2^0, z_2^0),\ldots$, $(s_{p+1}^0, z_{p+1}^0)$
annehmen und den Differentialgleichungen gen\"{u}gen, so hat man
dadurch die Differentialgleichungen allgemein integrirt.  Nun
l\"{a}sst sich die Gr\"{o}sse
$\displaystyle\frac{1}{\zeta}$
als einwerthige und folglich rationale Function von $(s, z)$
immer so bestimmen, dass sie nur f\"{u}r alle oder einige von den
$p + 1$ Werthenpaaren $(s_\mu^0, z_\mu^0)$ unendlich und f\"{u}r
diese nur unendlich von der ersten Ordnung wird, da sich in dem
Ausdrucke
\[      \sum_{\mu=1}^{\mu = p + 1} \beta_\mu  t (s_\mu^0, z_\mu^0)
      + \sum_{\mu=1}^{\mu = p}     \alpha_\mu w_\mu
      + \mbox{const.} \]
die Verh\"{a}ltnisse der Gr\"{o}ssen $\alpha$ und $\beta$ immer
so bestimmen lassen, dass die Periodicit\"{a}tsmoduln
s\"{a}mmtlich $0$ werden.  Es gen\"{u}gen dann, wenn kein $\beta
= 0$ ist, den zu l\"{o}senden Differentialgleichungen die $p + 1$
Zweige der $\overline{p + 1} \,$\-werthigen gleichverzweigten Functionen $s$
und $z$ von $\zeta$,
\[ (s_1, z_1),\enspace
   (s_2, z_2),\ldots,\enspace
   (s_{p+1}, z_{p+1}),\]
welche f\"{u}r $\zeta = 0$ die Werthe
\[ (s_1^0, z_1^0),\enspace
   (s_2^0, z_2^0),\ldots,\enspace
   (s_{p+1}^0, z_{p+1}^0) \]
annehmen.  Wenn aber von der Gr\"{o}ssen $\beta$ einige, etwa die
$p + 1 - m$ letzten gleich $0$ werden, so werden die zu
l\"{o}senden Differentialgleichungen befriedigt durch die $m$
Zwiege der $m$ werthigen Functionen $s$ und $z$ von $\zeta$,
\[ (s_1, z_1),\enspace
   (s_2, z_2),\ldots,\enspace
   (s_m, z_m),\]
welche f\"{u}r $\zeta = 0$ gleich
\[ (s_1^0, z_1^0),\enspace
   (s_2^0, z_2^0),\ldots,\enspace
   (s_m^0, z_m^0) \]
werden, und durch constante, also ihren Anfangswerthen
$s_{m+1}^0,\ldots, z_{p+1}^0$ gleiche Werthe der Gr\"{o}ssen
$s_{m+1}, z_{m+1};\ldots; s_{p+1}, z_{p+1}$.  Im letzteren Falle
sind von den $p$ linearen homogenen Gleichungen
\[ \sum_{\mu = 1}^{\mu = m}
         \frac{\varphi_\pi(s_\mu, z_\mu) \, \partial z_\mu}{\displaystyle
            \frac{\partial F(s_\mu, z_\mu)}{\partial s_\mu}}
   = 0 \]
f\"{u}r $\pi = 1,2,\ldots, p$ zwischen den Gr\"{o}ssen
\[ \frac{\partial z_\mu}{\displaystyle
      \frac{\partial F(s_\mu, z_\mu)}{\partial s_\mu}} \]
$p + 1 - m$ eine Folge der \"{u}brigen; es ergeben sich hieraus
$p + 1 - m$ Bedingungsgleichungen, welche, damit dieser Fall
eintritt, zwischen den Functionen
$(s_1, z_1),\enspace (s_2, z_2),\ldots,\enspace (s_m, z_m)$,
und also auch zwischen ihren Anfangswerthen
$(s_1^0, z_1^0),\enspace (s_2^0, z_2^0),\ldots,\enspace (s_m^0, z_m^0)$
erf\"{u}llt sein m\"{u}ssen, und es k\"{o}nnen daher von diesen,
wie oben (\S.~5) gefunden, nur $2m - p - 1$ beliebig gegeben
werden.

\medbreak

\centerline{15.}

\nobreak\medskip

Es sei nun
\[ \int  \frac{\varphi_\pi(s, z) \, \partial z}{\displaystyle
            \frac{\partial F(s, z)}{\partial s}}
      + \mbox{const.},\]
durch das Innere von $T'$ integrirt, gleich $w_\pi$ und der
Periodicit\"{a}tsmodul von $w_\pi$ f\"{u}r den
$\nu \,$\-ten Querschnitt gleich $k_\pi^{(\nu)}$, so dass sich die
Functionen $w_1, w_2,\ldots, w_p$ des Gr\"{o}ssenpaars $(s,z)$
beim Uebertritt des Punkts $(s,z)$ von der negativen auf die
positive Seite des $\nu \,$\-ten Querschnitts gleichzeitig um
$k_1^{(\nu)}, k_2^{(\nu)},\ldots k_p^{(\nu)}$
\"{a}ndern.  Zur Abk\"{u}rzung mag ein System von $p$ Gr\"{o}ssen
$(b_1, b_2,\ldots, b_p)$ einem andern
$(a_1, a_2,\ldots, a_p)$ \emph{congruent nach $2p$ Systemen
zusammengeh\"{o}riger Moduln} genannt werden, wenn es aus ihm
durch gleichzeitige Aenderungen s\"{a}mmtlicher Gr\"{o}ssen um
zusammengeh\"{o}rige Moduln erhalten werden kann.  Ist der Modul
der $\pi \,$\-ten Gr\"{o}sse im $\nu \,$\-ten Systeme
$= k_\pi^{(\nu)}$, so heisst demnach
\[ (b_1, b_2,\ldots, b_p) \equiv (a_1, a_2,\ldots, a_p),\]
wenn
\[ b_\pi = a_\pi + \sum_{\nu = 1}^{\nu = 2p} m_\nu k_\pi^{(\nu)} \]
f\"{u}r $\pi = 1,2,\ldots, p$ und $m_1, m_2,\ldots, m_{2p}$ ganze
Zahlen sind.

Da sich $p$ beliebige Gr\"{o}ssen $a_1, a_2,\ldots, a_p$ immer
und nur auf eine Weise in die Form
$a_\pi = \sum\limits_{\nu = 1}^{\nu = 2p} \xi_\nu k_\pi^{(\nu)}$
setzen lassen, so dass die $2p$ Gr\"{o}ssen $\xi$ reell sind, und
durch Aenderung dieser Gr\"{o}ssen $\xi$ um ganze Zahlen alle
congruenten Systeme und nur diese sich ergeben, so erh\"{a}lt man
aus jeder Reihe congruenter Systeme eins und nur eins, wenn man
in diesen Ausdr\"{u}cken jede Gr\"{o}sse~$\xi$ alle Werthe von
einem beliebigen Werthe bis zu einem um $1$ gr\"{o}sseren, einen
der beiden Grenzwerthe eingeschlossen, stetig durchlaufen
l\"{a}sst.

Dieses festgesetzt, folgt aus den obigen Differentialgleichungen
oder aus den $p$ Gleichungen
\[ \sum_{\mu = 1}^{\mu = p + 1} dw_\pi^{(\mu)} = 0
   \mbox{ f\"{u}r } \pi = 1,2,\ldots, p \]
durch Integration
\[    \left(
         {\textstyle \sum} w_1^{(\mu)},
         {\textstyle \sum} w_2^{(\mu)},\ldots,
         {\textstyle \sum} w_p^{(\mu)}
      \right)
   \equiv (c_1, c_2,\ldots, c_p),\]
worin $c_1, c_2,\ldots, c_p$ constante von den Werthen
$(s^0, z^0)$ abh\"{a}ngige Gr\"{o}ssen sind.

\medbreak

\centerline{16.}

\nobreak\medskip

Dr\"{u}ckt man $\zeta$ als Quotienten zweier ganzen Functionen
von $s$ und $z$,
$\displaystyle \frac{\chi}{\psi}$,
aus, so sind die Gr\"{o}ssenpaare
$(s_1, z_1),\ldots$, $(s_m, z_m)$
die gemeinschaftlichen Wurzeln der Gleichungen $F = 0$ und
$\displaystyle \frac{\chi}{\psi} = \zeta$.
Da die ganze Function
\[ \chi - \zeta \psi = f(s,z) \]
f\"{u}r alle Werthenpaare, f\"{u}r welche $\chi$ und $\psi$
gleichzeitig verschwinden ebenfalls, was auch $\zeta$ sei,
verschwindet, so k\"{o}nnen die Gr\"{o}ssenpaare
$(s_1, z_1),\ldots$, $(s_m, z_m)$
auch definirt werden als gemeinschaftliche Wurzeln der Gleichung
$F = 0$ und einer Gleichung $f(s,z) = 0$, deren Coefficienten so
sich \"{a}ndern, dass alle \"{u}brigen gemeinschaftlichen Wurzeln
constant bleiben.  Wenn $m < p + 1$, kann $\zeta$ in der Form
$\displaystyle \frac{\varphi^{(1)}}{\varphi^{(2)}}$
dargestellt werden (\S.~10) und $f$ in der Form
\[ \varphi^{(1)} - \zeta \varphi^{(2)} = \varphi^{(3)}.\]
Die allgemeinsten Werthe der den $p$ Gleichungen
\[ \sum_{\mu = 1}^{\mu = p} dw_\pi^{(\mu)} = 0
   \mbox{ f\"{u}r } \pi = 1,2,\ldots, p \]
gen\"{u}genden Functionenpaare
$(s_1, z_1),\ldots, (s_p, z_p)$
werden daher gebildet durch $p$ gemeinschaftliche Wurzeln der
Gleichungen $F = 0$ und $\varphi = 0$, welche so sich \"{a}ndern,
dass die \"{u}brigen gemeinschaftlichen Wurzeln constant bleiben.
Hieraus folgt leicht der sp\"{a}ter n\"{o}thige Satz, dass die
Aufgabe, $p - 1$ von den $2p - 2$ Gr\"{o}ssenpaaren
$(s_1, z_1),\ldots$, $(s_{2p-2}, z_{2p-2})$ als Functionen der
$p - 1$ \"{u}brigen so zu bestimmen, dass die $p$ Gleichungen
\[ \sum_{\mu = 1}^{\mu = 2p - 2} dw_\pi^{(\mu)} = 0
   \mbox{ f\"{u}r } \pi = 1,2,\ldots, p \]
erf\"{u}llt werden, v\"{o}llig allgemein gel\"{o}st wird, wenn
man f\"{u}r diese $2p - 2$ Gr\"{o}ssenpaare die von den $r$
Wurzeln $s = \gamma_\varrho$, $z = \delta_\varrho$ (\S.~6)
verschiedenen gemeinschaftlichen Wurzeln der Gleichungen $F = 0$
und $\varphi = 0$ oder die $2p - 2$ Werthenpaare nimmt, f\"{u}r
welche $dw$ unendlich klein von der zweiten Ordnung wird, und
dass diese Aufgabe daher nur \emph{eine} L\"{o}sung zul\"{a}sst.
Solche Gr\"{o}ssenpaare sollen \emph{durch die Gleichung
$\varphi = 0$ verkn\"{u}pft} heissen.  In Folge der Gleichungen
\[ \sum_1^{2p - 2} dw_\pi^{(\mu)} = 0 \]
wird
\[ \left(
      \sum_1^{2p - 2} w_1^{(\mu)}, \,
      \sum_1^{2p - 2} w_2^{(\mu)}, \, \ldots, \,
      \sum_1^{2p - 2} w_p^{(\mu)}
   \right),\]
die Summe \"{u}ber solche Gr\"{o}ssenpaare ausgedehnt, congruent
einem constanten Gr\"{o}ssensysteme $(c_1, c_2,\ldots, c_p)$,
worin $c_\pi$ nur von der additiven Constante in der
Function~$w_\pi$ oder dem Anfangswerthe des sie ausdr\"{u}ckenden
Integrals abh\"{a}ngt.

\begin{center}
\large\bfseries
Zweite Abteilung.
\end{center}

\centerline{17.}

\nobreak\medskip

F\"{u}r die ferneren Untersuchungen \"{u}ber Integrale von
algebraischen,\break
$\overline{2p + 1} \,$\-fach zusammenh\"{a}ngenden
Functionen ist die Betrachtung einer $p \,$\-fach unendlichen
$\vartheta$-Reihe von grossem Nutzen, d.~h.\ einer $p \,$\-fach
unendlichen Reihe, in welcher der Logarithmus des allgemeinen
Gliedes eine ganze Function zweiten Grades der Stellenzeiger ist.
Es sei in dieser Function f\"{u}r ein Glied, dessen Stellenzeiger
$m_1, m_2,\ldots, m_p$
sind, der Coefficient des Quadrats $m_\mu^2$ gleich
$a_{\mu, \mu}$, des doppelten Products $m_\mu m_{\mu'}$ gleich
$a_{\mu, \mu'} = a_{\mu', \mu}$, der doppelten Gr\"{o}sse $m_\mu$
gleich $v_\mu$, und das constante Glied $= 0$.  Die Summe der
Reihe, \"{u}ber alle ganzen positiven oder negativen Werthe der
Gr\"{o}ssen $m$ ausgedehnt, werde als Function der $p$
Gr\"{o}ssen $v$ betrachtet und durch
$\vartheta(v_1, v_2,\ldots, v_p)$
bezeichnet, so dass
\begin{equation}
\label{eqn-1}
   \vartheta(v_1, v_2,\ldots, v_p)
   = \left( \sum_{-\infty}^\infty \right)^p
      e^{\left( \sum\limits_1^p \right)^2 a_{\mu, \mu'} m_\mu m_{\mu'}
            + 2 \sum\limits_1^p v_\mu m_\mu},
\end{equation}
worin die Summationen im Exponenten sich auf $\mu$ und $\mu'$,
die \"{a}usseren Summationen auf $m_1, m_2,\ldots, m_p$ beziehen.
Damit diese Reihe convergirt, muss der reelle Theil von
$\left( \sum\limits_1^p \right)^2 a_{\mu, \mu'} m_\mu m_{\mu'}$
wesentlich negativ sein oder, als eine Summe von positiven oder
negativen Quadraten reeller linearer von einander
unabh\"{a}ngiger Functionen der Gr\"{o}ssen~$m$ dargestellt, aus
$p$ negativen Quadraten zusammengesetzt sein.

Die Function $\vartheta$ hat die Eigenschaft, dass es Systeme von
gleichzeitigen Aenderungen der $p$ Gr\"{o}ssen $v$ giebt, durch
welche $\log \vartheta$ nur um eine lineare Function der
Gr\"{o}ssen~$v$ ge\"{a}ndert wird, und zwar $2p$ von einander
unabh\"{a}ngige Systeme (d.~h.\ von denen keins eine Folge der
\"{u}brigen ist).  Denn man hat, die unge\"{a}ndert bleibenden
Gr\"{o}ssen~$v$ unter dem Functionszeichen $\vartheta$
weglassend, f\"{u}r $\mu = 1,2,\ldots,p$
\begin{equation}
\label{eqn-2}
\vartheta = \vartheta( v_\mu + \pi i) \quad \mbox{und}
\end{equation}
\begin{equation}
\label{eqn-3}
\vartheta = e^{2 v_\mu + a_{\mu, \mu}}
         \vartheta( v_1 + a_{1,\mu}, v_2 + a_{2,\mu},\ldots, v_p + a_{p,\mu}),
\end{equation}
wie sich sofort ergiebt, wenn man in der Reihe f\"{u}r
$\vartheta$ den Stellenzeiger $m_\mu$ in $m_\mu + 1$ verwandelt,
wodurch sie, w\"{a}hrend ihr Werth unge\"{a}ndert bleibt, in den
Ausdruck zur Rechten \"{u}bergeht.

Die Function $\vartheta$ ist durch diese Relationen und durch die
Eigenschaft, allenthalben endlich zu bleiben, bis auf einen
constanten Factor bestimmt.  Denn in Folge der letzteren
Eigenschaft und der Relationen (\ref{eqn-2}) ist sie eine
einwerthige, f\"{u}r endliche $v$ endliche Function von
$e^{2 v_1}, e^{2 v_2},\ldots, e^{2v_p}$
und folglich in eine $p \,$\-fach unendliche Reihe von der Form
\[ \left( \sum_{-\infty}^\infty \right)^p
      A_{m_1, m_2,\ldots, m_p}
      e^{2 \sum\limits_1^p v_\mu m_\mu} \]
mit den constanten Coefficienten $A$ entwickelbar.  Aus den
Relationen (\ref{eqn-3}) ergiebt sich aber
\[ A_{m_1,\ldots, m_\nu + 1,\ldots, m_p}
   = A_{m_1,\ldots, m_\nu,\ldots, m_p}
      e^{2 \sum\limits_1^p a_{\mu, \nu} m_\mu + a_{\nu, \nu}} \]
folglich
\[ A_{m_1,\ldots, m_p}
   = \mbox{const.} e^{ \left( \sum\limits_1^p \right)^2
      a_{\mu, \mu'} m_\mu m_{\mu'}},
\quad\mbox{w.~z.~b.~w.} \]

Man kann daher diese Eigenschaften der Function zu ihrer
Definition verwenden.  Die Systeme gleichzeitiger Aenderungen der
Gr\"{o}ssen~$v$, durch welche sich $\log \vartheta$ nur um eine
lineare Function von ihnen \"{a}ndert, sollen \emph{Systeme
zusammengeh\"{o}riger Periodicit\"{a}tsmoduln der unabh\"{a}ngig
ver\"{a}nderlichen Gr\"{o}ssen} in dieser $\vartheta$-Function
genannt werden.

\medbreak

\centerline{18.}

\nobreak\medskip

Ich substituire nun f\"{u}r die $p$ Gr\"{o}ssen
$v_1, v_2,\ldots, v_p$, $p$ immer endlich bleib\-ende Integrale
$u_1, u_2,\ldots, u_p$ rationaler Functionen einer
ver\"{a}nderlichen Gr\"{o}sse~$z$ und einer $\overline{2p + 1} \,$\-fach
zusammenh\"{a}ngenden algebraischen Function~$s$ dieser
Gr\"{o}sse, und f\"{u}r die zusammengeh\"{o}rigen
Periodicit\"{a}tsmoduln der Gr\"{o}ssen~$v$ zusammengeh\"{o}rige
(d.~h.\ an demselben Querschnitte stattfindende)
Periodicit\"{a}tsmoduln dieser Integrale, so dass
$\log \vartheta$ in einer Function \emph{einer}
Ver\"{a}nderlichen $z$ \"{u}bergeht, welche sich, wenn $s$ und
$z$ nach beliebiger stetiger Aenderung von $z$ den vorigen Werth
wieder annehmen, um lineare Functionen der Gr\"{o}ssen $u$
\"{a}ndert.

Es soll zun\"{a}chst gezeigt werden, dass eine solche
Substitution f\"{u}r jede $\overline{2p + 1} \,$\-fach zusammenh\"{a}ngende
Function~$s$ m\"{o}glich ist.  Die Zerschneidung der
Fl\"{a}che~$T$ muss zu diesem Zwecke so durch $2p$ in sich
zur\"{u}cklaufende Schnitte
$a_1, a_2,\ldots, a_p$, $b_1, b_2,\ldots, b_p$
geschehen, dass folgende Bedingungen erf\"{u}llt werden.  Wenn
man $u_1, u_2,\ldots, u_p$ so w\"{a}hlt, dass der
Periodicit\"{a}tsmodul von $u_\mu$ an dem Schnitte $a_\mu$ gleich
$\pi i$, an den \"{u}brigen Schnitten $a$ gleich $0$ ist, und man
den Periodicit\"{a}tsmodul von $u_\mu$ an dem Schnitte $b_\nu$
durch $a_{\mu, \nu}$ bezeichnet, so muss
$a_{\mu, \nu} = a_{\nu, \mu}$ und der reelle Theil von
$\sum\limits_{\mu, \mu'} a_{\mu, \mu'} m_\mu m_{\mu'}$
f\"{u}r alle reellen (ganzen) Werthe der $p$ Gr\"{o}ssen $m$
negativ sein.

\medbreak

\centerline{19.}

\nobreak\medskip

Die Zerlegung der Fl\"{a}che~$T$ werde nicht wie bisher nur durch
in sich zur\"{u}cklaufende Querschnitte, sondern folgendermassen
ausgef\"{u}hrt.  Man mache zuerst einen in sich
zur\"{u}cklaufenden die Fl\"{a}che nicht zerst\"{u}ckelnden
Schnitt~$a_1$ und f\"{u}hre dann einen Querschnitt~$b_1$ von der
positiven Seite von $a_1$ auf die negative zum Anfangspunkte
zur\"{u}ck, worauf die Begrenzung aus \emph{einem} St\"{u}cke
bestehen wird.  Einen dritten die Fl\"{a}che nicht
zerst\"{u}ckelnden Querschnitt kann man demzufolge (wenn die
Fl\"{a}che noch nicht einfach zusammenh\"{a}ngend ist) von einem
beliebigen Punkte dieser Begrenzung bis zu einem beliebigen
Begreunzungspunkte, also auch zu einem fr\"{u}heren Punkte dieses
Querschnitts f\"{u}hren.  Man thue das Letztere, so dass dieser
Querschnitt aus einer in sich zur\"{u}cklaufenden Linie $a_2$ und
einem dieser Linie voraufgehenden Theile~$c_1$ besteht, welche das
fr\"{u}here Schnittsystem mit ihr verbindet.  Den folgenden
Querschnitt $b_2$ ziehe man von der positiven Seite von $a_2$ auf
die negative zum Anfangspunkte zur\"{u}ck, worauf die Begrenzung
wieder aus \emph{einem} St\"{u}cke besteht.  Die weitere
Zerschneidung kann daher, wenn n\"{o}thig, wieder durch zwei in
demselben Punkte anfangende und endende Schnitte $a_3$ und $b_3$
und eine das System der Linien $a_2$ und $b_2$ mit ihnen
verbindende Linie~$c_2$ geschehen.  Wird dieser Verfahren
fortgesetzt, bis die Fl\"{a}che einfach zusammenh\"{a}ngend ist,
so erh\"{a}lt man ein Schnittnetz, welches aus $p$ Paaren von
zwei in einem und demselben Punkte anfangenden und endenden
Linien $a_1$ und $b_1$, $a_2$ und $b_2$,\dots, $a_p$ und $b_p$
besteht und aus $p - 1$ Linien $c_1, c_2,\ldots, c_{p-1}$, welche
jedes Paar mit dem folgenden verbinden.  Es m\"{o}ge $c_\nu$ von
einem Punkte von $b_\nu$ nach einem Punkte von $a_{\nu + 1}$
gehen.  Das Schnittnetz wird als so entstanden betrachtet, dass
der $\overline{2\nu - 1} \,$\-te Querschnitt aus $c_{\nu - 1}$ und
der von dem Endpunkte von $c_{\nu - 1}$ zu diesem
zur\"{u}ckgezogenen Linie $a_\nu$ besteht, und der $2\nu \,$\-te
durch die von der \emph{positiven} auf die \emph{negative} Seite
von $a_\nu$ gezogene Linie~$b_\nu$ gebildet wird.  Die Begrenzung
der Fl\"{a}che besteht bei dieser Zerschneidung nach einer
geraden Anzahl von Schnitten aus \emph{einem}, nach einer
ungeraden aus zwei St\"{u}cken.

Ein allenthalben endliches Integral~$w$ einer rationalen Function
von $s$ und $z$ nimmt dann zu beiden Seiten einer Linie~$c$
denselben Werth an.  Denn die ganze fr\"{u}her entstandene
Begrenzung besteht, wie bemerkt, aus \emph{einem} St\"{u}cke und
bei der Integration l\"{a}ngs derselben von der einen Seite der
Linie~$c$ bis auf die andere wird $\int \partial w$ durch jedes
fr\"{u}her entstandene Schnittelement zweimal, in
entgegengesetzter Richtung, erstreckt.  Eine solche Function ist
daher in $T$ allenthalben ausser der Linien $a$ und $b$ stetig.
Die durch diese Linien zerschnittene Fl\"{a}che~$T$ m\"{o}ge
durch $T''$ bezeichnet werden.

\medbreak

\centerline{20.}

\nobreak\medskip

Es seien nun $w_1, w_2,\ldots, w_p$ von einander unabh\"{a}ngige
solche Functionen, und der Periodicit\"{a}tsmodul von $w_\mu$ an
dem Querschnitte $a_\nu$ gleich $A_\mu^{(\nu)}$ und an dem
Querschnitte $b_\nu$ gleich $B_{\mu}^{(\nu)}$.  Es ist dann das
Integral $\int w_\mu \, dw_{\mu'}$, um die Fl\"{a}che~$T''$
positiv herum ausgedehnt, $= 0$, da die Function unter dem
Integralzeichen allenthalben endlich ist.  Bei dieser Integration
wird jede der Linien $a$ und $b$ zweimal, einmal in positiver und
einmal in negativer Richtung durchlaufen, und es muss w\"{a}hrend
jener Integration, wo sie als Begrenzung des positiverseits
gelegenen Gebiets dient, f\"{u}r $w_\mu$ der Werth auf der
positiven Seite oder $w_\mu^+$, w\"{a}hrend dieser der Werth auf
der negativen oder $w_\mu^-$ genommen werden.  Es ist also dies
Integral gleich der Summe aller Integrale
$\int (w_\mu^+ - w_\mu^-) \, dw_{\mu'}$
durch die Linien $a$ und $b$.  Die Linien~$b$ f\"{u}hren von der
positiven zur negativen Seite der Linien $a$, und folglich die
Linien~$a$ von der negativen zur positiven Seite der Linien~$b$.
Das Integral durch die Linie $a_\nu$ ist daher
\[ = \int A_\mu^{(\nu)} \, dw_{\mu'}
   = A_\mu^{(\nu)} \int dw_{\mu'}
   = A_\mu^{(\nu)} B_{\mu'}^{(\nu)},\]
und das Integral durch die Linie $b_\nu$
\[ = \int B_\mu^{(\nu)} \, dw_{\mu'}
   = - B_\mu^{(\nu)} A_{\mu'}^{(\nu)}.\]
Das Integral $\int w_\mu \, dw_{\mu'}$, um die Fl\"{a}che~$T''$
positive herum erstreckt, ist also
\[ = \sum_\nu ( A_\mu^{(\nu)} B_{\mu'}^{(\nu)}
      - B_\mu^{(\nu)} A_{\mu'}^{(\nu)} ),\]
und diese Summe folglich $= 0$.  Diese Gleichung gilt f\"{u}r je
zwei von den Functionen $w_1, w_2,\ldots, w_p$ und liefert also
$\displaystyle \frac{p(p - 1)}{1 \mathbin{.} 2}$
Relationen zwischen deren Periodicit\"{a}tsmoduln.

Nimmt man f\"{u}r die Functionen $w$ die Functionen $u$ oder
w\"{a}hlt man sie so, dass $A_\mu^{(\nu)}$ f\"{u}r ein von $\mu$
verschiedenes $\nu$ gleich $0$ und $A_\nu^{(\nu)} = \pi i$ ist,
so gehen diese Relationen \"{u}ber in
$B_{\mu'}^{(\mu)} \pi i - B_\mu^{(\mu')} \pi i = 0$
oder in $a_{\mu, \mu'} = a_{\mu', \mu}$.

\medbreak

\centerline{21.}

\nobreak\medskip

Es bleibt noch zu zeigen, dass die Gr\"{o}ssen~$a$ die zweite
oben n\"{o}thig gefundene Eigenschaft besitzen.

Man setze $w = \mu + \nu i$ und den Periodicit\"{a}tsmodul dieser
Function an dem Schnitte $a_\nu$ gleich
$A^{(\nu)} = \alpha_\nu + \gamma_\nu i$
und an dem Schnitte $b_\nu$ gleich
$B^{(\nu)} = \beta_\nu + \delta_\nu i$.
Es ist dann das Integral
\[ \int  \left(
            \left( \frac{\partial \mu}{\partial x} \right)^2
          + \left( \frac{\partial \mu}{\partial y} \right)^2
         \right) \, dT \]
oder\footnote{Dies Integral dr\"{u}ckt den Inhalt der Fl\"{a}che
aus, welche die Gesammtheit der Werthe, die $w$ innerhalb $T''$
annimmt, auf der $w$-Ebene repr\"{a}sentirt.}
\[ \int  \left(
            \frac{\partial \mu}{\partial x}
            \frac{\partial \nu}{\partial y}
          - \frac{\partial \mu}{\partial y}
            \frac{\partial \nu}{\partial x}
         \right) \, dT \]
durch die Fl\"{a}che $T''$ gleich dem Begrenzungsintegral
$\int \mu \, d\nu$ um $T''$ positiv herum erstreckt, also gleich
der Summe der Integrale $\int ( \mu^+ - \mu^- ) \, d\nu$ durch
die Linien $a$ und $b$.  Das Integral durch die Linie~$a_\nu$ ist
$= \alpha_\nu \int d\nu = \alpha_\nu \delta_\nu$,
das Integral durch die Linie $b_\nu$ gleich
$\beta_\nu \int d\nu = - \beta_\nu \gamma_\nu$,
und folglich
\[ \int  \left(
            \left( \frac{\partial \mu}{\partial x} \right)^2
          + \left( \frac{\partial \mu}{\partial y} \right)^2
         \right) \, dT
   = \sum_{\nu = 1}^{\nu = p}
         (\alpha_\nu \delta_\nu - \beta_\nu \gamma_\nu).\]
Diese Summe ist daher stets positiv.

Hieraus ergiebt sich die zu beweisende Eigenschaft der
Gr\"{o}ssen~$\alpha$, wenn man f\"{u}r $w$ setzt
$u_1 m_1 + u_2 m_2 + \cdots + u_p m_p$.
Denn es ist dann
$A^{(\nu)} = m_\nu \pi i$,
$B^{(\nu)} = \sum\limits_\mu a_{\mu , \nu} m_\mu$,
folglich $a_\nu$ stets $= 0$ und
\[ \int  \left(
            \left( \frac{\partial \mu}{\partial x} \right)^2
          + \left( \frac{\partial \mu}{\partial y} \right)^2
         \right) \, dT
   = - \sum \beta_\nu \gamma_\nu = - \pi \sum m_\nu \beta_\nu \]
oder gleich dem reellen Theile von
$-\pi \sum\limits_{\mu, \nu} a_{\mu, \nu} m_\mu m_\nu$,
welche also f\"{u}r alle reellen Werthe der Gr\"{o}ssen $m$
positiv ist.

\medbreak

\centerline{22.}

\nobreak\medskip

Setzt man nun in der $\vartheta$-Reihe (1) \S.~17 f\"{u}r
$a_{\mu, \mu'}$ den Periodicit\"{a}tsmodul der Function $u_\mu$
an dem Schnitt $b_{\mu'}$ und, durch $e_1, e_2,\ldots, e_p$
beliebige Constanten bezeichnend, $u_\mu - e_\mu$ f\"{u}r
$v_\mu$, so erh\"{a}lt man eine in jedem Punkte von $T$ eindeutig
bestimmte Function von $z$,
\[ \vartheta( u_1 - e_1, u_2 - e_2,\ldots, u_p - e_p ),\]
welche ausser den Linien~$b$ stetig und endlich und auf der
positiven Seite der Linie $b_\nu$
$\left( e^{-2(u_\nu - e_\nu)} \right) \,$\-mal so gross als auf der
negativen ist, wenn man den Functionen~$u$ in den Linien~$b$
selbst den Mittelwerth von den Werthen zu beiden Seiten beilegt.
F\"{u}r wie viele Punkte von $T'$ oder Werthepaare von $s$ und
$z$ diese Function unendlich klein von der ersten Ordung wird,
kann durch Betrachtung des Begrenzungsintegrals
$\int d \log \vartheta$,
um $T'$ positiv herum erstreckt, gefunden werden; denn dieses
Integral ist gleich der Anzahl dieser Punkte multiplicirt mit
$2\pi i$.  Andererseits ist dies Integral gleich der Summe der
Integrale
$\int ( d \log \vartheta^+ - d \log \vartheta^-)$
durch s\"{a}mmtliche Schnittlinien $a$, $b$ und $c$.  Die
Integrale durch die Linien $a$ und $c$ sind $= 0$, das Integral
durch $b_\nu$ aber gleich $- 2 \int du_\nu = 2\pi i$, die Summe
aller also $= p \, 2 \pi i$.  Die Function $\vartheta$ wird daher
unendlich klein von der ersten Ordnung in $p$ Punkten der
Fl\"{a}che~$T'$, welche durch $\eta_1, \eta_2,\ldots, \eta_p$
bezeichnet werden m\"{o}gen.

Durch einen positiven Umlauf des Punktes $(s,z)$ um einen dieser
Punkte w\"{a}chst $\log \vartheta$ um $2\pi i$, durch einen
positiven Umlauf um das Schnittepaar $a_\nu$ und $b_\nu$ um
$-2 \pi i$.  Um daher die Function $\log \vartheta$ allenthalben
eindeutig zu bestimmen, f\"{u}hre man von jedem Punkte~$\eta$
einen Schnitt durch das Innere nach je einem Linienpaar, von
$\eta_\nu$ den Schnitt $l_\nu$ nach $a_\nu$ und $b_\nu$, und zwar
nach ihrem gemeinschaftlichen Anfangs- und Endpunkte, und nehme
in der dadurch enstanden Fl\"{a}che $T^*$ die Function
allenthalben stetig an.  Sie ist dann auf der positiven Seite der
Linien $l$ um $-2\pi i$, auf der positiven Seite der Linie
$a_\nu$ um $g_\nu \, 2 \pi i$ und auf der positiven Seite der
Linie $b_\nu$ um
\[ - 2 (u_\nu - e_\nu) - h_\nu \, 2 \pi i \]
gr\"{o}sser, als auf der negativen, wenn $g_\nu$ und $h_\nu$
ganze Zahlen bezeichnen.

Die Lage der Punkte $\eta$ und die Werthe der Zahlen $g$ und $h$
h\"{a}ngen von den Gr\"{o}ssen $e$ ab, und diese Abh\"{a}ngigkeit
l\"{a}sst sich auf folgendem Wege n\"{a}her bestimmen.  Das
Integral
$\int \log \vartheta \, du_\mu$, um $T^*$
positiv herum erstreckt, ist $= 0$, da die Function
$\log \vartheta$ in $T^*$ stetig bleibt.  Dieses Integral ist
aber auch gleich der Summe der Integrale
$\int ( \log \vartheta^+ - \log \vartheta^-) \, du_\mu$
durch s\"{a}mmtliche Schnittlinien $l$, $a$, $b$ und $c$ und
findet sich, wenn man den Werth von $u_\mu$ im Punkte $\eta_\nu$
durch $\alpha_\mu^{(\nu)}$ bezeichnet,
\[ = 2 \pi i \left( \sum_\nu \alpha_\mu^{(\nu)}
      + h_\mu \, \pi i
      + \sum_\nu g_\nu a_{\nu, \mu}
      - e_\mu + k_\mu \right),\]
worin $k_\mu$ von den Gr\"{o}ssen $e$, $g$, $h$ und der Lage der
Punkte $\eta$ unabh\"{a}ngig ist.  Dieser Ausdruck ist also
$= 0$.

Die Gr\"{o}sse $k_\mu$ h\"{a}ngt von der Wahl der
Function~$u_\mu$ ab, welche durch die Bedingung, an dem
Schnitte~$a_\mu$ den Periodicit\"{a}tsmodul $\pi i$, an den
\"{u}brigen Schnitten~$a$ den Periodicit\"{a}tmodul $0$
anzunehmen, nur bis auf eine additive Constante bestimmt ist.
Nimmt man f\"{u}r $u_\mu$ eine um die Constante~$c_\mu$
gr\"{o}ssere Function und zugleich $e_\mu$ um $c_\mu$
gr\"{o}sser, so bleiben die Function~$\vartheta$ und folglich die
Punkte~$\eta$ und die Gr\"{o}ssen $g$, $h$ unge\"{a}ndert, der
Werth von $u_\mu$ im Punkte~$\eta_\nu$ aber wird
$\alpha_\mu^{(\nu)} + c_\mu$.  Es geht daher $k_\mu$ in
$k_\mu - (p - 1) c_\mu$ \"{u}ber und verschwindet, wenn
$\displaystyle c_\mu = \frac{k_\mu}{p - 1}$ genommen wird.

Man kann folglich, wie f\"{u}r die Folge geschehen soll, die
additiven Constanten in den Functionen~$u$ oder die Anfangswerthe
in den sie ausdr\"{u}ckenden Integralen so bestimmen, dass man
durch die Substitution von $u_\mu - \sum \alpha_\mu^{(\nu)}$ f\"{u}r
$v_\mu$ in $\log \vartheta (v_1,\ldots, v_p)$ eine Function erh\"{a}lt,
welche in den Punkten $\eta$ logarithmisch unendlich wird und,
durch $T^*$ stetig fortgesetzt, auf der positiven Seite der
Linien $l$ um $- 2 \pi i$, der Linien $a$ um $0$ und der Linie
$b_\nu$ um
$-2 ( u_\nu - \sum\limits_1^p \alpha_\nu^{(\mu)} )$
gr\"{o}sser wird, als auf der negativen.  Zur Bestimmung dieser
Anfangswerthe werden sich sp\"{a}ter leichtere Mittel darbieten,
als der obige Integralausdruck f\"{u}r $k_\mu$.

\medbreak

\centerline{23.}

\nobreak\medskip

Setzt man
$(u_1, u_2,\ldots, u_p)
   \equiv (\alpha_1^{(p)}, \alpha_2^{(p)},\ldots, \alpha_p^{(p)})$
nach den $2p$ Modulsystemen der Functionen $u$ (\S.~15), also
\[ (v_1, v_2,\ldots, v_p) \equiv
         \left(
          - \sum_1^{p-1} \alpha_1^{(\nu)}, \,
          - \sum_1^{p-1} \alpha_2^{(\nu)},\ldots, \,
          - \sum_1^{p-1} \alpha_p^{(\nu)}
         \right),\]
so wird $\vartheta = 0$.  Wird umgekehrt $\vartheta = 0$ f\"{u}r
$v_\mu = r_\mu$, so ist $(r_1, r_2,\ldots, r_p)$ einen
Gr\"{o}ssensysteme von der Form
\[       \left(
          - \sum_1^{p-1} \alpha_1^{(\nu)}, \,
          - \sum_1^{p-1} \alpha_2^{(\nu)},\ldots, \,
          - \sum_1^{p-1} \alpha_p^{(\nu)}
         \right) \]
congruent.  Denn setzt man
$v_\mu = u_\mu - \alpha_\mu^{(p)} + r_\mu$,
indem man $\eta_p$ beliebig w\"{a}hlt, so wird die Function
$\vartheta$ ausser in $\eta_p$ noch in $p - 1$ andern Punkten
unendlich klein von der ersten Ordnung, und bezeichnet man diese
durch $\eta_1, \eta_2,\ldots, \eta_{p-1}$, so ist
\[       \left(
          - \sum_1^{p-1} \alpha_1^{(\nu)}, \,
          - \sum_1^{p-1} \alpha_2^{(\nu)},\ldots, \,
          - \sum_1^{p-1} \alpha_p^{(\nu)}
         \right)
   \equiv (r_1, r_2,\ldots, r_p).\]

Die Function $\vartheta$ bleibt unge\"{a}ndert, wenn man
s\"{a}mmtliche Gr\"{o}ssen $v$ in's Entgegengesetzte verwandelt;
denn verwandelt man in der Reihe f\"{u}r
\[ \vartheta(v_1, v_2,\ldots, v_p) \]
s\"{a}mmtliche Indices $m$ in's Entgegengesetze, wodurch der
Werth der Reihe unge\"{a}ndert bleibt, da $- m_\nu$ dieselben
Werthe wie $m_\nu$ durchl\"{a}uft, so geht
$\vartheta(v_1, v_2,\ldots, v_p)$
\"{u}ber in
$\vartheta(-v_1, -v_2,\ldots, -v_p)$.

Nimmt man nun die Punkte $\eta_1, \eta_2,\ldots, \eta_{p-1}$
beliebig an, so wird
\[ \vartheta
         \left(
          - \sum\limits_1^{p-1} \alpha_1^{(\nu)},\ldots, \,
          - \sum\limits_1^{p-1} \alpha_p^{(\nu)}
         \right)
   = 0 \]
und folglich, da die Function $\vartheta$ wie eben bemerkt gerade
ist, auch
\[ \vartheta
         \left(
            \sum\limits_1^{p-1} \alpha_1^{(\nu)},\ldots, \,
            \sum\limits_1^{p-1} \alpha_p^{(\nu)}
         \right)
   =  0.\]
Es lassen sich also die $p - 1$ Punkte
$\eta_p, \eta_{p-1},\ldots, \eta_{2p - 2}$
so bestimmen, dass
\[       \left(
            \sum_1^{p-1} \alpha_1^{(\nu)},\ldots, \,
            \sum_1^{p-1} \alpha_p^{(\nu)}
         \right)
   \equiv
         \left(
          - \sum_p^{2p-2} \alpha_1^{(\nu)},\ldots, \,
          - \sum_p^{2p-2} \alpha_p^{(\nu)}
         \right) \]
und folglich
\[       \left(
            \sum_1^{2p-2} \alpha_1^{(\nu)},\ldots, \,
            \sum_1^{2p-2} \alpha_p^{(\nu)}
         \right)
   \equiv (0,\ldots, 0) \]
ist.  Die Lage der $p - 1$ letzten Punkte h\"{a}ngt dann von der
Lage der $p - 1$ ersten so ab, dass bei beliebiger stetiger
Aenderung derselben
$\sum\limits_1^{2p - 2} d\alpha_\pi^{(\nu)} = 0$
f\"{u}r $\pi = 1,2,\ldots, p$, und folglich sind (\S.~16) die
Punkte $\eta$ solche $2p - 2$ Punkte, f\"{u}r welche ein $dw$
unendlich klein von der zweiten Ordnung wird, oder wenn man den
Werth des Gr\"{o}ssenpaars $(s,z)$ im Punkte $\eta_\nu$ durch
$(\sigma_\nu, \zeta_\nu)$ bezeichnet, so sind
$(\sigma_1, \zeta_1),\ldots,$ $(\sigma_{2p-2}, \zeta_{2p-2})$
durch die Gleichung $\varphi = 0$ verkn\"{u}pfte Werthenpaare
(\S.~16).

{\itshape
Bei den hier gew\"{a}hlten Anfangswerthen der Integrale $u$ wird
also
\[       \left(
            \sum_1^{2p-2} u_1^{(\nu)},\ldots, \,
            \sum_1^{2p-2} u_p^{(\nu)}
         \right)
   \equiv (0,\ldots, 0),\]
wenn die Summationen \"{u}ber s\"{a}mmtliche von den
Gr\"{o}ssenpaaren $(\gamma_\varrho, \delta_\varrho)$} (\S.~6)
{\itshape
verschiedene gemeinschaftliche Wurzeln der Gleichung $F = 0$ und
der Gleichung
$c_1 \varphi_1 + c_2 \varphi_2 + \cdots + c_p \varphi_p = 0$
erstreckt werden, wobei die constanten Gr\"{o}ssen~$c$ beliebig
sind.}

Sind $\varepsilon_1, \varepsilon_2,\ldots, \varepsilon_m$
$m$~Punkte, f\"{u}r welche eine rationale Function $\xi$ von $s$
und $z$, die $m \,$\-mal unendlich von der ersten Ordnung wird,
denselben Werth annimmt, und $u_\pi^{(\mu)}$, $s_\mu$, $z_\mu$
die Werthe von $u_\pi$, $s$, $z$ im Punkte $\varepsilon_\mu$, so
ist (\S.~15)
$( \sum\limits_1^m u_1^{(\mu)}, \,
   \sum\limits_1^m u_2^{(\mu)},\ldots, \,
   \sum\limits_1^m u_p^{(\mu)} )$
congruent einem constanten, d.~h.\ vom Werthe der
Gr\"{o}sse~$\xi$ unabh\"{a}ngigen Gr\"{o}ssensysteme
$(b_1, b_2,\ldots, b_p)$, und es kann dann f\"{u}r jede beliebige
Lage eines Punktes $\varepsilon$ die Lage der \"{u}brigen so
bestimmt werden, dass
\[       \left(
            \sum_1^m u_1^{(\mu)},\ldots, \,
            \sum_1^m u_p^{(\mu)}
         \right)
   \equiv (b_1,\ldots, b_p).\]
Man kann daher, wenn $m = p$, $(u_1 - b_1,\ldots, u_p - b_p)$
und, wenn $m < p$,
\[    \left(
         u_1 - \sum_1^{p - m} \alpha_1^{(\nu)} - b_1,\ldots, \,
         u_p - \sum_1^{p - m} \alpha_p^{(\nu)} - b_p
      \right) \]
f\"{u}r jede beliebige Lage des Punktes $(s,z)$ und der $p - m$
Punkte $\eta$ auf die Form
$( - \sum\limits_1^{p - 1} \alpha_1^{(\nu)},\ldots, \,
   - \sum\limits_1^{p - 1} \alpha_p^{(\nu)} )$
bringen, indem man einen der Punkte $\varepsilon$ mit $(s, z)$
zusammenfallen l\"{a}sst, und folglich ist
\[ \vartheta
      \left(
         u_1 - \sum_1^{p - m} \alpha_1^{(\nu)} - b_1,\ldots, \,
         u_p - \sum_1^{p - m} \alpha_p^{(\nu)} - b_p
      \right) \]
f\"{u}r jedwede Werthe des Gr\"{o}ssenpaars $(s,z)$ und der
$p - m$ Gr\"{o}ssenpaare $(\sigma_\nu, \zeta_\nu)$ gleich~$0$.

\medbreak

\centerline{24.}

\nobreak\medskip

Aus der Untersuchung des \S.~22 folgt als Corollar, dass ein
beliebig gegebenes Gr\"{o}ssensystem $(e_1,\ldots, e_p)$ immer
einem und nur einem Gr\"{o}ssensysteme von der Form
$( \sum\limits_1^p \alpha_1^{(\nu)},\ldots,
   \sum\limits_1^p \alpha_p^{(\nu)} )$
congruent ist, wenn die Function
$\vartheta( u_1 - e_1,\ldots, u_p - e_p )$
nicht identisch verschwindet; denn es m\"{u}ssen dann die
Punkte~$\eta$ die $p$ Punkte sein, f\"{u}r welche diese
Function~$0$ wird.  Wenn aber
$\vartheta( u_1^{(p)} - e_1,\ldots, u_p^{(p)} - e_p )$
f\"{u}r jeden Werth von $(s_p, z_p)$ verschwindet, so l\"{a}sst
sich
\[ (u_1^{(p)} - e_1,\ldots, u_p^{(p)} - e_p)
   \equiv
         \left(
          - \sum_1^{p - 1} u_1^{(\nu)},\ldots, \,
          - \sum_1^{p - 1} u_p^{(\nu)}
         \right) \]
setzen (\S.~23), und es lassen sich also f\"{u}r jeden Werth des
Gr\"{o}ssenpaars $(s_p, z_p)$ die Gr\"{o}ssenpaare
$(s_1, z_1),\ldots,$ $(s_{p-1}, z_{p-1})$
so bestimmen, dass
\[       \left(
            \sum_1^p u_1^{(\nu)},\ldots, \,
            \sum_1^p u_p^{(\nu)}
         \right)
   \equiv (e_1,\ldots, e_p),\]
und folglich, bei stetiger Aenderung von $(s_p, z_p)$,
$\sum\limits_1^p du_\pi^{(\nu)} = 0$
ist f\"{u}r $\pi = 1,2,\ldots,p$.  Die $p$ Gr\"{o}ssenpaare
$(s_\nu, z_\nu)$ sind daher $p$ von den Gr\"{o}ssenpaaren
$(\gamma_\varrho, \delta_\varrho)$ verschiedene Wurzeln einer
Gleichung $\varphi = 0$, deren Coefficienten so sich \"{a}ndern,
dass die \"{u}brigen $p - 2$ Wurzeln constant bleiben.
Bezeichnet man die Werthe von $u_\pi$ f\"{u}r diese $p - 2$
Werthenpaare von $s$ und $z$ durch
$u_\pi^{(p+1)}, u_\pi^{(p+2)},\ldots, u_\pi^{(2p - 2)}$,
so ist
\[       \left(
            \sum_1^{2p - 2} u_1^{(\nu)},\ldots, \,
            \sum_1^{2p - 2} u_p^{(\nu)}
         \right)
   \equiv (0,\ldots, 0),\]
und folglich
\[ (e_1,\ldots, e_p) \equiv
         \left(
          - \sum_{p+1}^{2p - 2} u_1^{(\nu)},\ldots, \,
          - \sum_{p+1}^{2p - 2} u_p^{(\nu)}
         \right).\]
Umgekehrt ist, wenn diese Congruenz stattfindet,
\[ \vartheta( u_1^{(p)} - e_1,\ldots, u_p^{(p)} - e_p)
   = \vartheta
         \left(
            \sum_p^{2p - 2} u_1^{(\nu)},\ldots, \,
            \sum_p^{2p - 2} u_p^{(\nu)}
         \right)
   = 0.\]

{\itshape
Ein beliebig gegebenes Gr\"{o}ssensystem $(e_1,\ldots, e_p)$ ist
also nur einem Gr\"{o}ssensystem von der Form
$( \sum\limits_1^p \alpha_1^{(\nu)},\ldots,
   \sum\limits_1^p \alpha_p^{(\nu)} )$
congruent, wenn es nicht einem Gr\"{o}ssensysteme von der Form
$( - \sum\limits_1^{p - 2} \alpha_1^{(\nu)},\ldots,
   - \sum\limits_1^{p - 2} \alpha_p^{(\nu)} )$
congruent ist, und unendlich vielen, wenn dieses stattfindet.}

Da
\[ \vartheta
      \left(
         u_1 - \sum_1^p \alpha_1^{(\mu)},\ldots,
         u_p - \sum_1^p \alpha_p^{(\mu)}
      \right)
   = \vartheta
      \left(
         \sum_1^p \alpha_1^{(\mu)} - u_1,\ldots,
         \sum_1^p \alpha_p^{(\mu)} - u_p
      \right),\]
so ist $\vartheta$ eine ganz \"{a}hnliche Function wie von
$(s,z)$ auch von jedem der $p$ Gr\"{o}ssenpaare
$(\sigma_\mu, \zeta_\mu)$.  Diese Function von
$(\sigma_\mu, \zeta_\mu)$ wird $= 0$ f\"{u}r das Werth\-enpaar
$(s,z)$ und f\"{u}r die den \"{u}brigen $p - 1$ Gr\"{o}ssenpaaren
$(\sigma, \zeta)$ durch die Gleichung $\varphi = 0$
verkn\"{u}pften $p - 1$ Punkte.  Denn bezeichnet man den Werth
von $u_\pi$ in diesen Punkten mit
$\beta_\pi^{(1)}, \beta_\pi^{(2)},\ldots, \beta_\pi^{(p-1)}$,
so ist
\[       \left(
            \sum_1^p \alpha_1^{(\mu)},\ldots, \,
            \sum_1^p \alpha_p^{(\mu)}
         \right)
   \equiv
         \left(
            \alpha_1^{(\mu)} - \sum_1^{p - 1} \beta_1^{(\nu)},\ldots, \,
            \alpha_p^{(\mu)} - \sum_1^{p - 1} \beta_p^{(\nu)}
         \right) \]
und folglich $\vartheta = 0$, wenn $\eta_\mu$ mit einem dieser
Punkte oder mit dem Punkte $(s,z)$ zusammenf\"{a}llt.

\medbreak

\centerline{25.}

\nobreak\medskip

Aus den bisher entwickelten Eigenschaften der
Function~$\vartheta$ ergiebt sich der Ausdruck von
$\log \vartheta$ durch Integrale algebraischer Functionen von
$(s,z)$, $(\sigma_1, \zeta_1),\ldots,$ $(\sigma_p, \zeta_p)$.

Die Gr\"{o}sse
\[    \log \vartheta \left(
            u_1^{(2)} - \sum_1^p \alpha_1^{(\mu)},\ldots \right)
    - \log \vartheta \left(
            u_1^{(1)} - \sum_1^p \alpha_1^{(\mu)},\ldots \right) \]
ist, als Function von $(\sigma_\mu, \zeta_\mu)$ betrachtet, eine
Function von der Lage des Punkts $\eta_\mu$, welche im Punkte
$\varepsilon_1$, wie $- \log (\zeta_\mu - z_1)$, im Punkte
$\varepsilon_2$, wie $\log (\zeta_\mu - z_2)$ unstetig wird und
auf der positiven Seite einer von $\varepsilon_1$ nach
$\varepsilon_2$ zu ziehenden Linie um $2\pi i$, auf der positiven
Seite der Linie $b_\nu$ um $2(u_\nu^{(1)} - u_\nu^{(2)})$
gr\"{o}sser ist, als auf der negativen, ausser den Linien~$b$ und
der Verbindungslinie von $\varepsilon_1$ und $\varepsilon_2$ aber
allenthalben stetig bleibt.  Bezeichnet nun
$\varpi^{(\mu)} (\varepsilon_1, \varepsilon_2)$
irgend eine Function von $(\sigma_\mu, \zeta_\mu)$, welche ausser
den Linien~$b$ ebenso unstetig ist und auf der einen Seite einer
solchen Linie ebenfalls um eine Constante gr\"{o}sser ist, als
auf der andern, so unterscheidet sie sich (\S.~3) von dieser nur
um eine von $(\sigma_\mu, \zeta_\mu)$ unabh\"{a}ngige Gr\"{o}sse,
und folglich ist sie von
$\sum\limits_1^p \varpi^{(\mu)} (\varepsilon_1, \varepsilon_2)$
nur um eine von s\"{a}mmtlichen Gr\"{o}ssen $(\sigma, \zeta)$
unabh\"{a}ngige und also bloss von $(s_1, z_1)$ und $(s_2, z_2)$
abh\"{a}ngende Gr\"{o}sse verschieden.
$\varpi^{(\mu)} (\varepsilon_1, \varepsilon_2)$
dr\"{u}ckt den Werth einer Function
$\varpi (\varepsilon_1, \varepsilon_2)$
des \S.~4 f\"{u}r $(s,z) = (\sigma_\mu, \zeta_\mu)$ aus, deren
Periodicit\"{a}tsmoduln an den Schnitten~$a$ gleich $0$ sind.
Aendert man diese Function um die Constante $c$, so \"{a}ndert
sich
$\sum\limits_1^p \varpi^{(\mu)} (\varepsilon_1, \varepsilon_2)$
um $pc$; man kann daher, wie f\"{u}r die Folge geschehen soll,
die additive Constante in der Function
$\varpi (\varepsilon_1, \varepsilon_2)$
oder den Anfangswerth in dem sie darstellenden Integrale dritter
Gattung so bestimmen, dass
\[ \log \vartheta^{(2)} - \log \vartheta^{(1)}
   = \sum_1^p \varpi^{(\mu)} (\varepsilon_1, \varepsilon_2).\]
Da $\vartheta$ von jedem der Gr\"{o}ssenpaare $(\sigma, \zeta)$
auf \"{a}hnliche Art, wie von $(s,z)$ abh\"{a}ngt, so kann die
Aenderung von $\log \vartheta$, wenn irgend eins der
Gr\"{o}ssenpaare
$(s,z)$, $(\sigma_1, \zeta_1),\ldots,$ $(\sigma_p, \zeta_p)$
eine endliche Aenderung erleidet, w\"{a}hrend die \"{u}brigen
constant bleiben, durch eine Summe von Functionen
$\varpi$ ausgedr\"{u}ckt werden.  Offenbar kann man also,
indem man nach und nach die einzelnen Gr\"{o}ssenpaare
$(s,z)$, $(\sigma_1, \zeta_1),\ldots,$ $(\sigma_p, \zeta_p)$
\"{a}ndert, $\log \vartheta$ ausdr\"{u}cken durch eine Summe von
Functionen~$\varpi$ und
$\log \vartheta (0,0,\ldots,0)$
oder den Werth von $\log \vartheta$ f\"{u}r ein beliebiges
anderes Werthensystem.  Die Bestimmung von
$\log \vartheta (0,0,\ldots, 0)$
als Function der $3p - 3$ Moduln des System rationaler Function
von $s$ und $z$ (\S.~12) erfordert \"{a}hnliche Betrachtungen,
wie sie von \emph{Jacobi} in seinen Arbeiten \"{u}ber
elliptische Functionen zur Bestimmung von $\Theta(0)$ angewandt
worden sind.  Man kann dazu gelangen, indem man mit H\"{u}lfe
der Gleichungen
\[ 4 \frac{\partial \vartheta}{\partial a_{\mu, \mu}}
   = \frac{\partial^2 \vartheta}{\partial v_\mu^2}
     \quad \mbox{ und } \quad
   2 \frac{\partial \vartheta}{\partial a_{\mu, \mu'}}
   = \frac{\partial^2 \vartheta}{\partial v_\mu \, \partial v_{\mu'}},\]
wenn $\mu$ von $\mu'$ verschieden ist, die Differentialquotienten
von $\log \vartheta$ nach den Gr\"{o}ssen~$a$ in
\[ d \log \vartheta
   = \sum \frac{\partial \log \vartheta}{\partial a_{\mu, \mu'}}
      \, da_{\mu, \mu'} \]
durch Integrale algebraischer Functionen ausdr\"{u}ckt.  F\"{u}r
die Ausf\"{u}hrung dieser Rechnung scheint jedoch eine
ausf\"{u}hrlichere Theorie der Functionen, welche einer linearen
Differentialgleichung mit algebraischen Coefficienten
gen\"{u}gen, n\"{o}thig, die ich nach den hier angewandten
Principien n\"{a}chstens zu liefern beabsichtige.

Ist $(s_2, z_2)$ unendlich wenig von $(s_1, z_1)$ verschieden, so
geht
$\varpi (\varepsilon_1, \varepsilon_2)$
\"{u}ber in $\partial z_1 \, t(\varepsilon_1)$, worin
$t(\varepsilon_1)$ ein Integral zweiter Gattung einer rationaler
Function von $s$ und $z$ ist, welches in $\varepsilon_1$ wie
$\displaystyle \frac{1}{z - z_1}$
unstetig wird und an den Schnitten~$a$ den
Periodicit\"{a}tmodul~$0$ hat; und es ergiebt sich, dass der
Periodicit\"{a}tsmodul eines solchen Integrals an dem
Schnitte~$b_\nu$ gleich
$\displaystyle 2 \frac{\partial u_\nu^{(1)}}{\partial z_1}$
ist und die Integrationsconstante sich so bestimmen l\"{a}sst,
dass die Summe der Werthe von $t(\varepsilon_1)$ f\"{u}r die $p$
Werthenpaare
$(\sigma_1, \zeta_1),\ldots,$ $(\sigma_p, \zeta_p)$
gleich
$\displaystyle \frac{\partial \log \vartheta^{(1)}}{\partial z_1}$
wird.  Es ist dann
$\displaystyle \frac{\partial \log \vartheta^{(1)}}{\partial \zeta_\mu}$
gleich der Summe der Werthe von $t(\eta_\nu)$ f\"{u}r die den
$p - 1$ von $(\sigma_\mu, \zeta_\mu)$ verschiedenen
Gr\"{o}ssenpaaren $(\sigma, \zeta)$ durch die Gleichung
$\varphi = 0$ verkn\"{u}pften $p - 1$ Werthenpaare und f\"{u}r
das Werthenpaar $(s,z)$, und man erh\"{a}lt f\"{u}r
\[ \frac{\partial \log \vartheta^{(1)}}{\partial z_1} \, dz_1
   + \sum_1^p \frac{\partial \log \vartheta^{(1)}}{\partial \zeta_\mu}
         \, d \zeta_\mu
   = d \log \vartheta^{(1)},\]
einen Ausdruck, welchen \emph{Weierstrass} f\"{u}r den Fall,
wenn $s$ nur eine zweiwerthige Function von $z$ ist, gegeben hat
(d.~J.  Bd.~47  S.~300 Form.~35).

Die Eigenschaften von
$\varpi (\varepsilon_1, \varepsilon_2)$
und $t(\varepsilon_1)$ als Functionen von $(s_1, z_1)$ und $(s_2, z_2)$
ergeben sich aus den Gleichungen
\[ \varpi (\varepsilon_1, \varepsilon_2)
   = \frac{1}{p}
      \left(
         \log \vartheta (u_1^{(2)} - p u_1,\ldots )
       - \log \vartheta (u_1^{(1)} - p u_1,\ldots )
      \right) \]
und
\[ t(\varepsilon_1)
   = \frac{1}{p} \frac{
      \partial \log \vartheta (u_1^{(1)} - p u_1,\ldots)}{
      \partial z_1},\]
welche in den obigen Ausdr\"{u}cken f\"{u}r
$\log \vartheta^{(2)} - \log \vartheta^{(1)}$
und
$\displaystyle \frac{\partial \log \vartheta^{(1)}}{\partial z_1}$
als specielle F\"{a}lle enthalten sind.

\medbreak

\centerline{26.}

\nobreak\medskip

Es soll jetzt die Aufgabe behandelt werden, algebraische
Functionen von $z$ als Quotienten zweier Producte von
gleichvielen Functionen
$\vartheta (u_1 - e_1,\ldots, )$
und Potenzen der Gr\"{o}ssen $e^u$ darzustellen.

Ein solcher Ausdruck erlangt bei den Ueberg\"{a}ngen von $(s,z)$
\"{u}ber die Querschnitte constante Factoren, und diese
m\"{u}ssen Wurzeln der Einheit sein, wenn er algebraisch von $z$
abh\"{a}ngen und also bei stetiger Fortsetzung f\"{u}r dasselbe
$z$ nur eine endliche Anzahl von Werthen annehmen soll.  Sind
alle diese Factoren $\mu \, \,$\-te Wurzeln der Einheit, so ist die
$\mu \, \,$\-te Potenz des Ausdrucks eine einwerthige und folglich
rationale Function von $s$ und $z$.

Umgekehrt l\"{a}sst sich leicht zeigen, dass jede algebraische
Function $r$ von $z$, die innerhalb der ganzen Fl\"{a}che~$T'$
stetig fortgesetzt, allenthalben nur \emph{einen} bestimmten
Werth annimmt und beim Ueberschreiten eines Querschnitts einen
constanten Factor erlangt, sich auf mannigfaltige Art als
Quotient zweier Producte von $\vartheta$-Functionen und Potenzen
der Gr\"{o}ssen $e^u$ ausdr\"{u}cken l\"{a}sst.  Man bezeichne
einen Werth von $u_\mu$ f\"{u}r $r = \infty$ durch $\beta_\mu$
und f\"{u}r $r = 0$ durch $\gamma_\mu$ und nehme $\log r$, indem
man von jedem Punkte, wo $r$ unendlich von der ersten Ordnung
wird, nach je einem Punkte, wo $r$ unendlich klein von der ersten
Ordnung wird, eine Linie durch das Innere von $T'$ zieht, ausser
diesen Linien in $T'$ allenthalben stetig an.  Ist dann $\log r$
auf der positiven Seite der Linie $b_\nu$ um $g_\nu \, 2\pi i$
und auf der positiven Seite der Linie $a_\nu$ um
$- h_\nu \, 2\pi i$ gr\"{o}sser, als auf der negativen, so
ergiebt sich durch die Betrachtung des Begrenzungsintegrals
$\int \log r \, du_\mu$
\[ \sum \gamma_\mu - \sum \beta_\mu
   = g_\mu \, \pi i + \sum_\nu h_\nu a_{\mu, \nu} \]
f\"{u}r $\mu = 1,2,\ldots, p$, worin $g_\nu$ und $h_\nu$ nach dem
oben Bemerkten rationale Zahlen sein m\"{u}ssen und die Summen
auf der linken Seite der Gleichung \"{u}ber s\"{a}mmtliche
Punkte, wo $r$ unendlich klein oder unendlich gross von der
ersten Ordnung wird, auszudehnen sind, indem man einen Punkt, wo
$r$ unendlich klein oder unendlich gross von einer h\"{o}heren
Ordnung wird, als aus mehreren solchen Punkten bestehend
betrachtet (\S.~2).  Wenn diese Punkte bis auf $p$ gegeben sind,
so lassen sich diese $p$ immer und, allgemein zu reden, nur auf
eine Weise so bestimmen, dass die $2p$ Factoren
$e^{g_\nu \, 2 \pi i}$, $e^{- h_\nu \, 2 \pi i}$,
gegebene Werthe annehmen (\S\S.~15, 24).

Wenn man nun in dem Ausdrucke
\[ \frac{P}{Q} e^{- 2 \sum h_\nu u_\nu},\]
worin $P$ und $Q$ Producte von gleichvielen Functionen
$\vartheta( u_1 - \sum \alpha_1^{(\pi)},\ldots )$
mit demselben $(s, z)$ und verschiedenen $(\sigma, \zeta)$ sind,
die Werthenpaare von $s$ und $z$ f\"{u}r welche $r$ unendlich
wird, f\"{u}r Gr\"{o}ssenpaare $(\sigma, \zeta)$ in den
$\vartheta$-Functionen des Nenners und die Werthenpaare, f\"{u}r
welche $r$ verschwindet, f\"{u}r Gr\"{o}ssenpaare
$(\sigma, \zeta)$ in den $\vartheta$-Functionen des Z\"{a}hlers
substituirt und die \"{u}brigen Gr\"{o}ssenpaare
$(\sigma, \zeta)$ im Nenner und im Z\"{a}hler gleich annimmt, so
stimmt der Logarithme dieses Ausdrucks in Bezug auf die
Unstetigkeiten im Innern von $T'$ mit $\log r$ \"{u}berein, und
\"{a}ndert sich beim Ueberschreiten der Linien $a$ und $b$, wie
$\log r$, nur um rein imagin\"{a}re l\"{a}ngs diesen Linien
constante Gr\"{o}ssen; er unterscheidet sich also von $\log r$
nach dem \emph{Dirichlet}'schen Princip nur um eine Constante
und der Ausdruck selbst von $r$ nur durch einen constanten
Factor.  Bei dieser Substitution darf selbstredend keine der
$\vartheta$-Functionen identisch, f\"{u}r jeden Werth von $z$,
verschwinden.  Dieses w\"{u}rde geschehen (\S.~23), wenn
s\"{a}mmtliche Werthenpaare, f\"{u}r welche eine einwerthige
Function von $(s,z)$ verschwindet, f\"{u}r Gr\"{o}ssenpaare
$(\sigma, \zeta)$ in einer und derselben $\vartheta$-Function
substituirt w\"{u}rden.

\medbreak

\centerline{27.}

\nobreak\medskip

Als Quotient \emph{zweier} $\vartheta$-Functionen, multiplicirt
mit Potenzen der Gr\"{o}ssen $e^u$, l\"{a}sst sich demnach eine
einwerthige oder rationale Function von $(s,z)$ nicht darstellen.
Alle Functionen~$r$ aber, die f\"{u}r dasselbe Werthenpaar von
$s$ und $z$ mehrere Werthe annehmen und nur f\"{u}r $p$ oder
weniger Werthenpaare unendlich von der ersten Ordnung werden,
sind in dieser Form darstellbar und umfassen alle in dieser Form
darstellbaren algebraischen Functionen von $z$.  Man erh\"{a}lt,
abgesehen von einem constanten Factor, jede und jede nur einmal,
wenn man in
\[ \frac{
   \vartheta( v_1 - g_1 \pi i - \sum\limits_\nu h_\nu a_{1,\nu},\ldots)}{
   \vartheta( v_1,\ldots, v_p)}
   e^{-2 \sum\limits_\nu v_\nu h_\nu} \]
f\"{u}r $h_\nu$ und $g_\nu$ rationale \"{a}chte Br\"{u}che und
$\displaystyle u_\nu - \sum\limits_1^p \alpha_\nu^{(\mu)}$
f\"{u}r $v_\nu$ setzt.

Diese Gr\"{o}sse ist zugleich eine algebraische Function von
jeder der Gr\"{o}ssen~$\zeta$ und die (im vor.~\S.) entwickelten
Principien reichen v\"{o}llig hin, um sie durch die Gr\"{o}ssen
$z, \zeta_1,\ldots, \zeta_p$ algebraisch auszudr\"{u}cken.

In der That: Als Function von $(s,z)$ nimmt sie, durch die ganze
Fl\"{a}che $T'$ stetig fortgesetzt, allenthalben \emph{einen}
bestimmten  Werth an, wird unendlich von der ersten Ordnung
f\"{u}r die Werthenpaare
$(\sigma_1, \zeta_1),\ldots,$ $(\sigma_p, \zeta_p)$
und erlangt an dem Schnitte $a_\nu$ beim Uebergange von der
positiven zur negativen Seite den Factor $e^{h_\nu \, 2 \pi i}$,
an dem Schnitte $b_\nu$ den Factor $e^{- g_\nu \, 2 \pi i}$; und
jede andere dieselben Bedingungen erf\"{u}llende Function von
$(s,z)$ unterscheidet sich von ihr nur durch einen von $(s,z)$
unabh\"{a}ngigen Factor.  Als Function von
$(\sigma_\mu, \zeta_\mu)$ nimmt sie, durch die ganze
Fl\"{a}che~$T'$ stetig fortgesetzt, allenthalben \emph{einen}
bestimmten Werth an, wird unendlich von der ersten Ordnung
f\"{u}r das Werth\-enpaar $(s,z)$ und f\"{u}r die den \"{u}brigen
$p - 1$ Gr\"{o}ssenpaaren $(\sigma, \zeta)$ durch die Gleichung
$\varphi = 0$ verkn\"{u}pften $p - 1$ Werthenpaare
$(\sigma_1^{(\mu)}, \zeta_1^{(\mu)}),\ldots$
$(\sigma_{p-1}^{(\mu)}, \zeta_{p-1}^{(\mu)})$
und erlangt an dem Schnitte~$a_\nu$ den Factor
$e^{- h_\nu \, 2 \pi i}$, an den Schnitte~$b_\nu$ den Factor
$e^{g_\nu \, 2 \pi i}$; und jede andere dieselben Bedingungen
erf\"{u}llende Function von $(\sigma_\mu, \zeta_\mu)$
unterscheidet sich von ihr nur durch einen von
$(\sigma_\mu, \zeta_\mu)$ unabh\"{a}ngigen Factor.  Bestimmt man
also eine algebraische Function von
$z, \zeta_1,\ldots, \zeta_p$
\[ f( (s,z);\enspace
      (\sigma_1, \zeta_1),\ldots,\enspace
      (\sigma_p, \zeta_p) ) \]
so, dass sie als Function von jeder dieser Gr\"{o}ssen dieselben
Eigenschaften besitzt, so unterscheidet sie sich von dieser nur
durch einen von s\"{a}mmtlichen Gr\"{o}ssen
$z, \zeta_1,\ldots, \zeta_p$
unabh\"{a}ngigen Factor und wird also $= Af$, wenn $A$ diesen
Factor bezeichnet.  Um diesen Factor zu bestimmen, dr\"{u}cke man
in $f$ die von $(\sigma_\mu, \zeta_\mu)$ verschiedenen
Gr\"{o}ssenpaare $(\sigma, \zeta)$ durch
$(\sigma_1^{(\mu)}, \zeta_1^{(\mu)}),\ldots,$
$(\sigma_{p-1}^{(\mu)}, \zeta_{p-1}^{(\mu)})$
aus, wodurch er in
\[ g( (\sigma_\mu, \zeta_\mu);\enspace
      (s,z),\enspace
      (\sigma_1^{(\mu)}, \zeta_1^{(\mu)}),\ldots,\enspace
      (\sigma_{p-1}^{(\mu)}, \zeta_{p-1}^{(\mu)}) ) \]
\"{u}bergehe; offenbar erh\"{a}lt man dann den inversen Werth der
dazustellenden Function und also einen Ausdruck, welcher
$\displaystyle= \frac{1}{Af}$
sein muss, wenn man in $Ag$ f\"{u}r $(\sigma_\mu, \zeta_\mu)$ das
Grossenpaar $(s,z)$ und f\"{u}r die Gr\"{o}ssenpaare
$(s,z)$, $(\sigma_1^{(\mu)}, \zeta_1^{(\mu)}),\ldots,$
$(\sigma_{p-1}^{(\mu)}, \zeta_{p-1}^{(\mu)})$
die Werthenpaare von $(s,z)$ substituirt, f\"{u}r welche die
darzustellende Function und also $f = 0$ wird.  Hieraus ergiebt
sich $A^2$ und also $A$ bis auf das Vorzeichen, welches durch
directe Betrachtung der $\vartheta$-Reihen in dem darzustellenden
Ausdrucke gefunden werden kann.

G\"{o}ttingen, 1857.

\end{document}

